First Lecture
Monday, April 16, 1804
Honorable Guests: 1
Nothing about the venture which we now jointly begin is as difficult as the
beginning; indeed, even the escape route which I just started to take by begin-
ning with a consideration of the difficulty of beginning has its own difficul-
ties. No means remain except to cut the knot boldly and to ask you to accept
that what I am about to say is aimed at the wide world generally and applies
to it but not at all to you. 2
Namely: in my view the chief characteristic of our time 3 is that in it life
has become merely historical and symbolic, while real living is scarcely ever
found. One not insignificant aspect of life is thinking. Where all life has degen-
erated into a strange tale, the same must also happen to thinking. Of course one
will have heard and made a note of the fact that, among other things, human
beings can think; indeed that many of them have thought, the first in this way,
the second differently, and the third and fourth each in yet another way and
that all have failed in some fashion. It is not easy to decide to undertake this
thought process again for oneself. One who assumes responsibility for arousing
an era like ours to this decision must face this discomfort among others: he 4
doesn’t know where he might look for, and find, the people who need arous-
ing. Whomever he accuses has a ready answer, “Yes, that is certainly true of
others, but not of us”; and they are right to the extent that, as well as knowing
the criticized form of thought they are also familiar historically with the oppo-
site schools, so that if someone attacks them on one side, they are ready to flee
to the very position they currently reject. So, for example, if one speaks in the
way I did just now, rebuking our historical superficiality, dispersion into the
The lecture begins at GA II, 8, pp. 2–3. Double numbers in square brackets refer to cor-
responding pages in GA II, 8. In that volume the two existing texts are printed on fac-
ing pages. Because translation sometimes alters word order, and page breaks in SW and
Copia can differ, these page-break markers are sometimes only approximate. Through-
out, I have worked primarily from SW and follow it as closely as is practicable.
2122
The Science of Knowing
most multifarious and contradictory opinions, indecisiveness about everything
altogether, and absolute indifference to truth [4–5] in the way I have just
rebuked these things, then everyone will be sure that he does not recognize
himself in this picture, that he knows very well that only one thing can be true
and that all contradictions must necessarily be false. If we accused him of dog-
matic rigidity and one-sidedness for his adherence to the one, the very same
person would praise his all-around skepticism. 5
In such a state of affairs there is nothing for it but to assert briefly and
sweetly at a single stroke that I presuppose here in all seriousness that there is
a truth which alone is true and everything apart from this is unconditionally
false; further, that this truth can actually be found and be immediately evident
as unconditionally true; but that not even the least spark of it can be grasped
or communicated historically as an appropriation from someone else’s mind.
Rather, whomever would have it must produce it entirely out of himself. The
presenter can only provide the terms for insight; each individual must fulfill
these terms in himself, applying his living spirit to it with all his might, and
then the insight will happen of itself without any further ado. There is no
question here of an object which is already well known from other contexts,
but of something completely new, unheard of, and totally unknown to anyone
who has not already studied the science of knowing 6 thoroughly. No one can
arrive at this unknown unless it produces itself in him, but it does this only
under the condition that this very person 7 produces something, namely the
conditions for insight’s self-production. Whoever does not do this will never
obtain the object about which we will speak here. And since our whole dis-
course will be about this object, he will have no object at all ; for him our entire
discourse will be words about pure, bare nothingness, an empty vessel, word
breath, the mere movement of air, and nothing more. And so let this, taken
rigorously and literally, serve as the first prolegomenon.
I have still more to add, which however first presupposes the following.
From now on, honorable guests, I 8 wish to be considered silenced and erased,
and you yourselves must come forward and stand in my place. From now on,
everything which is to be thought in this assembly should be thought and be
true only to the extent that you yourselves have thought it and seen it to be true. 9
I said that I have more to add by way of introduction, and I will devote
this week’s four lectures to this task. Repeated experience compels me to
remind you explicitly that these introductory remarks should not be viewed as
most are, that is, as a simple [6–7] approach which the lecturer takes and
whose content has no very great significance. The introductory remarks I will
present have meaning, and what follows will be entirely lost without them.
They are designed to call your spiritual eye away from the objects over which
it has heretofore glided to and fro, directing it to the point which we must
consider here, and indeed to give this point its existence for the very first time.First Lecture
23
They are designed to initiate you into the art which we will subsequently prac-
tice together, the art of philosophy. They should simultaneously acquaint you,
and make you fluent, with a system of rules and maxims of thinking whose
employment will recur in every lecture. 10
I hope the matters which are to be handled in the introduction will
become easily comprehensible to everyone of even moderate attention; but
past experience requires that I add a word about just this comprehending.
First, one must not assume that the standard of comprehensibility for the sci-
ence of knowing in general is comparable to the standard of study and atten-
tion which the introduction requires; since anyone assuming this will be
unpleasantly disappointed later on. Thus, whomever has heard and under-
stood the introductory remarks has acquired a true and fitting concept of the
science of knowing sanctioned by the very originator of this science, but the
listener has not thereby acquired the tiniest spark of the science itself. This
universally applicable distinction between the mere concept and the real, true
substance has particular bearing here. Possessing the concept has its uses; for
example, it protects us from the absurdity of underestimating and misjudging
what we do not possess; only no one should believe that by this possession,
which is not in any case all that rare, one has become a philosopher. One is
and remains a sophist, only to be sure less superficial than those who do not
even possess the concept.
Following these prefatory remarks about the prefatory remarks, let’s get
to work.
I have promised a discourse on the science of knowing, and science of
knowing you expect. What is the science of knowing?
First, in order to start with what everyone would admit, and to speak of
it as others would: without doubt it is one of the possible philosophical sys-
tems, one of the philosophies. So much as an initially stated genus, according
to the rules of definition.
[8–9] So what generally is philosophy and what is it commonly taken to
be; or, as one could more easily say, what should it be according to what is gen-
erally required of it? 11
Without doubt: philosophy should present the truth. But what is truth, 12
and what do we actually search for when we search for it? Let’s just consider
what we will not allow to count as truth: namely when things can be this way
or equally well the other; for example the multiplicity and variability of opin-
ion. Thus, truth is absolute oneness and invariability of opinion. So that I can
let go of the supplemental term “opinion,” since it will take us too far afield, let
me say that the essence of philosophy would consist in this: to trace all multi-
plicity (which presses itself upon us in the usual view of life) back to absolute one-
ness. I have stated this point briefly; and now the main thing is not to take it
superficially, but energetically and as something which ought in all seriousness24
The Science of Knowing
to hold good. All multiplicity—whatever can even be distinguished, or has its
antithesis, or counterpart 13 —absolutely without exception. Wherever even the
possibility of a distinction remains, whether explicitly or tacitly, the task is not
completed. Whoever can point out the smallest distinction in or with regard to
what some philosophical system has posited as its highest principle has refuted
that system.
As is obvious from this, absolute oneness is what is true and in itself
unchangeable, its opposite purely contained within itself. 14 “To trace back”—
precisely in the continuing insight of the philosopher himself as follows: that
he reciprocally conceives multiplicity through oneness and oneness through
multiplicity. That is, that, as a principle, Oneness = A illuminates such multi-
plicity for him; 15 and conversely, that multiplicity in its ontological ground can
be grasped only as proceeding from A.
The science of knowing has this task in common with all philosophy.
All philosophers have intended this consciously or unconsciously; and if one
could show historically that some philosophers didn’t have this objective, then
one can offer a philosophical proof that, to the extent they wished to exist (as
philosophers), they must have intended it. Because merely apprehending mul-
tiplicity, as such in its factical occurrence, is history. Whoever wants this alone
as the absolute one intends that nothing should exist except history. If this
person now says that there is something in addition to history, which he wants
to designate by the name “philosophy,” then he contradicts himself and
thereby destroys his entire statement.
[10–11] Since, as a result, absolutely all philosophical systems must
agree, to the extent that they wish to exist on their own apart from history, the
difference between them, taken initially in a superficial and historical manner,
can only reside in what they propose as oneness, the one true self-contained
in-itself (e.g., the absolute; therefore one could say in passing that the task of
philosophy could be expressed as the presentation of the absolute).
In this way, I say, various philosophies could be differentiated if one
looked at them superficially and historically. But let’s go further. I claim that,
to the extent that general agreement is possible among actually living individ-
uals in regard to any manifold, to that extent the oneness of principle is in
truth and in fact one. For divergent principles become divergent results, and
consequently yield thoroughly divergent and mutually incoherent worlds, so
that no sort of agreement about anything is possible. But if one principle alone
is right and true, it follows that only one philosophy is true, namely the one
which makes this one principle its own, and all others are necessarily false.
Therefore, in case there are several philosophies simultaneously presenting
different absolutes, either all, or all except one, are false. 16
Further, this significant consequence also follows: since there is only one
absolute, a philosophy which has not made this one true absolute its own sim-First Lecture
25
ply doesn’t have the absolute at all but only something relative, a product of
an unperceived disjunction which for this very reason must have an opposing
term. Such a philosophy leads all multiplicity back, not to absolute oneness as
the task requires, but only to a subordinate, relative oneness; and thus it is
refuted and shown in its insufficiency not just by the true philosophy but even
by itself, if only one is acquainted with philosophy’s true task and reflects it
more prudently than this system does. Therefore the entire differentiation of
philosophies according to their principles of oneness is only provisional and
historical, but cannot in any way be adequate by itself. However, since we must
start here with provisional and historical knowledge, let’s return to this prin-
ciple of classification. Again, the science of knowing may be one of the possible
philosophies. Therefore, if it makes the claim, as it already has, that it resem-
bles no previous system but is completely distinct from them all, new, and self-
sufficient, then it must have a different principle of oneness from all the rest.
What did these have as a principle of oneness? In passing, let me note that it
is not my intention here to discuss the history of philosophy and to let myself
in for all the controversies which would be aroused for me in this way, instead
I simply intend to progress gradually in developing my own concept. For this
purpose, what I will say could serve as well if it were assumed arbitrarily and
were not historically grounded, as if it were historically true. This can be
abundantly demonstrated if such demonstration is necessary and [if ] there are
people concerned about this. I claim that this much is evident from all
philosophies prior to Kant, the absolute was located in being, that is, in the
dead thing as thing. The thing should be the in-itself. (In passing: [12–13] I
can add that, except for the science of knowing, since Kant, philosophers
everywhere without exception, the supposed Kantians as well as the supposed
commentators on and improvers of the science of knowing, have stayed with
the same absolute being, and Kant has not been understood in his true, but
never clearly articulated principle. Because it is not a matter of what one calls
being, but of how one grasps and holds it inwardly. For all that one names it
[i.e., the absolute] “I,” if one fundamentally objectifies it, and separates it from
oneself, then it is the same old thing-in-itself ).—But surely everyone who is
willing to reflect can perceive that absolutely all being posits a thinking or con-
sciousness of itself; 17 and that therefore mere being is always only one half of a
whole together with the thought of it, and is therefore one term of an origi-
nal and more general disjunction, a fact which is lost only on the unreflective
and superficial. Thus, absolute oneness can no more reside in being than in its
correlative consciousness; [14–15] it can as little be posited in the thing as in
the representation of the thing. Rather, it resides in the principle, which we
have just discovered, of the absolute oneness and indivisibility of both, which
is equally, as we have seen, the principle of their disjunction. We will name this
principle pure knowing, knowing in itself, and, thus, completely objectless26
The Science of Knowing
knowing, because otherwise it would not be knowing in itself but would
require objectivity for its being. It is distinct from consciousness, which posits
a being and is therefore only a half. This is Kant’s discovery, and is what makes
him the founder of Transcendental Philosophy. Like Kantian philosophy, the
science of knowing is transcendental philosophy, and thus it resembles Kant’s
philosophy in that it does not posit the absolute in the thing, as previously, or
in subjective knowing—which is simply impossible, because whomever
reflects on this second term already has the first—but in the oneness of both.
But now, how does the science of knowing differentiate itself from Kan-
tianism? Before I answer, let me say this. Whoever has caught a genuine inner
glimpse of just this higher oneness has already achieved in this first hour an
insight into the true home for the principle of the sole true philosophy, which
is nearly entirely lacking in this philosophical era; and he has acquired a con-
ception of the science of knowing and an introduction to understanding it,
which has also been wholly lacking. [16–17] Namely, as soon as one has heard
that the science of knowing presents itself as idealism, one immediately infers
that it locates the absolute in what I have been calling thinking or consciousness
which stands over against being as its other half and which therefore can no
more be the absolute than can its opposite. Nevertheless, this view of the sci-
ence of knowing is accepted equally by friends and enemies and there is no
way to dissuade them from it.
In order to find a place for their superiority at improving things, the
improvers have switched the absolute from the term in which, according to
their view, it resides in the science of knowing to the other term to which they
append in addition the little word “I” which may well be the single net result
of Kant’s life and, if I may name myself after him, of my life as well, which has
also been devoted to science.C H A P T E R
2
Second Lecture
Wednesday, April 18, 1804
Honorable Guests:
We will begin today’s lecture with a brief review of the previous one. In con-
ducting this review I have an ulterior motive, namely to adduce what in gen-
eral can be said about the technique of fixing lectures like this in memory and
of reproducing them for oneself, and to discuss the extent to which transcrib-
ing is or is not useful. “In general,” I say, because in what concerns memory
and the possibility of directing the attention simultaneously to a number of
objects, one finds a great discrepancy among people; and I in particular am
among the least fortunate in this regard, since I am utterly lacking in what is
usually called memory, and my attention is capable of taking up no more than
one thing at a time. For this reason my recommendations are all the less
authoritative, and each of you must decide for yourselves how far they apply
to you and how you might employ them.
For me, the proper and favorite listener would be the one who is able to
reproduce the lecture for himself at home, not literally {unmittlebar} for that
would be mechanical memory, but by pondering and reflection; and indeed
following the course of the argument absolutely freely, moving backwards
from the results with which we concluded to the premises, forwards from the
premises with which we started deducing the result, and moving both ways at
once from the center. Further, this listener should be able to do so in absolute
independence from specific modes of speech; and, since we propose to con-
duct what is really only one whole, self-contained presentation of the science
of knowing in many lectures separated by hours and days, and with the single
lecture periods composing the integrative parts of this whole, just as the sin-
gle minutes of the lecture hour make up its parts, once again my favorite lis-
tener will be the one who is able to present all the single lecture sessions com-
prehensively as a whole, whether beginning with the first, or with the last he
has heard, or with any of those in the middle. This is the first point. 1
This lecture begins at GA II, 8, pp. 18–19.
2728
The Science of Knowing
Now, in the second place, what is most noteworthy in each lesson for
each person, and therefore what each grasps most surely, is whatever new ele-
ment that person learns and clearly understands. What we genuinely compre-
hend becomes part of ourselves, and if it is a genuinely new insight, it pro-
duces a personal transformation. It is impossible that one not be, or that one
cease to be, what one has genuinely become, and for exactly this reason the
science of knowing can, [20–21] more than any other philosophy, commit
itself to reactivate the dormant instinct for thinking, because it introduces new
concepts and insights. 2
What was new in the last hour for those who are not familiar with the
science of knowing, and which may have appeared in a new light for those
already acquainted with it, was this: if one reflects properly, absolutely all
being presupposes a thought or consciousness having that being as its object;
that consequently being is a term of a disjunction, the other half of which is
thinking; and that for this reason oneness is not to be found in the one or the
other but in the connection of both. 3 Oneness thus is the same as pure know-
ing in and for itself, and therefore it is knowledge of nothing; or in case you
find the following expression easier to remember, it (oneness) is found in truth
and certainty, which is not certainty about any particular thing, since in that
case the disjunction of being and knowing is already posited. So if, in the
effort to reproduce the first lecture from within, someone had clearly and
vividly hit upon just this single point, then it would have been possible with a
little logical reflection to develop all the rest from it. For example, the listener
might inquire: How did we get to the point of proving that being had a cor-
relative term? Did this arise in a polemical context? Was being taken to be
absolute oneness rather than a correlative term? Then each would have
recalled that this was the case up until Kant. He could then have asked him-
self: But how generally did we come to investigate what might or might not
be absolute oneness? Thus each would recall just from knowing why he
attended this lecture series that we were supposed to be doing philosophy and
that the essence of philosophy was said to be the assertion of absolute oneness
and the reduction of all multiplicity to it, and so the whole thought process
can be laid out without any problem: what is the science of knowing, etc. 4
But this restoration must not lack depth and thoroughness. For exam-
ple, “reducing multiplicity to oneness” is a brief formula, easy to remember, and
it is comfortable to use it in answering the common question “What is phi-
losophy?” which is a question to which one usually doesn’t quite know how to
respond. “But,” one asks oneself, “do you really know what you are saying: can
you clarify it to yourself to the extent of providing a lucid and transparent con-
struction of it? Has it been described? [22–23] How has it been described?
With such and such words. All right, the lecturer said this: and these are
words! 5 I will construct it.” Or—this thing which is neither being nor con-Second Lecture
29
sciousness but rather the oneness of both, which has been presented as the
absolute oneness of the transcendental philosophy, can be indicated with these
words. But this cannot yet be completely clear and transparent to you, because
all philosophy is contained in the transparency of this oneness, and from here
on we will do nothing else but work on heightening the clarity of this one
concept, in which at one stroke I have given you the whole. If it were already
completely clear to you, you would not need me any longer. Nevertheless, 6 I
would add that everyone must possess more than the bare empty formula.
One must also have a living image of this oneness which is firmly fixed and
which never leaves him. With my lectures I attend to this, your fixed image;
we will extend and clarify it together. If someone does not have this image,
there is no way for me to address them, and for this person my whole dis-
course becomes talk about nothing, since in fact I will discuss nothing except
this image.
And so that I finally say directly what it all comes down to—without this
free, personal re-creation of the exposition of the science of knowing in its living pro-
fundity, which I have just mentioned, one will have no possible use for these lec-
tures. “The subject cannot remain simply in the form in which I express myself
here”; even though you of course can recall it by, and out of, yourselves in just
this form. In short, a middle term must come between my act of exposition
and your active mastery of what has been expounded, and that term is your
own rediscovery; otherwise, everything ends with the act of exposition and you
never attain active mastery. 7
It is quite unimportant whether one undertakes this re-creation with
pen in hand (as I, for example, would do it because I have no memory but
rather an imagination, which can only be held in check by written letters) or
whether someone with more memory and a tamer imagination does it in free
thought. It is essential only that each one does it in the individually appropri-
ate way; and, in no case can re-creating it in writing cause harm. 8
In the light of what has been said, whatever is noted down during the
lecture cannot take the place of a proper re-creation; instead it can only serve
as resource for the act, which must be undertaken with or without this [24–25]
assistance. With the deliberate speech, considerable pauses after important
paragraphs, and the repetitions of significant phrases which I am using, it must
certainly be possible for your pens to catch in passing the main points of the
lecture for the required task. For my part, if I had to attend lectures like these,
I would not even begin to try to note down more than such main points, since
while writing I cannot listen energetically, and while listening energetically I
cannot write; for me it would be more a matter of the whole living lecture
rather than the isolated, dead words, and more especially a matter of the sel-
dom noticed, but very genuine and real, physical/spiritual action of keen think-
ing carried out in my presence. Yet I am completely certain that it can be quite30
The Science of Knowing
different with others and that more easily activated spirits could indeed do two
things at the same time equally well.
So much once and for all on this topic!—Now let’s carry on with the
investigation begun last time, that is, with provisionally answering the ques-
tion: What is the science of knowing? 9
All transcendental philosophy, such as Kant’s (and in this respect the
science of knowing is not yet different from his philosophy), posits the
absolute neither in being nor in consciousness but in the union of both. Truth
and certainty in and for itself = A. Thus it follows (and this is another point 10
through which today’s talk links up with the previous one and by means of
which in the general reproduction of all the lectures, the previous one is to be
produced from it and it from the previous one) . . . it follows, I say, that for
this kind of philosophy the difference between being and thinking, as valid in
itself, totally disappears. Indeed, everything that can arise in such a philoso-
phy is contained in the epiphany 11 which we already consummated in ourselves
during the last lecture; in the insight that no being is possible without think-
ing, and vice versa—wholly being and thinking at the same time, 12 and nothing
can occur in the manifest sphere of being without simultaneously occurring in
the manifest sphere of thinking, if one just considers rightly rather than sim-
ply dreaming, and vice versa. Thus in the vision, 13 which is given and granted,
there is nothing primarily of concern for us as transcendental philosophers. 14
But according to our insight that the absolute is not a half but indivisible one-
ness, an insight which reaches beyond all appearance, it is absolutely and in
itself neither being nor thinking, but rather: 15
A
/ \
A.—B T—
If now (in order [26–27] to apply what has just been said and to make
it even clearer it is assumed that, in addition to its absolute, fundamental divi-
sion into B and T, A also divides itself into x, y, z, then it follows:
1. that everything together in and for itself, including absolute A, and x, etc.,
is only just a modification of A; from this follows directly
2. that it all must occur in B as well as in T. 16
Assume also that there is a philosophical system that has no doubts that
the dichotomy of B and T, which arises from A, is mere appearance, and that
therefore it would be a genuine transcendental philosophy; but assume further
that this philosophy remained caught in such an absolute division of A into x,
y, z, just as we proposed. Thus, this system, for all its transcendentalism, wouldSecond Lecture
31
still not have penetrated through to pure oneness, nor would it have com-
pleted the task of philosophy. Having eluded one disjunction, it would have
fallen into another; and despite all the admiration one must show it for first
discovering the primordial illusion, nevertheless with the discovery of this new
disjunction, it is refuted as the true, fully accomplished, philosophy.
The Kantian system is exactly this very one, precisely characterized by
the outline I have just given. If Kant is studied, not as the Kantians without
exception have studied him (holding on to the literal text, which is often clear
as the heavens, but is also often clumsy, even at important times), but rather
on the basis of what he actually says, raising oneself to what he does not say
but which he must assume in order to be able to say what he does, then no
doubt can remain about his Transcendentalism, understanding the word in the
exact sense I have just explained. Kant conceived A as the indivisible union of
being and thinking.
But he did not conceive it in its pure self-sufficiency in and for itself, as
the science of knowing presents it, but rather only as a common basic determi-
nation or accident of its three primordial modes, x, y, z—these expressions are
meaningful; it cannot be said more precisely—as a result of which for him
there are actually three absolutes and the true unitary absolute fades to their
common property.
The way his decisive and only truly meaningful works, the three cri-
tiques, come before us, Kant has made three starts. In The Critique of Pure Rea-
son, his absolute (x) is sensible experience, and [28–29] in that text he actually
speaks in a very uncomplimentary way about Ideas and the higher, purely spir-
itual world. From earlier writings and a few casual hints in this Critique itself,
one might conclude that in his own view the matter didn’t end there. But I
would commit myself to showing that these hints are only one more inconsis-
tency, since if we correctly follow the implications of the premisses stated in
that text, then the supersensible world must totally disappear, leaving behind a
mere noumenon which has its complete realization in the empirical world.
Moreover, he has the correct concept of the noumenon, and by no means the
confused Lockean notion which his followers have imposed on him. The high
inner morality of the man corrected the philosopher, and he published the Cri-
tique of Practical Reason. In this text the I comes to light as something in itself
through the inherent categorical concept as [30–31] it could not possibly do in
the Critique of Pure Reason, which is solely based on, and drawn from, what is
empirically; and thus we get the second absolute, a moral world = z. Still, not
all the phenomena that are undeniably present in self-observation have been
accounted for; there still remains the notions of the beautiful, the sublime, and
the purposive, which are evidently neither theoretical cognitions nor moral
concepts. Further, and more significantly, with the recent introduction of the
moral world as the one world in itself, the empirical world is lost, as revenge32
The Science of Knowing
for the fact that the latter had initially excluded the moral world. And so the
Critique of Judgment appears, and in its Introduction, the most important part
of this very important book, we find the confession that the sensible and super-
sensible worlds must come together in a common but wholly unknown root,
which would be the third absolute = y. I say a third absolute, separate from the
other two and self-sufficient, despite the fact that it is supposed to be the con-
nection of both other terms; and I do not thereby treat Kant [32–33] unjustly.
Because if this y is inscrutable, then while it may indeed always contain the con-
nection, 17 I at least can neither comprehend it as such, nor collaterally conceive
the two terms as originating from it. If I am to grasp it, I must grasp it imme-
diately as absolute, and I remain trapped forever, now as before, in the (for me
and my understanding) three absolutes. Therefore, with this final decisive addi-
tion to his system, Kant did not in any way improve that which we owe to him,
he only generously admitted and disclosed it himself.
Let me now characterize the science of knowing in this historical move-
ment, from which, to be sure, my speculations, which are independent from
Kant, take their origin. Its essence consists first in discovering the root (indis-
cernible for Kant) in which the sensible and supersensible worlds come
together and then in providing the actual conceptual derivation of both worlds
from a single principle. The maxims which Kant so often repeated orally and
in writing, which his followers parroted (we must stop short again 18 and can-
not go farther), and which pre-Kantian dogmatists, too, could have used to
answer Kant (we must stand by our dogmatism and cannot go farther), are
here completely rejected as maxims of weakness or inertia, which are then
taken to apply to everyone. The science of knowing’s own maxim is to admit
absolutely nothing inconceivable and to leave nothing unconceived; and it is
satisfied to wish not to exist if something is pointed out to it which it hasn’t
grasped, since it will be everything or nothing at all. To avoid all misunder-
standing let me add that if it too must finally admit something inconceivable,
then at least it will conceive it as just what it is, i.e., absolutely inconceivable,
[34–35] and as nothing more; and thus too it will conceive the point at which
absolute conceiving is able to begin. Let this much suffice as an historical
characterization of the science of knowing vis-à-vis the sole neighbor against
which it is immediately juxtaposed and to which it can be compared: Kantian
philosophy. It cannot be directly compared to any earlier philosophies or
recent afterbirths, since it shares nothing with them and is totally different. It
shares the common genus of transcendentalism with Kantianism alone, and to
that extent must demarcate itself from it, a demarcation which has to do sim-
ply with the clarity of this property but not at all with its vain reputation.
Let me give this characterization in pure concepts, higher than, and
independently of, historical setting, and present its schema: A is admitted; it
divides itself both into B and T and simultaneously into x, y, z. Both divisionsSecond Lecture
33
are equally absolute; one is not possible without the other. Therefore, the
insight with which the science of knowing begins and which constitutes its
distinction from Kantianism is not at all to be found in the insight into the
division into B and T which we already also completed in the last lecture, nor
in the insight into the division of x, y, z which we have still not finished but
have problematically presupposed; rather it is to be found in the insight into
the immediate inseparability of these two modes of division. Therefore the two
divisions cannot be seen into {eingesehen} immediately, as it seemed from the
outset until now; instead they can be seen into only mediately through the
higher insight into their oneness.
I call the attention, especially of the returning listeners, to this most
important clue; there you have in its full simplicity a characterization of our
speculation much earlier and right from the beginning which did not come
into our first lecture series until the middle of our work. 19
(This schema is a summary of the whole lecture. How to re-create the
whole out of it is indicated above.)C H A P T E R
3
Third Lecture
Thursday, April 19, 1804
Honorable Guests:
First let me clarify a point raised at the end of the last lecture, which might
occasion misunderstanding. Absolute A is divided into B and T and into x, y,
z; all at one stroke. As into x, y, z, so into B and T; as into B and T, so into x,
y, z. But how have I expressed myself just now? Once commencing with x and
the other time commencing with B. This is just a perspective, a bias in my
speaking. I certainly know—and even expressly assert—that implicitly, beyond
the possibilities of my mode of expression and my discursive construction,
both are totally identical, completely comprehended in one self-contained
stroke. Therefore, I am constructing what cannot be constructed, with full
awareness that it cannot be constructed.
Let me continue now to characterize the science of knowing on the
basis of the indicators found in comparing it with Kantian transcendental-
ism—among other things, I said that Kant very well understood A as the link
between B and T, but that he did not grasp it in its absolute autonomy.
Instead, he made it the basic common property and accident of three
absolutes. In this the science of knowing distinguishes itself from him. There-
fore this science must hold that knowing (or certainty), as soon as we have
characterized it as A, must actually be a purely self-sustaining substance; that we
can realize it as such for ourselves; and it is in just this realization that the gen-
uine realization of the science consists. (A “genuine realization,” I say, with
which we aren’t yet concerned here, since we are still occupied with stating the
simple concept, which is not the thing itself.) This constitutes today’s thesis. 1
To begin with, we can demonstrate immediately that knowing can actu-
ally appear as something standing on its own. I ask you to look sequentially at
your own inner experience: 2 if you remember it accurately, you will find the
object and its representation, with all their modifications. But now I ask fur-
ther: Do you not know in all these modifications; and is not your knowing, as
This lecture begins at GA II, 8, pp. 36–37.
34Third Lecture
35
knowing, the same self-identical knowing in all variations of the object? As
surely as you say, “Yes,” to this inquiry (which you will certainly do if you have
carried out the given task), so surely will knowledge manifest and present itself
to you (as = A) whatever the variation of its objects (and hence [38–39] in total
abstraction from objectivity); remaining as (= A); and thus as a substantive, as
staying the same as itself through all change in its object; and thus as oneness,
qualitatively changeless in itself. This was the first point. 3
Thus it presents itself to you with impressive, absolutely irrefutable
manifestness. You understand it so certainly that you say: “It is absolutely thus,
I cannot conceive it differently.” And if you were asked for reasons, you could
refuse the request and still not give up this contention. It is manifest to you as
absolutely certain. During all possible variation in the object, you have said,
knowing always remains self-identical. Have you then run through and
exhausted all possible changes in object, testing in each case whether know-
ing remains the same? I do not think so, because how could you have done it?
Therefore, this knowing manifests itself independently of such experiments
and completely a priori as self-sustaining and self-identical independent from
all subjectivity and objectivity.
1. Now, note what actually belongs to this substantial knowing, so conceived,
and do so with the deepest sincerity of self-consciousness, so that the erro-
neous view of the science of knowing which was criticized at the end of our
first lecture, namely that it locates the Absolute in the knowing which
stands over against its object, doesn’t arise again here. It is true that in our
experiment we have started with this consciousness or presentation of an
object. T.B.–T.B. 4 and so on. In this part of the experiment, B made the T
different in every new moment, because the T was altogether nothing else
than the T for this B, and disappears with it. Now when we raised ourselves
to the second part by asking, “Is not knowing one and the same through-
out?” and finding it to be so, we raised ourselves above all differences of T
as well as of B. Therefore we could express ourselves much more accurately
and precisely: knowing, 5 which for this reason is not subjective, is
absolutely unalterable and self-identical not just independently from all
variability of the object, but also independently from all variability of the
subject without which the object doesn’t exist. The changeable is nothing
further, neither the object nor the subject, [40–41] but just the mere pure
changeable, and nothing else. Now this changeable in its continuing
changeability, which is itself unchangeable, divides itself into subject and
object; and the purely unchangeable, in which the division of subject and
object falls away, as does change, opposes itself to the changeable.
2. Here has been disclosed a splendid example of an insight that comes from
exhaustive, continuous searching, which cannot be derived from experience,36
The Science of Knowing
but which rather is absolutely a priori. And so past experiences obligate me
to entreat everyone here to whom this insight has really been evident (which
I think is the case with all of you given the simplicity of the task) to keep
this very example in mind, to hold on to it, and, if the old empiricist demon
shows up, to attack it and send it away promptly, until we succeed in com-
pletely exterminating it. I would very much like to be spared the eternal
struggle about whether in general there is manifestness or something a pri-
ori (for both are the same). Individuals come to the conclusion that this is
the case only by producing it somehow or other in themselves. This has
happened here today and I simply ask that you not forget it.
Result: knowing, in the mentioned meaning as A, has actually appeared
to us as self-sufficient, as independent of all changeability, and as self-same and
self-contained oneness, as was presupposed by previous historical reports of the
science of knowing. We therefore already seem to have realized the principle of
the science of knowing in ourselves and to have penetrated into it.
The second advance in today’s lecture is this. 6 We only seem to have
done this, but this is an empty seeming. We see merely {sehen bloss ein} that it
is so, but we do not see into what it authentically is as this qualitative oneness.
Precisely because we see into only such a that, we are trapped in a disjunction
and thus in two absolutes, changeability and unchangingness, to which we
might possibly append a third, the undiscoverable root of both, and thus end
up in the same shape as Kant’s philosophy. The ground of this duality, insur-
mountable in this way, is as follows: the that must seem self-creating, just as
our recent insight seemed to be. But this appearance is possible only under the
condition that a point of origin {terminus a quo} appears, which seems (as
opposed to this self-creation) to be produced by us, just as the [42–43] first
part of our previously conducted experiment actually and in fact appeared to
be. In a word, we grasp both changeability and unchangingness equally, and
are inwardly torn into two or three immediate terms. How then is this to be?
Obviously, it is clear without any further steps both that one of the two would
have to be grasped mediately, and that this term which is grasped mediately
cannot be unchangingness, which as the absolute can only be realized
absolutely, but rather must be changeability. The unchangeable would have to
be intuited not only in its being, which we have already done, but instead it
would have to be penetrated in its essence, its one absolute quality. It (the
unchangeable) would have to be worked through in such a way that change-
ability would be seen as necessarily proceeding from it and as mediately gras-
pable only by means of it.
Briefly, clearly, and to fix the point easily in memory: the insight that
knowing is a self-sustaining qualitative oneness (an insight which is purely
provisional and belongs to a theory of the science of knowing) leads to theThird Lecture
37
question, “What is it [i.e., knowing] in this qualitative oneness?” The true
nature of the science of knowing resides in answering this question. In order
to analyze this even further, it is clear that for this purpose one must inwardly
construct this essence of knowing. Or, as in this case is exactly the same thing,
this essence must construct itself. In this constructive act, it is without any
doubt and is what it is as existing; and, as existing, it is what it is. Therefore it
is clear that the science of knowing and the knowing that presents itself in its
essential oneness are entirely one and the same; that the science of knowing
and primordially essential knowing merge reciprocally into one another and
permeate each other; that in themselves they are not different; and that the
difference which we still make here is only a verbal difference of the exact kind
mentioned at the start of this lecture. The primordially essential knowing is
constructive, thus intrinsically genetic; this would be the original knowing or
certainty in itself. Manifestness in itself is therefore genetic.
And with this we have specified the deepest characteristic difference
demarcating the science of knowing from all other philosophies, particularly
from the most similar, the Kantian. All philosophy should terminate in know-
ing in and for itself. Knowing, or manifestness in and for itself, is actively
genetic. The highest appearance of knowing, which no longer expresses its
inner essence but instead just its external existence, is factical; and since it is
still the appearance of knowing, factical manifestness. All factical manifest-
ness, even if it is the [44–45] absolute, remains something objective, alien, self-
constructing but not constructed of knowing, and therefore something
inwardly unexplained, which an exhausted speculation, skeptical of its own
power, calls inexplicable. Kantian speculation ends at its highest point with
factical manifestness: the insight that at the basis of both the sensible and
supersensible worlds, there must be a principle of connection, thus a thor-
oughly genetic principle which creates and determines both worlds absolutely.
This insight, which is completely right in itself, could occur to Kant only as a
result of his reason’s absolute, but unconsciously operative law: that it [that is,
his reason] come to a stop only with absolute oneness, recognize only this as
absolutely substantial, and derive everything changeable from this one. This
basic law of oneness remains only factical for him, and therefore its object is
unexamined, because he allows it to work on him only mechanically; but he
does not bring this action itself and its law into his awareness anew. If he did
so, pure light would dawn on him and he would come to the science of know-
ing. Kantian factical manifestness is not even the highest kind, because he lets
its object emerge from two related terms and does not grasp it as we have
grasped the highest factical object, namely as pure knowing; instead he grasps
it with the qualification that it is the link between the sensible and supersen-
sible worlds. That is, he does not grasp it inwardly and in itself as oneness, but
as duality; his highest principle is a synthesis post factum. 7 Namely, this means38
The Science of Knowing
a case when by self-observation one discovers in one’s consciousness two terms
of a disjunction and, compelled by reason, sees that they must intrinsically be
one, disregarding the fact that one cannot say how, given this oneness, they
can likewise become two. Briefly, this is exactly the same procedure by which
in our first lecture we rose from the discovery of the duality of being and
thinking to A as their required necessary connection, in order initially to con-
struct for ourselves the transcendentalism common to the science of knowing
and Kantianism, but the matter was not to rest there. Additionally there
should be a synthesis a priori which is equally an analysis, since it simultane-
ously provides the basis for both oneness and duality.—Kant’s highest mani-
festness, I said, is factical, and not even the highest factical kind. The highest
factical manifestness has been presented today: the insight into knowing’s
absolute self-sufficiency, without any determination by anything outside itself,
anything changeable. This is contrary to the Kantian absolute, which [46–47]
is determined by the transition between the sensible and supersensible. Since
now this presently factical element in science is itself to become actively
genetic and developing, then, change in general will be grounded in it as a
genesis pure and simple. But by no means will any particular change be so
grounded. It seems that absolute facticity could be discovered only by those
who have raised themselves above all facticity, as I have actually {in der That}
discovered it and consistently made use of it only after discovering the true
inner principle of the science of knowing, and as I am using it now to lead the
audience from that point forward in the genetic process {in die Genesis}.
Kant’s manifestness is factical, we ourselves are presently also standing
in facticity; and I add that everywhere in the scientific world there is no other
kind of manifestness except the factical (namely in the first principles), except
in the science of knowing. As far as philosophy is concerned, after conducting
the demonstration with Kant, we can safely omit tests on other systems. After
philosophy, mathematics makes claims to manifestness, indeed in some of its
representatives it takes on airs by elevating itself above philosophy, an error
which can be excused in the light of today’s philosophical eclecticism. Now
abstracting here completely from the fact that things are not so wonderful for
mathematics—not even in regard to how it can and should be—this science
must confess that its principles admit nothing more than factical manifest-
ness, regardless of the fact that they will become actively genetic as we pro-
ceed. For let the arithmetician qua arithmetician simply tell me how he is able
to elicit a solid and permanent number one; or let [48–49] the geometer
explain what fixes and holds his space for him while he draws his continuous
lines through it; and whether these and ever so many other ingredients which
he needs for the possibility of his derivations are given in any other way than
through factical intuition. Of course this does not in any way constitute an
objection to mathematics; as mathematics it can and should be nothing else.Third Lecture
39
It is certainly not our business to obscure the boundaries of the sciences; but
one should simply recognize, and this science, like all the others, should know
that it is neither the first nor self-sufficient but rather that the principles of its
possibility lie in another, higher science.
Now, since in the actual sciences generally no other principles are avail-
able than those that are factically manifest, and since by contrast the science
of knowing intends to introduce entirely genetic manifestness and then to
deduce the factical from it, it is clear that essentially in its spirit and life the
science of knowing is wholly different from all previous scientific employment
of reason. It is clear that it is not known to anyone who has not studied it
directly and that nothing can take the place of such study. It is equally clear
that there is no perspective or premise which has appeared in previous life or
science from which it can be seen as true, attacked, or refuted, because what-
ever this perspective or premise is, and however certain it might be, still it is
nonetheless surely only factically manifest, and this science accepts nothing of
this kind unconditionally, but does so only under conditions which it deter-
mines in its genetic analysis. But whoever wants to argue against the science
of knowing using such a perspective as its principle, wants unconditional agree-
ment which is already once and for all ruled out in advance. Therefore he is
arguing from a premise that has not been accepted {ex non concessis} and makes
himself ridiculous. The science of knowing can only be judged internally, it
could be attacked and refuted only internally, by pointing out an internal con-
tradiction, an inner inconsistency or insufficiency. Therefore such activity
must be preceded by study and comprehension and must begin with that.
Until now to be sure the opposite order has been tried; first judging and refut-
ing and after that, God willing, understanding. As a result nothing has ensued
except that the blows have [50–51] completely missed the science of knowing,
which has remained hidden from view like an invisible spirit, and they have
struck instead the chimeras which these men have created with their own
hands. Following this, they have gone so far astray with these fantasies and
have spread confusion so extensively that today it can be expected that they
would at least understand that they are confused!C H A P T E R
4
Fourth Lecture
Friday, April 20, 1804
Honorable Guests:
It seems to me that we have succeeded, even in the prolegomenon, in gaining
a very clear and deep insight into the scientific form of the science that we
wish to pursue here. Let us continue the observations with which we have
achieved this insight.
Here are the results so far: that we have certainly grasped knowing as
unchangeable, self-same and self-sustaining, beyond all change and beyond the
subjectivity-objectivity, which is inseparable from change. But this insight was
not yet the science of knowing itself, but rather only its premise. The science
of knowing must still actually construct this inner, qualitatively unchangeable
being, and as soon as it does this, it will simultaneously create change, the sec-
ond term, as well.
The true authentic meaning of this simultaneous double construction of
the changeable and the unchangeable will become completely clear only when
we actually and immediately carry out the construction, something that
belongs within the sphere of the science of knowing itself, certainly not in pre-
liminary reports. Misunderstandings about this are unavoidable at the begin-
ning. In order to come as close as possible to complete accuracy right from the
start, I venture on a question that has already been raised.
On introducing the schema: 1
A
x y z • B–T
I said that the science of knowing stood in the point. 2 I have been asked
whether it doesn’t rather go in A. The most exact answer is that actually and
strictly it doesn’t belong in either of these but rather in the oneness of both.
By itself A is objective and therefore inwardly dead; it should not remain so,
This lecture begins at GA II, 8, pp. 52–53.
40Fourth Lecture
41
indeed it should become actively genetic. 3 The point, on the other hand, is
merely genetic. Mere 4 genesis is nothing at all; but this is not just mere gen-
esis but the determinate genesis that is required by the absolute qualitative
A; [it is] a point of oneness. Now to be sure this point of oneness can be real-
ized immediately, oscillating and expending itself in this point; and we, as sci-
entists of knowing, 5 are 6 this realization inwardly (I say inwardly and con-
cealed from ourselves). But this point can neither be expressed nor
reconstructed in its immediacy, since all expression or reconstruction is con-
ceiving and is intrinsically mediated. It is expressed and reconstructed
[54–55] just as we have expressed it at this moment: namely, that one begins
from A and, indicating that it cannot persist alone, links it to the point; or
one begins with the point and, indicating that it cannot persist alone, links
it to A, all the while, to be sure, knowing, saying, and meaning that neither
A nor the point can exist by itself and that all our talk could not express the
implicit truth, but instead that the implicitly unreconstructible something
which can only be pictured in an empty and objective image is the organic
oneness of both. Thus, since reconstruction is conceiving, and since this very
conceiving explicitly abandons its own intrinsic validity, this is precisely a
case of conceiving the inconceivable as inconceivable. Therefore—and I put
this here first as a clarification—this is a question of the organic split into
B–T and x, y, z, as I explained at the beginning of the last lecture; because
when I speak, I must always put one before the other. But is it actually so?
No, instead it is exactly the same stroke; and let me add this as well: this
deeper connection must indeed itself be a result and a lower expression of the
higher one just now described. Finally, in order to say this in its full meaning
and thereby to make your insight into the science of knowing, and into
knowing as such, much clearer: secondary knowing, or consciousness, with its
whole lawful play by means of fixed change and the manifold (within it or
outside it), of sensible and supersensible, and of time and space, comes to be,
in principle, through this recently noted and demonstrated division, taken
merely as a division and nothing else. Everything we attribute to the sub-
ject, as originating from it, derives from this. Because it is clear without fur-
ther ado that from a particular perspective, namely the science of knowing’s
synthetic perspective, the disjunction must be just as absolute as oneness;
otherwise we would be stuck in oneness and would never get outside it to
changeability. (Let me note in passing that this is an important characteris-
tic of the science of knowing and distinguishes it, e.g., from Spinoza’s sys-
tem, which also wants absolute oneness but does not know how to make a
bridge from it to the manifold; and, on the other hand, if it has the mani-
fold, cannot get from there to oneness.)
[56–57] As scientists of knowing, we never escape the principle of divi-
sion inwardly and empirically (i.e., by means of what we do and promote); but42
The Science of Knowing
we certainly escape it intellectually with regard to what is valid in itself, in
which very regard the principle of division surrenders and negates itself.
Or, so that I make the point at which we have arrived even clearer: since
we actually reason as we have just been doing, where does our reasoning stand,
if we remain exclusively {nur recht} in consciousness {Besinnung}? (“If we remain
exclusively in consciousness,” I say, because we can also lose ourselves in the
intelligible realm, and there is even, in its place, an art of consciously losing one-
self in it.) Manifestly in our construction by means of the principle of division,
it stands in the place not of that which is to be valid to the extent it is constructed
but of that which is intrinsically valid; thus it stands wholly autonomously, as has
been said, between the two principles of oneness and separation, simultaneously
annulling both and positing both. Thus, the standpoint of the science of know-
ing, which stands still in consciousness, is by no means a synthesis post factum;
but instead a synthesis a priori, taking neither division nor oneness as given, but
creating both at one stroke. Once again to adduce an even higher perspective
what is the absolute oneness of the science of knowing? Neither A nor the point,
but instead the inner organic oneness of both. Besides this given description of
the point of oneness is there also another? None at all, we have seen. Therefore,
this description is the original and absolutely authentic one. What are its con-
stituents? The organic oneness of both is a construction or a concept, and indeed
the single absolute concept, abstracted from nothing existing, since even its own
separate existence, and hence the existence of everything conceptual, is denied.
Further, the construction as such is denied by the manifestness of what exists
autonomously; thus even the inconceivable, as the inconceivable and nothing
more, is posited by this manifestness, posited through the negation of the
absolute concept, which must be posited just so that it can be annulled. 7
And so:
1. the necessary unification and indivisibility of the concept and the incon-
ceivable is clearly seen into {eingesehen}, and the result may be expressed
thus: if the absolutely inconceivable is to be manifest as solely self-sustain-
ing, then the concept must be annulled, but to be annulled, it must be
posited, because the inconceivable becomes manifest only with the nega-
tion of the concept. Supplement: hence inconceivable = unchangeable; con-
cept = [58–59] change. Therefore along with the foregoing it is evident that
if the unchangeable is to appear, there must be change.
2. Now to be sure inconceivability is only the negation of the concept, an
expression of its annulment. Therefore it is something which originates
from both the concept and knowing themselves; it is a quality transferred
by means of absolute manifestness. Noting this, and therefore abstracting
from this quality, nothing remains for oneness except absoluteness or pure
self-sufficiency in itself.Fourth Lecture
43
3. The following consideration makes this particularly important and rele-
vant: What is pure self-sufficient knowing in itself? The science of know-
ing has to answer this question, or, as we put it more precisely, it has to
construct the presupposed inner quality of knowing. We are undertaking
this construction here: negating the concept by means of manifestness,
and thus the self-creation of inconceivability is this living construction of
knowing’s inner quality. This inconceivability itself originates in the con-
cept and in pure immediate manifestness; likewise the whole quality of the
absolute, as well as the fact that a quality can be applied to it at all, origi-
nates in the concept. 8 The absolute is not intrinsically inconceivable, since
this makes no sense; it is inconceivable only when the concept itself tries
for it, and this inconceivability is its only property. Having recognized this
inconceivability as an alien quality introduced by knowing, I said before,
only pure self-sufficiency, or substantiality, remains in the absolute; and it
is quite true that at best this self-sufficiency does not originate in the con-
cept, since it enters only with the latter’s annulling. But it is clear that this
quality enters only within immediate manifestness, within intuition, and
thus is only the representative and correlate of pure light. This latter is its
genetic principle by which, first of all, according to our hypothesis all
manifestness opens up into genetic manifestness, since pure light mani-
fests itself implicitly as genesis. Secondly, the previously presented rela-
tionship of concept to being and vice versa is further determined as fol-
lows: If there is to be an expression and realization of the absolute light,
then the concept must be posited, so that it can be negated by the imme-
diate light, since the expression of pure light consists just in this negation.
But the result of this expression is being in itself, period. [This result] is
inconceivable precisely because pure light is simultaneously destruction of
the concept. Thus, pure light has prevailed as the one focus and the sole
principle of both being and the concept.
4. From the preceding it follows that this inconceivable, as the bearer of all
reality in knowing, which we grasp 9 in its principle, is absolute only as
inconceivable, and cannot be thought in any other way. No other [60–61]
additional hidden qualities are attributable to it. Just as little can any qual-
ity be added to light, beyond the previously mentioned characteristic,
namely that it annuls the concept and remains absolute being. If we made
such additions, we would, as Kant has been criticized for doing, run up
against something unexplained and perhaps inexplicable. As support for
this contention, notice that we have understood it as inconceivable purely
in its form and nothing more. We have no right to assert anything before
we have seen into it. 10 So if we posit some other hidden quality, we have
either invented it, or better, since pure invention from nothing is com-
pletely impossible, we have manufactured it by trying to supply a principle44
The Science of Knowing
for some facticity. This happened with Kant when he first factically dis-
covered the distinction between the sensible and supersensible worlds and
then added to his absolute the additional inexplicable quality of linking the
two worlds, a move which pushed us back from genetic manifestness into
merely factical manifestness, completely contravening the inner spirit of
the science of knowing. Therefore, it is important to note that whatever yet
to be determined characteristics the reality appearing in our knowing may
carry in itself, besides the common basic property of inconceivability, such
characteristics by no means require any new absolute grounding principle
besides the one principle of pure light, since this would multiply the num-
ber of absolutes. Rather, the multiplicity and change of these various traits
is to be deduced purely from the interaction of the light with itself, in its
multifarious relations to concepts, and to inconceivable being.
I invite you to the following reflections so that I can offer a hint about
this last point, raise what has been said today to a higher level, give the new
listeners a unified perspective for viewing everything that they can learn here,
and give the returning listeners the same perspective from which they can
again gather and reproduce everything they have heard up to now.
The focus of everything is pure light. To truly come to this requires that
the concept be posited and annulled and that an intrinsically inconceivable
being be posited. If it is granted that the light should exist, then in this judg-
ment everything else mentioned is possible as well. We have now seen {einge-
sehen} this; it is true; it remains true forever; and it expresses the basic princi-
ple of all knowing. We can so designate it for ourselves.
[62–63] Now, however, I want to completely ignore the content of
this insight and reflect on its form, on our actual situation of insight. I also
think that we, those of us present here who have actually seen it, are the
ones who had the insight. As I remember it and as I think we all do, the
process was that we freely constructed the concepts and premisses with
which we began, that we held them up to each other freely, and that in
holding them together we were gripped by the conviction that they
belonged together absolutely and formed an indivisible oneness. Thus we
created at least the conditions for the self-manifesting insight, and so we
likewise appeared to ourselves unconditionally.
But let us not go to work with too much haste; rather let us consider
things a little more deeply. Did we create what we created because we wished
to do so, and therefore as the result of some earlier knowledge, which we
would have created because we wished to create it as the result of an even ear-
lier knowledge, and so on to infinity, so that we might never arrive at a first
creation? Somewhere, if the concept is created it must absolutely and thor-
oughly create itself, without anything antecedent and without any necessity ofFourth Lecture
45
a “we”; because this “we,” as has been shown, always and everywhere requires
some previous knowing and cannot achieve immediate knowing. Therefore
we cannot create the conditions, they must emerge spontaneously. Reason
must create itself, independent from any volition or freedom, or self. But this
proposition, disclosed through analysis, [contradicts] 11 the first which is given
reflexively, and so immediately. Which one is true, and on which should we
rely? Before trying to answer, let us return again to the matter of the insight
which has become controversial in regard to its principle, in order to grasp its
meaning and true worth clearly. We realized that if light is to be, 12 then the
concept must be posited and negated. Therefore, the light itself is not imme-
diately present in this insight, and the insight does not dissolve into light and
coincide with it; instead it is only an insight in relation to the light, an insight
which objectifies it, grasping it by its inner quality only. Thus, whatever the
principle and true bearer of this insight might be, whether we ourselves, as it
seems to be, or pure self-creating reason, as it also seems to be, the light is not
immediately present in this bearer, instead it is present merely mediately in a
representative and likeness of itself. First of all, [64–65] that this light occurs
merely mediately applies not just in the science of knowing but in any possi-
ble consciousness that has to posit a concept so it can annul it; and the science
of knowing rests on a completely different point than the one on which many
may have assumed it to rest after the last lecture, because undoubtedly know-
ing was understood too simply.
And now to answer the question: both are clear, therefore both are
equally true; and so, as was said at the start on another occasion, manifestness
rests neither in one nor in the other, but entirely between both. We arrive here,
and this is the first important and significant result, at the principle of division,
not as before a division between two terms, which in that case are to be intrin-
sically distinct like A and the point, but instead a division of something which
remains always inwardly self-same through all division. In a word, we have to
do with constructing and creating the very same primal concepts which
appear one time as immanent, in the unconditionally evident final being, the
I; and appear the other time as emanent, 13 in reason, absolute and in itself,
which nevertheless is completely objectified. Thus, it is a division pure and in
itself, without any result or alteration in the object.
Further, manifestness oscillates between these two perspectives: if it is
to be really constructed, then it must be constructed in that way. Thus, it must
be constructed as oscillating from a to b and again from b to a and as com-
pletely creating both; thus as oscillating between the twofold oscillation,
which was the first point, and which gives rise to a three or fivefold synthesis.
What is the common element in all these determinations? The very
same representative of the light, seen in its familiar inner quality. Here it stays.
Therefore everything is the same one common consciousness of light. This46
The Science of Knowing
consciousness, which is held in common and therefore cannot be really con-
structed but instead can only be thought 14 by means of the science of know-
ing, can be regarded or represented in the three or fivefold modifications only
from another standpoint, which this science alone oversees.
So much for now.C H A P T E R
5
Fifth Lecture
Monday, April 23, 1804
Honored Guests:
It might be appropriate to cancel the lecture this Wednesday, because of the
general day of prayer. 1 I would have taken up preliminary matters again
today had I not also seen the necessity, because of this, for sparing you the
strongest nourishment.
Indeed, I have already adduced and shared with you everything which
is conducive for understanding these lectures and which helps one enter their
standpoint except for two things: first, what really cannot be imparted, namely
the knack for grasping them; and second, some observations which tend not
to be received well and which I had hoped to be able to omit this time.
As far as concerns the first item, the knack for grasping these lectures is
the knack of full, complete attention. This should be acquired and exercised
before one enters on the study of the science of knowing. For this reason, in
the written prospectus for these lectures—available at the place of subscrip-
tion—I have established as the sole, but serious precondition for understand-
ing this science the requirement that participants should have experience with
fundamental scientific inquiry. Not, of course, for the sake of the specific
information so gained, since none of that is presupposed or even accepted here
without qualification; rather, I did so because this study also awakens and
exercises full, complete attention. Collaterally, one gains a knowledge of sci-
entific terminology, which we are using freely here. Full, complete attention, I
have said, which throws itself into the present object with all its spiritual
power, puts itself there and is completely absorbed in it, so that no other
thought or fancy can occur; since there is no room for anything strange in a
spirit totally absorbed in its object: full, complete attention as distinct from
that partial attention which hears with half an ear and thinks with half its
thinking power, interrupted and criss-crossed by all kinds of fugitive thoughts
This lecture begins at GA II, 8, pp. 66–67. It could have possibly been given on Tues-
day, April 24, 1804.
4748
The Science of Knowing
and fantasies, which eventually succeed in totally overwhelming the mind so
that the person gradually falls into a dreamy fog with eyes wide open. And if
he should chance to come back to awareness, he will wonder where he is and
what he has heard. This full, complete attention of which I speak, and which
only those who possess can recognize, has no degrees. It is distinguished from
that scattered attention, which is capable of many degrees [68–69], not merely
quantitatively; but is totally different and even logically contradictory to it. It
fills the spirit completely, while incomplete attention does not.
For understanding these lectures everything depends on one’s possess-
ing this kind of attention; everything which makes understanding difficult
or impossible follows solely from its absence; if one is freed from this lack,
then all these things are ripped out by the roots. So, for example, if this lack
is removed, the phenomenon of believing one cannot intuit the particular
theorems presented in the lecture because one is too quickly confused will
fall away. As I like to repeat frequently so no one will lose heart, in the
nature of our science the same thing is constantly repeated in the most vari-
ous terms and for the most diverse purposes, so that an insight which is
missed on one occasion can be produced or made good on another occasion.
But strictly speaking it should be, and is actually, demanded of everyone that
they see into each theorem when it is initially presented. So those for whom
things don’t happen as we expect have not used these lectures as they should
be used; and if things don’t flow smoothly, they have only themselves to
blame. To give the most decisive proof that what I demand is possible under
conditions of complete attentiveness—any distinction between faster or
slower mental capacity has no place in the science of knowing, and the pre-
sentation of this science aims neither at good nor slow minds but at minds
as such, if only they can pay attention. For this is our procedure as it has gone
up to now and as it will remain: first we are required to construct a specific
concept internally. This is not difficult: anyone just paying attention to the
description can do it; and we construct it in front of him. Next, hold together
what has been constructed; and then, without any assistance from us, an
insight will spring up by itself, like a lightning flash. The slowness or speed
of one’s mind has nothing more to do in this final event, because the mind
in general has no role in it. For we do not create the truth, and things would
be badly arranged if we had to do so; rather, truth creates itself by its own
power, and it does so wherever the conditions of its creation are present, in
the same way and at the same rate. And in case the ensuing manifestness did
not arise for someone who had really carried out the construction which we
postulated, this would only mean that he did not sustain the construction in
all clarity and power, but instead that it faded because some distraction
intervened. That is [70–71], he did not place his total attention on the pre-
sent operation.Fifth Lecture
49
Or, if this lack is removed, another equally common phenomenon
would be destroyed at its root: namely, that an illusion which we have already
revealed as an illusion, nevertheless can return and deceive us again, as if it had
truth and meaning, or at least confuse us and make us uncertain about insights
that we had otherwise already achieved. For example, if you have really seen
that in intuiting the one, eternally self-same knowing, all differentiation into
subjective and objective completely disappears as arising only in what is
changeable, then how can you ever again allow yourself to be deceived by the
illusion (which, to be sure, as an illusion, can always recur) that you yourself
are the very thing that objectively posits {objectiviert} this one knowing (that
you therefore are the subject with it as your object)? Because you indeed have
seen once and for all that this disjunction is always and everywhere the same
illusion and never the truth, no matter in what form or in what place it might
be manifest. If you have seen this, then you have attained this insight and dis-
solved into it. How then could you possibly cease being what you are, unless,
because you have not really entirely become it but only entered half way, you
never threw yourself into, and rooted yourself in, this insight which now
remains wavering and deceptive for you. In this case the old illusion returns at
the first opportunity. But note well the sequence: the insight does not leave
you because the illusion steps in, rather the illusion enters because the insight
has left you! So much as regards the talent of total attentiveness as the sure
and unerring means of correctly grasping the science of knowing. Second, I
want to mention a few more things that block apprehension of the science of
knowing, because they do not allow proper attention to arise. I take these
things up collectively in their oneness, as is my custom (as will likely be the
custom of anyone who becomes familiar with the science). They arise together
from a lack in one’s love of science, which is either a simple lack: a weak, pow-
erless, and distracted love; or a secret hatred of knowledge because of some
other love already present in the mind.
Let us first take up the last: the other love which leads to a secret hatred
of knowledge is the same one from which hatred arises against every good,
namely, a perverted self-love for the empirically arisen self instead of for the
self which is immersed in the good, the true and the beautiful. This love is
either that of self-valuation, which therefore becomes pride, or that of self-
enjoyment, which therefore becomes spiritual lasciviousness {Wohllustigkeit}.
[72–73] The first of these is unwilling to admit that anything could
occur in the domain of knowing that it had not itself discovered, and long
been aware of. Whether it explicitly says so or not, to such a one, the science
of knowing’s claim to absolute novelty seems to be a statement of contempt
for itself. It would very much like to humiliate this arrogance on the part of
this science—for this is how things must seem to it. Therefore, instead of giv-
ing itself freely and with complete attention from the start, it focuses on50
The Science of Knowing
whether it can possibly catch this science in some failing; and ambivalence of
purpose distracts it so that it misses the right idea, does not enter into the
true subject matter, but rather finds just what it was looking for in the con-
fused concepts which it obtains of the enterprise, weaknesses in its “science
of knowing.” 2
The other mode of thought, love for the empirical self ’s self-enjoyment,
loves the free play of its mental capacities (which it partially becomes) with
the objects of knowing (which it in the same manner partially becomes). I
think it can best be characterized in the following way: it calls making some-
thing up “thinking” and it names the invention of a truth for oneself in one’s
own body “thinking for oneself.” A science which brings all thinking without
exception under the most stringent rules and annuls all freedom of spirit in
the one, eternal, self-sustaining truth can hardly please such a disposition, and
it must also incite this later mode of thought to the same secret polemic
against itself, producing the very consequences we have just described. More-
over—just to take this opportunity to make the point decisively—I do not
warn each of you against this secret inner polemic for my own sake, but rather
for yours, because one cannot achieve correct attention, let alone understand-
ing, while doing it. If one will only first understand and master the science of
knowing and then feel a desire to argue against it, I will have nothing more to
say against this.
Or again it could be a ruling passion for the merely empirical and for
the absolute impossibility of feeling and enjoying one’s spirit in any way except
as trained memory. These personified memories are not capable of such secret
hate; but they necessarily become very ill-tempered in this setting. They want
what they call results: namely what can be observed and can be reproduced in
similar circumstances; “[that is] a sufficient statement and one that says some-
thing.” Now when they think they have grasped something of this kind, the
next lecture arrives further qualified, differently arranged, [74–75] 3 symbols
and expressions change, so that not much remains from the hard-won trea-
sure. “What eccentricity! Why couldn’t the man just say what he meant from
the start?” For people like that the most extreme confusion and contradiction
must arise from what has the purest oneness and strongest coherence, simply
because it is the true inner coherence and not the merely external schematic
coherence, which is all they really want. 4
Originally, I first mentioned cold, weak love of science (which is not
exactly hate) as an obstacle to attention. Namely, whoever seeks, desires, or
wishes in science for something besides science itself does not love it as it
ought to be loved and will never enjoy its complete love and favor in return.
Even the most beautiful of all purposes (that of moral improvement) is too
lowly in this case; what should I say of other, obviously inferior ones! Love of
the absolute (or God) is the rational spirit’s true element, in which alone itFifth Lecture
51
finds peace and blessedness; but science is the absolute’s sweet expression; and,
like the absolute, this can be loved only for its own sake. It is self-evident that
there is no room for anything common or ignoble in a soul given over to this
love, and that its purification and healing are intrinsic to it.
This love, like every absolute, recognizes only the one who has it. To
those who are not yet possessed by it, it can give only the negative advice to
remove all false loves and subordinate purposes, and to allow nothing of that
kind to arise in them so that the right will spontaneously manifest itself with-
out any assistance from them. This much should be remembered once and for
all on this subject.
Now to the topic set for today. When I presented it, I already suspected
that my last talk might seem too rigorous and deep for a fourth lecture
[76–77], and it was made so in part to help me discover what mode of pre-
sentation I would need to follow with this new audience. Now I will repeat it
in a suitable form:
1. First, a remark that is valid for all previous and subsequent lectures, and
that will be very useful in order to reproduce and review them. Our proce-
dure is almost always this:
a. we perform something, undoubtedly led in this process by a rule of rea-
son which operates immediately in us. What in this case we really are in
our highest peak, and that in which we culminate, is still only facticity.
b. we then search out and reveal the law which guided us mechanically in
the initial action. Hence, we see mediately into what we previously had
seen into immediately, on the basis of its principle and the ground of its
being as it is; and we penetrate it in the origin {Genesis} of its determi-
nateness. In this way we will ascend from factical terms to genetic ones.
These genetic elements can themselves become factical in another per-
spective, in which case we would be compelled again in connection with
this new facticity to ascend genetically, until we arrive at the absolute
source, the source of the science of knowing. This is now noted and can
be clarified in reference to the consequences: x is nothing but the devel-
opmental link to y, and y in turn to z. 5
Now, whoever either has not comprehended z from the start, or has lost
and forgotten this understanding in the process, for that person neither x nor
y exists and the entire lecture has become an oration about nothing, through
no fault of the lecturer. This, I say, has been and will continue to be our pro-
cedure for some time. It was so in the last lecture. Whoever may have recog-
nized this process—it was obvious for everyone to see, and the earlier distinc-
tion between factical and genetic manifestness should have led right to
it—could have reproduced the entire lecture and made it intrinsically clear by52
The Science of Knowing
simply asking: Was any such factical term presented, and which one was it?
Which could it be after the earlier ones? Did the presentation succeed in pre-
senting the genetic term after the factical one? Assume that I may have com-
pletely forgotten this second step or perhaps never heard of it. Then I will
have to discover it for myself just as it was discovered in the talk, because the
rule of reason is unitary, and all reason which simply collects itself is self-same.
So what was the factical term? It was not in A and not in the point, but
unconditionally in both. We have now grasped {eingesehen} this, it has made
itself evident, and so it is. Analyze [78–79] it however you wish: it contains
A, the point, and, in the background, a union of both. With the first two
terms denied as the true point of oneness, the other one is thereby posited,
and in this fashion you will not arrive at any other term. It is so, factically.
But now I ask on another level: How have we brought it about that this
insight has arisen for us? We did not reflect further on the content, which we
completely abandoned; but focused instead on the procedure, asking about the
origin {Genesis}. In this way, as I indicated earlier, the initial, materially con-
stituted, immediacy 6 becomes mediately visible: once such an origination is
posited, this factical insight is posited, but solely by means of our establish-
ment of the origin.
How did we do it? Apparently we made a division in something which
on the other hand ought to be a oneness. 7 I say division and disjunction in a
general way, because one can ignore the fact that the terms separated are “A”
and the point, when it is a question of the act qua act. 8 This division shows
itself to be invalid in an immediate insight {Einleuchten}. We did not produce
this insight because we wished to, instead it produced itself absolutely (not
from any ground or premise) in an absolutely self-generating and self-pre-
senting manifestness, or pure light. The distinction, in the sense that it should
be valid by itself, would therefore be annulled by the [one’s] manifestness. On
the contrary, the same manifestness posits a self-same, intrinsically valid one-
ness which is incapable of any inner disjunction. The principle of division
equals the principle of construction, and thus of the concept as well. [Now,
consider] this principle in its absoluteness 9 (and by that I mean, the principle
as dividing the wholly and intrinsically one, which is seen into as one, [work-
ing] wholly and absolutely by itself without other ground and doing so rather
in contradiction to the truth)—this principle is negated in its absoluteness,
i.e., in its intrinsic validity. 10 It is seen unconditionally as negated, and there-
fore it is negated in and through the absolute light. Thus, in this annulling of
the absolute concept in relation to intrinsic being {das Wesen an sich}, this being
is inconceivable. Without this relation it is not even inconceivable but rather is
only absolute self-sufficiency. But further even this predicate “is” derives from
manifestness. Hence the sole remaining ground and midpoint is the pure light,
and so on.Fifth Lecture
53
This was by far the greatest part of the earlier talk’s content. That this
all has intrinsic clarity and incomparable manifestness is obvious to anyone
who sees it at all. [80–81] I am convinced that it could not be presented with
greater order, distinctness, clarity, and precision than in this case. Whomever
doesn’t see it now must be lacking in the undivided attention required here.
The part still to be added, which I will repeat now, is another develop-
mental analysis {Genetisieren} of the insight we have achieved: I said that we
saw into [the fact] that the light was the sole midpoint. In this reflexion, the
process we initially unfolded itself becomes factical. Now, since in this case we
have produced nothing, and [since] rather the insight as insight has produced
itself, we cannot ask as we did before how we did this, but we can rise to
greater clarity. It is clear; if only we see {sehen wir nur ein} that it is the light,
then we are not immediately consumed in this light, instead we have the light
present through its agent or representative, that is, through an insight into the
light, into its originality or absolute quality. We must disregard the fact that
we cannot now ask without contradiction how the light itself is produced;
since it is recognized as the principle of absolute creation, the question would
deny the insight again. We can certainly still ask how the insight into light
(which we called not light itself but rather its agent and representative) has
been produced; that is a different question. Therefore we only need to pay
attention to how the production of this insight has taken place. 1. We have put
ourselves in the condition; 2. how could we do this?—Both are true:—not the
light and not even the insight into light, but the insight into the insight into
the light stands between both. [The emanence and the immanence: these are
matters with which we must concern ourselves.
Regarding the entire distinction between the immanence and emanence
of the production of an insight into light, one must not forget that the same
thing extends to the insight and to the light itself. As before, the objective light
qua objective neither is, nor can become, the one true light. Instead pure light
enters insight under this aspect. But here we have won this: that the highest
object is no longer substance for itself, but light. Substance is only the form of
light as self-sufficient. On the other hand, insight (subjectivity), actually the
inner expression and life of light, disengages itself from the negation of the
concept, and of division. Can you penetrate into the true midpoint more
deeply in any other way? Into the entirely unique concept that is nevertheless
required here? 1. It is clear “that its being is not grasped except in immediate
doing.” 2. It can be made clear [82–83] here that immediate doing is a dissolv-
ing into immanence (the initially uncovered making of his being, as this sort
of making). First of all, doing deposes being and being deposes doing—or stated
otherwise: here is the fundamental reversal; and this must be understood : doing
replaces actual being—being beyond all being (not actual or material being)
deposes doing. Now it is also very clear that—to posit it without any actuality54
The Science of Knowing
(just in barely hypothetical form and in a “should”) and thereby to deduce and
materialize doing itself, thus also to intellectualize and idealize—being
negates itself in the other as a result of its own doing. Thus we once again
come back to the previous point, and we find the previous principle again in
this self-negation. Perhaps this is just the concept of being, dead in-itself: clearly
there is a division in it between being (what endures) and doing: and indeed
as a division this is intrinsic to constructing non-separation or oneness. Thus
this negation would be true in a certain respect. It is right since primordially
the division into being and doing is nothing at all.] 11C H A P T E R
6
Sixth Lecture
Thursday, April 26, 1804
Honorable Guests:
In today’s and tomorrow’s talks I will continue with the further development
of what has already been presented. In doing so I aim at an end useful to both
sorts of listeners. That is, since, like all philosophy, the science of knowing has
the task of tracing all multiplicity back to absolute oneness (and, correlatively,
to deduce all multiplicity from 1 oneness), it is clear that it itself stands neither
in oneness nor in multiplicity, but rather stays persistently between both. It
never descends into absolute multiplicity, which must after all exist and indeed
does exist (as mere empirical givenness), but rather it maintains the perspec-
tive from above, from the standpoint of its origin. Therefore, in the science of
knowing we will be very busy with multiplicity and disjunctions.
Now, these disjunctions, or differences and distinctions, which the sci-
ence of knowing has to make are new and previously unrecognized. There-
fore, in the usual modes of representation and speech from which we begin,
these differences collapse unnoticed into oneness, and when we are required
to draw them, they seem very minute. (It is hairsplitting, 2 as the literary rab-
ble has put it; and it is necessary that it be so, since if a science that is to trace
everything that is multiple—that is, everything in which a distinction can be
constructed—back to oneness allows any distinction that the science could
possibly make to remain hidden, then it has failed in its purpose.) Therefore
one of main problems for the science of knowing consists in just this: mak-
ing its very precise distinctions visible and distinct; so that when this prob-
lem is finally solved, these distinctions will be fixed and established in the
mind of those studying it, so that they will never again confuse them. I think
that both difficulties will be significantly reduced if I lay out for you in
advance (so far as this is possible) the general schema and basic rule in terms
of which these divisions will come about—although in an empty and purely
formal way. And, so that this schema can be correctly understood and
This lecture begins at GA II, 8, pp. 84–85.
5556
The Science of Knowing
noticed, I will deduce it in its unity and from its principles, to the extent that
this can be done with what we know so far.
To begin I mention in general the following:
[86–87] 1. Since, according to the nature of our science, we must stand
neither in oneness nor in multiplicity but instead between the two, it is clear—
and I focus on this because I believe I have detected several of you making this
error—that no oneness at all that appears to us as a simple oneness, or that
will appear to us as such in what follows, can be the true oneness. Rather, the
true and proper oneness can only be the principle simultaneously of both the
apparent oneness and the apparent multiplicity. And it cannot be this as some-
thing external, such that it merely projects oneness and the principle of mul-
tiplicity, throwing off an objective appearance; rather it must be so inwardly
and organically, so that it cannot be a principle of oneness without at the same
moment being a principle of disjunction, and vice versa; and it must be com-
prehended as such. Oneness consists in just this absolute, inwardly living,
active and powerful, and utterly irrepressible essence.—To put it simply, one-
ness cannot in any way consist in what we see or conceive as the science of
knowing, because that would be something objective; rather it consists in what
we are, and pursue, and live.—Let this be introduced once and for all to char-
acterize the oneness which we seek and to eliminate all the errors about this
central point which, if they continued, would necessarily be very confusing in
what follows. And be warned, not only so that you don’t content yourself
merely with that sort of oneness, taken just relatively and one-sidedly, as if it
was the absolute, but also so that if I in this lecture, or any other philosopher,
remain content with such a oneness, you will know and state strongly that this
philosopher has stopped half way and has not made things clear.
2. In consequence: Since the true oneness is the principle simultaneously
of the (apparent) oneness and of disjunction, and not one without the other,
it therefore makes no difference whether we regard what we will present pro-
visionally as our highest principle at each juncture in the progress of our lec-
tures as a principle of oneness or of disjunction. Both perspectives are one-sided,
merely our necessary point of view, but not true in themselves. Implicitly the
principle is neither one nor the other; rather it is both as an organic oneness and
is itself their organic oneness.
Therefore, so that I can say it even more clearly—first of all only princi-
ples can enter the circle of our science. Whatever is not in any possible
[88–89] respect a principle, but is instead only a principled result {Principiat}
and phenomenal, falls to the empirical level, which, of course, we understand
on the basis of its principle, but which we never scientifically construct, as this
cannot ever be done. Then, every principle that enters our science (and indeed
every principle qua principle) is simultaneously a principle of oneness and
multiplicity, and it is truly understood only insofar as it is conceived in thoseSixth Lecture
57
terms. Our own scientific life and activity, therefore, to the extent that it is a
process of penetrating and merging with the principle, never enters into that
oneness which is opposed to multiplicity nor into multiplicity; rather it main-
tains itself undisturbed between both, just like the principle. Finally each prin-
ciple in which we stand (and we never stand anywhere but in a principle)
yields an absolutely self-differentiating oneness: 3
x
a = a—y
z
The only question is whether this oneness is the highest. If not, and there
are several such a’s (a 1 ,a 2 ,a 3 , . . .), then not just in the former case but in the lat-
ter as well, “a” is still in this regard a principle of disjunction for unities, which
to be sure would be unities in relation to
x
y
z
but in connection to one another, they would by no means be so. For these a’s
we need a new a, until we have uncovered the highest oneness, which would
be the absolute disjunction, just as we have described it in relation to the
absolute oneness. This gives us the first general model for the procedure of the
science of knowing.
One comment here: the interchangeability of direction from a to x, y, z
and vice versa is evident, and this greatly aids their linkage.
3. Now the same point, from another side and deeper. As regards the
explanation we have been pursuing up to now in this hour (not about the prin-
ciple of disjunction, since strictly speaking there is no such thing, but rather our
view of the one implicit principle as a principle of disjunction, a one-sided
view that we undoubtedly must start with since the science of knowing finds
us completely trapped in this one-sidedness and starts from there), we find
ourselves trapped in the familiar, frequently cited inexpressibility: that the
oneness is to separate itself at one stroke into being and thinking and [90–91]
into x, y, z, both equally immediately. In expressing this verbally and in dia-
grams, we were compelled to make one of the two the immediate term,
though our inner insight contradicted this, negating the intrinsic adequacy of
the construction of our mode of expression. Expressing this curious relation in
its logical form will help us to speak precisely: in this actual disjunction there
are two distinguishing grounds {fundamenta divisionis} neither of which can
occur without the other. Therefore, expressing the matter just in the way we58
The Science of Knowing
have done is probably an empirically discovered turn of phrase, since we
found it on the occasion of explaining Kant’s philosophy and in adding a dis-
junction, which has been demonstrated neither by Kant 4 nor yet by us, 5 not
only between being and thinking, but between sensible and supersensible
being and thinking. And the claim that both distinguishing grounds are
absolutely inseparable would therefore be grounded simply on this: If what is
evident in empirical self-observation is to be explained, then we must assume
that the distinguishing grounds are inseparable. This “must” grounds itself
directly on a law of reason which operates in us mechanically and without our
awareness [of it] {eigene Einsicht}. Thus at bottom we had only an empirical
basis on top of which we postulated a supersensible one; 6 that is, we began a
synthesis post factum. This cannot be blamed on the science of knowing as
long as it is the science of knowing; it is not permitted simply to report this
inseparability of the distinguishing grounds, instead it must grasp this ground
conceptually in its principle and from its principle as necessary. It must there-
fore see into it genetically and mediately. “It grasps this ground conceptually”
means it sees the distinguishing grounds (and by no means just the actual
factically evident distinctions—whoever remains with these has simply not
finished the climb we have just completed) as themselves disjunctive terms of
a higher oneness, in which they are one and inseparable as they are when
enacted {im Akte}, so that, as we have said, it remains one and the same
stroke. But they are separable and conceptually distinguishable, as we may
provisionally think in order thereby to have something to think. “Separable”
so that, for example, the ground for distinguishing being and thinking can
appear as a further determination and modification of the ground for distin-
guishing sensible and supersensible, and so that likewise from another point
of view [92–93] the latter can appear vice versa as a further determination
and modification of the former. As has been said, when it is enacted, this dis-
junction in the oneness of the mere concept concresces {concresciert} into a fac-
tical oneness which is not further distinguishable, and in this concrete union,
every eye that remains factical is entirely closed to the higher world of the
conceptual beyond.
(And now a number of additional remarks. I ask that you not allow
yourself to be distracted while I state them):
1. I have now specified the boundary point between absolutely all factical
insight and truly philosophical and genetic insight entirely and exactly, and
I have opened up the sources of the entirely new world in concepts which
appear only in the science of knowing. The creation and essence of this new
world is found just here, in the negation of the primordial disjunctive act
as immediate and in the insight into this primordial act’s principle—mate-
rially, that it is thus, and formally, that it is at all.Sixth Lecture
59
2. I have here explained the essence of the final scientific form of the science
of knowing from a single point more precisely than I have been able to do
previously. The main point of this scientific form resides in seeing into the
oneness of the distinguishing ground—of being and thinking as one and of
sensible and supersensible (as I will say in the meantime) as one. Whoever
has understood this—as it is to be understood up to now, namely as an
empty form—and holds on to it firmly, can scarcely make any further errors
in the subsequent actual employment of this form.
3. In order to assist both your memories and repeated reproduction: in the
last hour I said that the path of our lectures was, and would for a long time
continue to be, that we first present something in factical manifestness
and then would ascend to a genetic insight into this object on the basis of
its principles. This is exactly what we have done in the just-completed
explanation. Already since the second lecture, we had developed the
inseparability of both recognized distinguishing grounds historically out
of Kant’s own statement, and we admitted the factical correctness of this
statement. Now we raise ourselves—to be sure not to a genetic insight
into the principle of this inseparability (since we do not yet actually know
this inseparability itself, nor its terms, but have only assumed all this tem-
porarily and for the time being—but to the genetic insight, which must be
the form of this principle, if such inseparability and such a genetic princi-
ple are to exist).
Now back to our project. It is also not at all our intention to see directly
into this inseparability and its principle, since these do not allow themselves
to be “seen into” directly. And in fact:
4. to take the process further—by our beginning we have already jumped past
this principle, which was discussed here in its form simply for the intelli-
gibility of what really concerns us, in order to derive it deductively; and
[94–95] indeed we have already uncovered good preliminaries for this
derivation. Namely, you recall that we have already presented a point of
oneness and difference, which covers the oneness of these distinguishing
grounds: the one between A and the point, and, in connection with the
deeper distinguishing grounds which are materially different from oneness
and difference, we have said that this might be only a profounder view of
this same higher principle, disregarding the fact that we could not yet prove
this contention. 7
[Let me] repeat a third time this oneness which has already been
constructed twice before our eyes. To that end I recall only that an absolute
dividing principle was evident there—not A and the mentioned disjunctive
point, since these are principled results {Principiate} of absolute division,60
The Science of Knowing
which disappear when one looks to the principle; but rather the living
absolute separation within us. I stress again what I said before about this
essential point, which is designed to tear our eyes away from facticity and
to lead them into the world of the pure concept, if I can succeed in mak-
ing it clearer. I hope nobody assumes here that the act of thinking the dis-
tinction between A and . is actually grounded in an original distinction in
these things, independent of our thinking. Or, in case someone is led to this
conclusion by the previous factical ascent with which we had to begin, he
will recover from this idea if he considers that in A and . he thinks only the
oneness which, according to him, should be unconditionally one contain-
ing no distinction within itself; that he himself thus makes clear that the
distinction is not based in the object itself, since he could not think the
object except by virtue of this distinction; that he thus expressly makes his
own thinking as thinking into the distinguishing principle. But the valid-
ity and result of this product of thought expressly surrenders and dies in
relation to the thing itself. With it as the root, its products A and . are also
doubtlessly uprooted and destroyed as intrinsically valid. Thus away with
all words and signs! Nothing remains except our living thinking and
insight, which can’t be shown on a blackboard nor be represented in any
way but can only be surrendered to nature.
[We intuit, I say, that it rests neither in A nor in . but rather in the
absolute oneness of both; we intuit it unconditionally without sources or
premisses. Absolute insight therefore presents itself here. Pure intuition,
pure light, from nothing out of nothing, going nowhere. To be sure bring-
ing oneness with it, but in no way based on it.] 8
[96–97] Here everything depends on this: that each person correctly
identify with this insight, in this pure light; if each one does, then nothing
will happen to extinguish this light again and to separate it from yourself.
Each will see that the light exists only insofar as it intuits vitally in him,
even intuits what has been established. The light exists only in living self-
presentation as absolute insight, and whomever it does not thus grasp,
hold, and fix in the place where we now stand, that one never arrives at the
living light, no matter what apparent substitute for it he may have.
5. Consideration of the light in its inner quality, and what follows from it, to
which we will proceed after this step, is entirely different from this surren-
der and disappearance into the living light. This consideration as such will
inwardly objectify and kill the light, as we will soon see more precisely. But
first we said: only the light remains as eternal and absolute; and this [light],
through its own inner immediate essence, sets down what is self-subsistent,
and this latter loses its previously admitted immediacy to the light, whose
product it is. But there is no life or expression of this light except through
negating the concept, and hence through positing it. As we said, no expres-Sixth Lecture
61
sion or life of the light could arise unless we first unconditionally posit and
see a life as a necessary determination of the light’s being, without which
no being is ever reached, except, that is, in the light itself—its essence in
itself and its being, which can only be a living being. Thereby, however,
what matters to us, since we add life to light, is that we have nevertheless
divided the two, have therefore, as I said, actually killed the light’s inner
liveliness by our act of distinction; that is, by the concept. Now, to be sure,
we contradict ourselves, ipso facto denying that life can be distinguished
from light, the very thing we have just accomplished. This is a contradic-
tion which may well be essential and necessary, since it may implicitly in
itself be the negation of the concept to which, according to the foregoing,
it must someday certainly come. (What I’m saying now is added paren-
thetically for future use. It is easy to remember; since it connects with our
reflections on the objectifying consideration of the light, and allows itself
to be reproduced from it for anyone who has paid even a little attention to
our proceedings, in case he has otherwise completely forgotten.)
To review: in this consideration of the light, light shows itself through
its mere positedness, absolutely and without anything further, as the ground
for a self-subsisting being—and at the same time for the concept ; and, to be
sure, for the concept [98–99] in a twofold sense: in part as negated, precisely
in its intrinsic validity; and in part as posited, posited as absolute but not real-
ized (though still actual); that is, as appearance and as the light’s vitality, but in
no way an appearance that conditions its inner essence. By the concept’s being
posited, A and . are also posited—to be sure as appearance and certainly not as
primordial appearance, but rather as conditioning appearance and the inner life
of the primordial appearance = b, thus appearance’s appearance. [In its inner life,
appearing should occur again as the unity of the above mentioned distin-
guishing grounds, its life comes from the livingness of the concept, this in turn
stems from the light’s livingness, thus an appearance of appearance of appear-
ance. Everything is brought together again when enacted {im Akte}. This
would be the schema of an established, rule-governed descent, in no way like
[the one outlined] yesterday, one equally possible on all sides and therefore
very exposed to error.] 9C H A P T E R
7
Seventh Lecture
Friday, April 27, 1804
Honorable Guests:
[Our] purpose [here is] to give a brief account of the rules according to which
the disjunction we will have to make proceeds.
1. [It is a disjunction into] principles, with each being equally a principle of
unity and of disjunction.
2. [It provides a deduction from this to a general] schema of the total empir-
ical domain according to the form of its genetic principle 1 —this will be an
entirely new explanation, because I observed to my pleasure that some of
you had seen that there is something else even more deeply hidden, despite
the fact that you could not assist yourselves [to find it], which, to be sure,
was not even required. 2
3
3. With the remark that our investigation has already gone beyond this prin-
ciple to a higher one and that it has already begun to deduce this principle
itself, [I repeat] this achieved insight. Neither in A nor .—for us the one-
ness beyond is nothing in itself, although it is posited as in itself ; rather it
exists only through the light and in the light, and (is) its projection—light
itself—contemplation of light. Now back to our former topic— 4
[There is] one further step which opens up a whole new side of our
investigation; as I said: we have already previously begun to derive the princi-
ple of oneness and disjunction of materially different principles of division,
only without recognizing that this was what we were doing. 5
So (dropping what we have done so far until I take it up again) recall
with me and consider the following: when we observe the light, the light is
objectified, alienated from us and killed as something primordial. We have
explained what is attributed materially to the light in this observation of the
light, and we have connected this explanation to the schema under consideration. 6
This lecture begins at GA II, 8, pp. 100–101.
62Seventh Lecture
63
Now we will explain this observing itself 7 in its inner form, that is, no longer
asking what it contains and leads to, but rather how it itself inwardly occurs,
while also rising to its principle and viewing it genetically to some extent. It is
immediately clear that: 1. the light is in us (that is, in what we ourselves are and
do in observing it) not immediately but rather through a representative or proxy, 8
which objectifies it as such, and so kills it. So then, where does the highest one-
ness and the true principle now rest? No longer, as above, in the light itself,
since we, as living, dissolve in the light. Neither [is it] in the representative and
image of the light which is to be identified now: because it is clear that a rep-
resentative without the representation of what is represented or an image with-
out the imaging of what it images, is nothing. In short, an image as such,
according to its nature, has no intrinsic self-sufficiency, but rather points
toward some external, primordial source. Here, therefore, we have not only, as
above, factical manifestness, as [102–103] with A and .; instead [we even have]
conceptual 9 manifestness: oneness only with disjunction, and vice versa. “Even
conceptual manifestness,” I say: something imaged—like the light, in this
case—is not thinkable without an image, nor likewise an image, qua image,
without something imaged. Notice this important fact, which will take one
deep into the subject matter, if it is properly grasped here. In this case you carry
out an act of thinking {ein Denken}, which has essence, spirit, and meaning and
is fully and completely self-identical and unchangeable in relation to this
essence. I cannot share this with you directly, nor can you share it with me; but
we can construct it, either from the concept of something imaged which then
posits an image, or from the image which then posits something imaged. I ask:
apart from the arrangement of the terms, which is irrelevant here, have we then
thought two different things in the two concepts thus fulfilled, or have we not
rather thought exactly the same thing in both, an issue that genuinely touches
the inner content of thinking? The listener must be able to elevate himself to
the required abstract level out of the irrelevance of the arrangement to the
essential matter of the content, of spirit, and of meaning, and then the insight
which is intended will immediately manifest itself to him. Should this indeed
be the case, then an absolute oneness 10 of content is manifest here that remains
unaltered as oneness but that splits itself only in the vital fulfillment 11 of think-
ing into an inessential disjunction, which neither spoils the content in any way
nor is grounded in it. Either [there is] an objective disjunction into something
imaged and its image, or, if you prefer, [there is] a subjective-objective disjunction
into a conception of something imaged on the basis of the directly posited
image, and a conception of the image on the basis of the directly posited imaged
something. I do advise you to prefer the latter, since in that case you have the
disjunction first hand.—And thus in this case our principle in the genetically-
oriented {genesirten} view of the light would be the concealed oneness, which
cannot be described further, but which is lived only immediately in this act of64
The Science of Knowing
seeing, and which, as the primordial concept’s content, presents itself as
absolute oneness, and as absolute disjunction, in its living fulfillment. Now, the
something imaged in the concept’s content should be the light: therefore our
principle (i.e., we, ourselves) rests no longer either in the light or in the light’s
representative, but rather in the oneness in and between the two, a oneness
realized in our act of thinking. 12 Therefore, I have called the concept situated
here the primordial concept {Urbegriff }. [I call it this] because what up to now we
have taken as the source of the absolutely self-sufficient, and which therefore
appeared as the original and was original for us, actually first arises in the way
it appears, in its objectivity, from this concept as one of its disjunctive terms.
Therefore, this concept is [104–105] more original than the light itself; hence,
so far as we have yet gone, it is in this sense what is truly original. Thus we have
given a deeper genetic explanation of the hint, given only as a fact in the lec-
ture just eight days ago, 13 concerning the representatives of the original light,
although, to be sure, we have done so for the particular end we intend here.
Thereby you will see that the concept is determined further and grasped
more deeply than it has been heretofore. Until now it was a dividing principle
which, as self-sufficient, expired in the light, and which preserved only a bare
factical existence as an appearance, qualifying the appearance of the original
light. Further, it had no contents and acquired no contents except that which
pure light added to it in immediate intuition through a higher synthetic unity.
Now, however, the concept has its own implicit content, which is self-subsis-
tent, totally unchangeable and undeniable; and the principle of division (which
arises again in this case, and as before is negated as intrinsically valid) is no
longer essential to it, but instead only conditions its life, i.e., its appearance. The
concept’s content, I say, is self-subsistent; thus it is exactly the same substantial
being which was previously projected out of intuition and which manifests here
in the concept as prior to 14 all intuition and as the principle of the objective and
objectifying intuition itself. Previously, the concept qualified both life and the
appearance of light, and these conversely qualified the concept’s being. There-
fore, it was a reciprocal influence, and every [act of ] thinking the two terms was
qualified externally. Now the same single concept grounds its appearance
through its own essential being; therefore, in this concept the image and what
it images are posited absolutely, things which are constructed organically only
through one another. And, hence, its appearance announces, and is the expo-
nent of, its inner being, as an organic unity of the through-one-another
{Durcheinander}, which must be presupposed. Its being for itself, permanent
and unchanging, and as an inner organization of the through-one-another—
essential, but in no way externally constructed—are completely one: therefore,
in this case, absolute oneness is grounded and explained through itself.
We will achieve a great deal if here and now we see fundamentally into
what I meant by the inner organic oneness of the primordial concept, whichSeventh Lecture
65
I mentioned just now; since this oneness is the very thing we will need as we
continue. In this regard I ask: Does the image, as image, completely and
unconditionally 15 posit something imaged? And if you answer “Yes,” does
not the something imaged likewise posit such an image? Now [106–107]
without further ado I admit that both can be seen (by you) as posited imme-
diately by the other, but only if you posit one of the two as prior. But I ask
you for once to abstract from your own insight. This is possible in the way
that I will preconstruct for you now, and in ordinary life it happens con-
stantly when it shouldn’t. Further, one could not ever enter the science of
knowing without this abstraction. 16 That is, I ask about the truth in itself,
which we recognize as being and remaining true even if no one saw it, and
we ask: Is it not true in itself that the image entails something imaged and
vice versa? And, in this case, what exactly is true in itself ? Just reduce what
remains as a pure truth to the briefest expression. Perhaps that a posits b and
b, a? Do we want to divide the true into two parts and then link these parts
by the empty expletive “and,” a word which we scarcely understand and
which is the least understandable word in all language, a word which is
unexplicated by any previous philosophy? (It is indeed the synthesis post fac-
tum.) How could we, since beyond this it is certainly clear that the determi-
nation of the terms derives solely from their place in the sequence—for
example, that image is the consequent because something imaged is the
antecedent, and vice versa.—Further, if one enters more deeply into the
meaning and sense of both terms, it is clear that their meaning simply
changes itself into the expression “antecedent” and “consequent,” while
something imaged is really antecedent and so forth: thus all this dissolves
into appearance. So then, what common element remains behind as the con-
dition for the whole exchange? Obviously only the through-one-another that
initially holds together every inference 17 however it might have been
grasped, and which, as through-one-another, leaves the consequence rela-
tion exactly as free in general as it has appeared to be. 18
[Let me say this in a preliminary way (there is not time today to explain
the deeper view which is possible here and which I will go into in the next
hour)—. The focal point = the concept of a pure enduring through-one-another
in living appearance: 19
A
|
E
/ \
I
R
B–A a–B
A–B B–a66
The Science of Knowing
a—act and consequence either ideally, or [108–109] really—I say
either/or, it always remains act, the concept proceeds from it alive, but not fin-
ished and complete; whomever wants this must do it.
On the other side, the concept projects the one eternally self-same light
as intuition, from which follows (and this is its absolute essence): what stands
beneath it are parts of its externalization and thus are further modifications
not of the light but of its appearance {Ansicht} in the living concept. 1. Only
through life to the concept and only through the concept to life and the
appearance of the light in itself: but its first modification, never as pure but in
one or another variant. 2. Creation of the science of knowing in its possible
modifications.]C H A P T E R
8
Eighth Lecture
Monday, April 30, 1804
Honorable Guests:
I believe we have arrived at a focal point in our lectures, a point which more
effectively facilitates a clear insight and overview than does any other, and
which therefore will permit us greater brevity in what follows. Therefore, let
us not economize on time now, and from here on out we will make ourselves
more secure. Today we will do this with the contents of the last lecture.
We have seen {eingesehen} actually and in fact, and not just provisionally,
that an absolute, self-grounded insight negates an equally absolutely created
division (that is, one not grounded in things) as invalid, and that this insight
posits in the background a self-subsistence, which cannot itself be described
more precisely.—Let this be today’s first observation: at this point the main
thing is that all of us assembled here together have really and actually seen
into this just as it has been presented, and that we will never again forget this
self-insight or allow it to fade; but rather we will take root in it and flow
together with it into one. 1
Thus—what has been said is not just my report or that of any other
philosopher, but rather it exists {ist} unconditionally, and it remains always
true, before anyone actually sees into it, and even if no one ever does. We, in
our own persons, have penetrated to the core and have viewed the truth with
our own eyes. Likewise, as has been evident from the first, what has been said
is in no way proposed as a hypothetical proposition, which is shown to be true
in itself only by way of its usefulness in explaining phenomena, as is the case
in the Kantian, and every other philosophy. Instead it is immediately true
independently of all phenomena and their explicability. (A good reason for
making this more precise!) Therefore, what genuinely follows from this, if
only it is itself completely enough determined, is also as unconditionally true
as it [i.e., the original insight] is. And everything which contradicts it, or the
least of its results, is unconditionally false and should be abandoned as false
This lecture begins at GA II, 8, pp. 100–111.
6768
The Science of Knowing
and deceptive. This categorical decisiveness between truth and error
[112–113] is the condition for our, and every, science; and it is presupposed.
It is far removed from that skeptical paralysis which in our days parades as
“wisdom,” doubts what is unconditionally manifest, and wants to make the
latter clearer and more manifest by the most derivative means.
And as particularly concerns the explanation of phenomena from a
principle that is manifest, it is obvious right away that if the principle is sound,
and likewise the inferences, then things will go well with the explanation; we
need only note that, since the principle first grants us a true insight into the
phenomenon’s essence as such, it may well happen that in this proof process
many things do not even have the honor of being genuine, orderly phenom-
ena, but rather dissolve into deceptions and phantasms, although all ages have
held them to be phenomena, or—God forbid—even to be self-subsisting real-
ities. Therefore, it may happen that in this regard science, far from acquiring
some law or orientation from the factical apprehension of appearances, on the
contrary rather legislates for them. This situation can also be expressed as fol-
lows. Only what can be derived from the principle counts as a phenomenon;
what cannot be derived from it is an error simply because of its non-deriva-
tion, although it may perhaps incidentally also be immediate, if one wants to
boast of this direct proof. 2
This is the second observation—already as a result of this just repeated
insight, a new world of light has opened for us, which transcends our entire
actual knowledge, and a world of error, in which nearly every mortal without
exception finds himself, has perished, especially if we appropriate what follows
and is recorded below. If we take up this result directly here, it will be invigo-
rating for the attention and throw a very beneficial light on what comes next.
[114–115] 1. By annihilating the formal concept, which is the condition
for its own real appearance and vivacity, the light, as the one true self-suffi-
ciency, posits a self-sufficient being, which is not further determinable and
which, as a result of the insufficiency {Nichtgültigkeit} of the concept that
attempts to grasp it, is inconceivable. The light is simply one, the concept that
disappears in the light is one (the division of what is one in-itself ), and being
is one; it can never be an issue of anything other than these three. 3 The one
existence arises 4 in the intuition of what is independent, and in the concept’s
negation (and it will turn out that whatever genuinely concerns true existence
will also rest in this). If, as is customary, you want to call the absolutely inde-
pendent One, the self consuming being, God, 5 then [you could say that] all
genuine existence is the intuition of God. But at the same time note well—
and already a world of errors will be extinguished—that this being, despite the
fact that the light posits it as absolutely independent (because the light loses
itself in its life), is actually not so, just because it bears within itself the pred-
icate “is,” “persistence,” {des Bestehens} and therefore death. Instead, what isEighth Lecture
69
truly absolutely independent is just the light; and thus divinity must be posited
in the living light and not in dead being. And not, I certainly hope, in us, as
the science of knowing has often been misunderstood to say; for however one
may try to understand that, it is senseless. This is the difficulty with every phi-
losophy that wants to avoid dualism 6 and is instead really serious about the
quest for oneness: either we must perish, or God must. We will not, and God
ought not! The first brave thinker who saw the light about this must have
understood full well that if the negation is to be carried out, we must undergo
it ourselves: Spinoza was that thinker. It is clear and undeniable in his system
that every separate existence vanishes as [something] independently valid and
self-subsistent. But then he kills even this, his absolute or God. Substance =
being without life—because he [116–117] forgets his very own act of insight—
the life in which the science of knowing as a transcendental philosophy makes
its entrance. (“Atheist or not atheist?” Only those can accuse the science of
knowing of atheism—I am not concerned here with real events, 7 because
regarding all those the science of knowing is not at issue, since in fact no one
knew anything about it—who want a dead God, inwardly dead at the root,
notwithstanding that after this it is dressed up with apparent life, temporal
existence, will, and even sometimes with blind caprice, whereby neither its life
nor ours becomes comprehensible; and nothing is gained except that one more
number is added to the crowd of finite beings, of whom there are more than
enough in the apparent world: one more that is just as constrained and finite
as themselves and that is in no way different from them in kind.—[I mention]
this in passing to state clearly and in a timely way a significant basic quality of
the science of knowing.) 8
One term is being, the other—the negated concept—is without doubt
subjective thought, or consciousness. Therefore we now have one of the two
basic disjunctions, that into B and T (being and thinking), we have grasped
this in its oneness, as we should, and as proceeding completely and simply
from its oneness, (L = light); and thereby, so that I can add this, too, paren-
thetically, we would simultaneously have the schema for the negation of the I
in the pure light and even have it intuitively. 9 For if, as everyone could easily
agree, one posits that the principle of the negated concept is just the I (since
I indeed appear as freely constructing and sketching out the concept in
response to an invitation), then its destruction in the face of what is valid in
itself is simultaneously my destruction in the same moment, since I as its prin-
ciple no longer exist. My being grasped and torn apart by the manifestness
which I do not make, but which creates itself, is the phenomenal image of my
being negated and extinguished in the pure light.
2. This, I say, is a result of the light itself and its inward living expres-
sion: things must remain here as a consequence of this insight and in case they
simply follow it, and we will never get beyond it. But I claim that, if only we70
The Science of Knowing
reflect correctly, we are already beyond it: we have certainly considered the
light and objectified it: the light therefore—whomever has forgotten this cir-
cumstance during the previous explanation should recall it now—[118–119]
has a twofold expression and existence, partly its inner expression and exis-
tence (conditioned by the negation of the concept, conditioning and positing
absolute being), partly an external and objective expression and existence, in
and for our insight.
As concerns the latter, 10 that at first we speak of it alone; we surely
remember that we did not possess it, and everything that lies within it, imme-
diately; instead we raised ourselves to it from the beginning of our investiga-
tion, initially by abstracting from the whole variety of objective knowings [and
rising] to the absolute, self-presenting insight that genuine knowing must
always remain self-same in this variation; and then [we continued] by means
of deeper, genetic examination of this insight itself. So far this has been our
procedure; the new and unknown spiritual world in which we pursue our path
has arisen by this procedure alone, and without it we would be speaking of
nothing. It now further appears to us that we could very properly have
neglected this procedure, just as we have undoubtedly neglected it every day
of our lives before we came to the science of knowing. Taking up this appear-
ance now—while not making any further inquiry into its validity or invalid-
ity—[we find that] it contains the following: the light’s external existence in
an insight directed to it, as the one absolute, 11 eternally self-same in its funda-
mental division into being and thinking; [it] is conditioned by a series of
abstractions and reflections that we have conducted freely, in short by the pro-
cedure that we state as the free, artificially created science of knowing; this
external existence arises only in this way and for it, and otherwise not at all. This
the first point here. 12
But, on the contrary, we assert, as concerns the inner existence and
expression of the light, that if the light exists unconditionally—and in partic-
ular whether we have insight into it or not, and this is the very insight which
depends on appearing freedom—it is in and for itself the very same one, eter-
nally self-identical, and thoroughly necessary, if only the light exists. There-
fore we assert a meaningful consequence, something I bid you to mark well,
that there are two different modes of light’s life and existence: the one medi-
ately and externally in the concept, the other simply immediate through itself,
even if no one realizes it. Strictly speaking in actual fact no one ever does real-
ize it, but instead this inner life of the light is completely inconceivable. This is
the second point. 13
Light is originally divided into being and thinking. That the light lives
unconditionally therefore means that it splits itself completely originally, and
also inexplicably, into being and concept, which persists, even though it is
negated qua concept. To be sure, insight can follow this very split, just as it isEighth Lecture
71
now on our [120–121] side following reconstructively the split into concept qua
concept and being qua being; but at the same time insight must leave the inner
split standing as impenetrable to it. This yields, in addition to the previously
discovered and well-conceived form of inconceivability, a material content of
light that remains ever inconceivable.
(I have just expressed myself on a major point in the science of knowing
more clearly than I have previously succeeded in doing. We would accomplish
a great deal if this became clear to us right now on the spot.—That the light
lives 14 absolutely through itself must mean: it splits itself absolutely into B
(being) and T (thinking). But “absolutely through itself ” also means “indepen-
dent of any insight [into itself ] and absolutely negating the possibility of
insight.” Nevertheless for the last several lectures we have seen, and had insight
into, the fact that light splits itself into B and T: consequently this split as such
no longer resides in the light, as we had thought, but in the insight into the
light. What then still remains? The inward life of the light itself in pure iden-
tity, from itself, out of itself, through itself without any split; a life which exists
only in immediate living and has itself and nothing else. “It lives”; and thus it
will live and appear and otherwise no path leads to it.—“Good, but can you not
provide me with a description of it?” Very good, and I have given it to you; it
is precisely what cannot be realized, what remains behind after the completely
fulfilled insight which penetrates to the root, and therefore what should exist
through itself. “How then do you arrive at these predicates of what cannot be
realized—i.e., is not to be constructed from disjunctively 15 related terms as being
is from thinking and vice versa—predicates such as that it is “what remains
behind after the insight,” “that which ought to exist through itself,” qualities
which are the content or the reality you have claimed to deduce fundamentally?
Manifestly only by negation of the insight: hence all these predicates, leading
with the most powerful—absolute substance—are only negative criteria, in
themselves null and void. “Then your system begins with negation and death?”
By no means; rather it pursues death all the way to its last resort in order to
arrive at life. This lies in the light that is one with reality, and reality opens up
in it. And the whole of reality as such according to its form is nothing more
than the graveyard of the concept, which tries to find itself in the light.) 16
It is obvious that our entire enterprise has achieved 17 a new standpoint
and that we have penetrated more deeply into its core. The light, which up to
now has been understood {eingsehen} only in its form as self-creating mani-
festness [122–123], and hence assigned only a merely formal being, has trans-
formed itself into one living being without any disjunctive terms. What we
have so far assumed to be the original light has now changed itself into mere
insight and representation of the light; and we have not merely negated the
concept that has been recognized as a concept, but even light and being as
well. Previously only the mere being of the concept was to have been negated;72
The Science of Knowing
but how were we supposed to have arrived at this being even though it is
empty? It was to be negated by something that itself was nothing. How could
that be possible? Now we have an absolute reality in the light itself, out of
which perhaps the being that appears, as well as its nonexistence in the face of
the absolute, might be made comprehensible.
I now explicitly mention in addition what direct experience teaches in
any case, that this reality in the original light, as it has been described, is
unconditionally and completely one and self same, and that no insight is
allowed into how it might arrive inwardly at a division and at multiplicity.
Observe: the division into being and thinking—as well as what might, fol-
lowing previously given hints, depend on this—resides in the concept which
perishes in the presence of reality and so has nothing at all to do with reality
and the light. Now, according to the testimony of appearance in life to which
our system has provisionally granted phenomenological truth, another dis-
junction ought to arise which stands either higher, or at least on the same level
with being and thinking, since it ranges across both of the latter; and this is
taken for a disjunction in reality. Since this last contradicts our previous
insight it is thus certainly false; hence this new grounds for disjunction must
lie in a determination of the concept that has not yet been recognized or is not
yet sufficiently explored. The concept, as a concept, must itself be conceivable,
and so no new inconceivability can appear here. But if this determination of
the concept is grasped conceptually, then everything that it contains can be
derived conceptually from it. Whatever range of differences may come for-
ward in appearing reality now and for all time, yet it is once and for all clear
a priori that they are B—T + C + L; 18 one-and-the-same, remaining eternally
self-identical, and only different in the concept. Therefore, it is clear that,
since everything true must begin with it and that falsity and illusion must be
turned away, reality (with which alone true philosophy can be concerned), not
only is generally completely deduced and made comprehensible, but also
divided and analyzed a priori into all its possible parts. “Into its parts,” I say,
excluding from this L (= the light). For in fact this is not a part, but the one
true essence.—It is hereby likewise clear how far the deduction and re-con-
struction of true [124–125] knowing goes in the science of knowing: insight
can have insight into itself, the concept can conceive itself; as far as one
reaches, the other reaches. The concept finds its limits; conceives itself as lim-
ited, and its completed self-conceiving is the conceiving of this limit. The
limit, which no one will transgress, even without any request or command
from us, it recognizes exactly; and beyond it lies the one, pure living light;
insight points therefore beyond itself to life, or experience, but not to that mis-
erable assembly of empty and null appearances in which the honor of existing
has no part: but rather to that experience which alone contains something
new: to a divine life.C H A P T E R
9
Ninth Lecture
Wednesday, May 2, 1804
Honored Guests:
In the next three lectures I am about to enter into a deeper investigation than
has been made so far. This investigation will, as it happens, set out to secure a
stable focus and, leading from it, a permanent guide for our science, even
before we possess this guide. So, in order not to get confused, a lot depends
on our holding on to what we have laid down provisionally; therefore:
1. Formally, i.e., in relation to the material which we are investigating and
to the manner in which we take it, we are already located beyond the prole-
gomenon and actually inside the science of knowing; because (the previous lec-
ture began by recalling this) we have already actually created insights in our-
selves, which have transposed us into an entirely new world belonging to the
science of knowing and raised above all factical manifestness, in whose realm the
prolegomenon always remains. We have passed unnoticed out of the prole-
gomenon and into science; and indeed the transition started as follows: we had
to elucidate the procedure of the science of knowing by examples, and, since I
found that the state of the audience 1 made it possible, we made use of the actual
thing as the original example. Let us now drop this as a mere example and take
things up earnestly and for real; thus we are inside the science. Just as this has
so far happened tacitly, let us now proceed conscientiously and explicitly.
2. Here is how things stood in the hour before last: I—L—B. a (a = our
insight into the matter). I (Image), positing something imaged in it, = B
(being) and vice versa; united in the oneness of the light (L). Thus—on the
one side, the connection of I—L—B, the essential element of all light without
exception: on the other side, the modifications 2 without which it does not exist.
This effectively indicates the way in general, but nothing is still specially
known thereby. It is only the prolegomenon to our investigation. 3
Additionally, this gives a good hint concerning an important point
which is not to be handled without difficulty as to its form. Knowing should
This lecture begins at GA II, 8, pp. 126–127.
7374
The Science of Knowing
divide itself entirely at one stroke according to two distinct principles of divi-
sion: B–T / oneness, and x, y, z / oneness. Here we see that the light, in itself
eternally one and self-identical [128–129], does not divide itself in itself, but
rather divides itself, in its insight and as being seen, into this multiplicity,
whatever x, y, z may be; the [same] light, which, in itself and in its eternal
identity independent from insight into it (at least as we have posited this more
deeply), divides itself into being and thinking. Therefore, if the light does not
even exist except in being the object of insight, this again divided; likewise the
light does not exist in itself without dividing itself into being and thinking, so
this disjunction is absolutely one and indivisible according to both distin-
guishing grounds. At present things must remain here, and this proposition,
together with all further qualifications which it may yet receive, true in itself
and remaining true, will never be permitted to fall. ( Just by virtue of the fact
that one has fixed termini in the conduct of the investigation, one is able to
follow the investigation’s most divergent turns without confusion, and to ori-
ent oneself in it as long as the point at which everything ties together remains;
while otherwise one would very quickly be led into confusion.)
Now, in regard to the concept—which lies neither in the light, as what
is imaged for the concept, nor in the insight, as the image itself, but rather
between these two—we realize that formally in itself this concept is a mere
through-one-another {Durcheinander}, without any external consequences, i.e.,
without antecedent and without consequent, which two, and all their shifting
relations, arise only out of the living exhibition of this concept. This insight,
which, if I am not mistaken, has been presented with the highest clarity, is
presupposed and here only recalled for you. If I wished to add something here
to sharpen this insight, then it could only be this: since the concept, as
absolute relation of the imaged to the image and vice versa, is only this rela-
tionship, it makes no difference to it that the thing imaged should be self-suf-
ficient light and that the image should be this image. Something imaged and
the image, simply as such, are sufficient. Further, the imaged thing and the
image are also of little concern to the concept’s inner essence, the latter presup-
posed as absolutely self-sufficient; instead this inner essence is evidently a
mere through-one-another. That this through-one-another, as simply existing,
manifests in the image and the thing imaged has shown itself empirically. But
who then authorizes us to say, on the one hand either that this through-one-
another must manifest itself or exist, or, on the other hand in case the former
should be true, that it must construct itself directly in the image and the thing
imaged rather than, say, construct itself for another and under other condi-
tions in an endless variety of ways? Through this consideration we lose the
subordinate terms, and their distinguishing grounds, for a system of genetic
knowing. [130–131] On the other hand, in case someone is willing to grant
us this, who would then authorize us to assume that the thing imaged couldNinth Lecture
75
only be the light, and that therefore necessarily the image of the light, which
arises in the concept as its imaged object and by means of it, must thereby
bring in the other distinguishing ground ? Thereby we also lose the second half
in any system which does not rest content with factical manifestness, and
rejects everything that is not seen genetically as necessary.
Of course, this result comes out the only way it could, as soon as we
think seriously. If we posit the concept, the absolute through-one-another, as an
independent self-subsisting being, then everything external to it disappears
and no possibility of escaping it is to be found; just as things happened previ-
ously with the light when we likewise posited it. That is obvious. Any inde-
pendent being annuls any other being external to it. Whenever you might
wish to posit a being of this sort, it will always similarly have this result, which
resides in its form.
This observation provides exactly the right task for our further proce-
dure; and I wish that we could come to know this procedure in its unity right
now in advance, so that we would not go astray among the various forms and
changes which it may assume as we go along, would easily recognize the same
pathway in every possible circumstance only with this or that modification,
and would know which modification it was and from whence it comes.—The
genetic relation whose interruption has come to light must be completed. This
cannot simply be done by inserting new terms and thereby filling the gap, for
where would we get them? We are scarcely capable of adding something in
thought where nothing exists. Therefore, the genetic relation which is cur-
rently absent must be found in the terms already available; we have not yet
considered them correctly, i.e., completely genetically, but so far still only con-
sidered them in part factically. “In the terms already available,” I say; thus, if
the only important thing were to arrive at our goal by this path, it would not
matter with which available term we began. If we worked through only one of
these to its implicit, creative life, then the flood of light which simultaneously
overcomes and connects everything would of necessity dawn in us. But
beyond this we also have the task of following the shortest way; and so it is
quite natural for us to hold on to what has shown itself to us as the most
immediate, those terms in which we alternatingly have placed the absolute,
and in regard to which we now find ourselves in doubt as to which is the true
absolute, namely light and concept. [132–133] If we work through both so that
each 4 shows itself as the principle of the other, then it is clear that
a. in each we have grasped mediately the distinguishing ground which is
immediately present in the other, and that
b. beginning from both, we, in our scientific procedure, have obtained a yet
higher common principle of distinction and oneness for both on essential
grounds.76
The Science of Knowing
Therefore, both of these lose their absolute character and retain only relative
validity. Thus, our knowledge of the emerging science of knowing transcends
them as something absolutely presupposed; and according to its external form
this is a synthesis post factum. But since this transcendence is itself genetic in
its inner essence—and it is not simply as Kant, and indeed we ourselves speak-
ing preliminarily, have said, that “there must surely be some yet higher one-
ness,” but rather this oneness in its inner essence is actually constructed—it is
a genetic 5 synthesis. But again, the science of knowing, which is genetic in its
principles and which permeates the higher oneness, is permeated by it, and is
therefore itself identical with it, steps down into multiplicity and is simulta-
neously analytic and synthetic, i.e., truly, livingly genetic. Our task is discover-
ing this oneness of L [light] and C [concept], and discovering it in this briefly
but precisely prescribed way; this discovery is the common point to which the
whole of our next stage refers. This procedure’s modifications and various
turnings are grounded in the necessity now of properly permeating C by
genetically permeating L, and then again the other way. Thus, [it grounds
itself ] on constant shifts in standpoint and being tossed from one to the
other. 6 I will not conceal 7 the fact that this procedure is not without difficulty
and that it demands a particularly high degree of attention; instead I
announce this explicitly. But I am overwhelmingly convinced that whoever
has actually seen into what has been presented so far, and holds fast to the pre-
sent schema and the just asserted common point for our investigation, orient-
ing himself in terms of these from time to time, will not be led astray. On the
other hand, this is the only truly difficult part of our science. The other part,
deducing the mediate and secondary disjunctions, is a brief and easy affair for
those who have properly achieved the first, no matter how monstrous and mad
it may appear to those who know nothing of the first. This second part
namely, as is evident from the foregoing and, which I mention here only
redundantly, has the task of deducing all possible modifications of apparent
reality. The individual who has so far remained trapped in factical manifest-
ness wonders at this because it is the only difficulty which is accessible and
apparent [134–135] to him. But until it has its own openly declared principle,
this deduction (of the manifold of apparent reality) is nothing more than a
clever discovery, which has recourse to the reader’s genius and sense of truth,
but can never justify itself before rigorous reason, if it does not have, and
declare, its own principle. Now to discover and clarify this principle may cer-
tainly be the right work: for one who possesses it, the application will thus
surely be simple, and—since the most complete clarity and distinctness is to
be found here—it will be even simpler than the application of principles in
other cases. Indeed, one could, if necessary, simply rest satisfied to have shown
this application through a few examples.—Since I would gladly dispose of this
once and for all, let me take it down to specific cases: the deduction of timeNinth Lecture
77
and space with which the Kantian 8 philosophy exhausts itself and in which a
certain group of Kantians remain imprisoned for life as if in genuine wisdom,
or of the material world in its various levels of organization, or of the world
of the understanding in universal concepts, or of the realm of reason in moral
or religious ideas, or even the world of minds, presents no difficulty and is cer-
tainly not the masterpiece of philosophy. Because all these things, together
with whatever one might wish to add to them, are actually and in fact nonex-
istent; instead, in case you have only just understood their nonexistence, each
one is the very easily grasped appearance of the truly existing One. To be sure,
up until now, some have freely believed in the existence of bodies (i.e., truth-
fully, in the nothing which is presented as nothing) and—at the most—in the
existence of souls (i.e., truthfully, in ghosts), and perhaps they have even con-
ducted deep researches into the relation of body and soul or into the soul’s
immortality. Let me add that not for one moment do I support skepticism
about the latter or wish to wound faith. The science of knowing legislates
nothing about the immortality of the soul, since in its terms there is neither
soul, nor dying, nor mortality, and hence there is likewise no immortality;
instead there is only life, and this is eternal in itself. Whatever exists, is in life,
and is as eternal as life is. Thus, the science of knowing holds with Jesus that
“whoever believes in me shall never die,” but it is given him to have life in
himself. But I say, picking up the thread, whoever has believed something like
this [136–137] and is used to philosophical questions of this kind demands
that a science, which says the things ours does, address this point with him
and free him from error, if only by an induction on what he has so far taken
as reality. This is what Kant 9 for instance, did; but it did not help at all. Noth-
ing could help that did not address the problem at its roots. The science of
knowing does even better than they wish, according to rigorous methods and
in the shortest possible way. It does not cut off errors individually, since it is
evident that in this work as soon as error is removed on one side it springs up
on the other; rather it insists on cutting off the single root for all the various
branches. For now, the science of knowing asks for patience and that one not
sympathize with the individual appearances of disease, which [our science]
has no wish to heal: if only the inner man is first healed, then these individ-
ual appearances will take care of themselves.
What must be built up in us today is this declaration of our proper
standpoint, and to consider the unity of our following proceedings, its coher-
ence as a first part, with a later piece that can be seen as the second part; and
in relation with this you can view everything previous as a condition of clear
insight into today’s material. Nothing is thereby gained for material insight
into our investigation’s object; there is indeed a very important point in this
insight which we found last time while doing something else {aliud agendo},
dropped today as not relevant to our purpose, and which we will investigate78
The Science of Knowing
again tomorrow for the purpose we have announced clearly today. As concerns
the form, however, a general perspective and orientation has been achieved
which will guard us from any future confusions. The schema serves as provi-
sionally valid and it will endure only those revisions grounded in our growing
insight and no capricious ones.
In conclusion, in order to point out to those who have attended my pre-
vious course 10 where they find themselves in the process and thereby to put
them in a position [138–139] to view the science of knowing with the com-
plexity that repeatedly pursuing it allows: what I am now calling concept was
named in the first series “inner essence of knowing,” what is called here “light”
was there called “knowing’s formal being,” the former [was called] simply the
intelligible, the latter intuition. For it is clear that the inner essence of know-
ing can only be manifested in the concept, and indeed in the original concept;
likewise, [it is clear that] that this concept, as implicitly insight, must posit
insight, or light. Therefore, it is clear that the task here expressed as “finding
the oneness of C and L” is the same one expressed there by the sentence: “The
essence of knowing [is] not without its being, and vice versa, nor intellectual
knowing without intuition and vice versa, which are to be understood {einge-
sehen} so that the disjunction that lies within them must become one in the
oneness of the insight.” Recall that we have concerned ourselves with this
insight for a long time, and that it has returned under various guises and in
various relations, but always according to synthetic rules. Certainly it could
not happen any differently in this case, and it was this something, surrender-
ing itself even then, which I meant when I spoke previously about the mani-
fold shifts and modifications of the one selfsame process. The difference from
before—and it seems to me also an advantage of the present path—is this: that
already from the start, even before we plunge into the labyrinth of appearance,
we can recognize our various future observations in their spiritual oneness. It
is to be hoped—and this hope does not really concern my own knowledge and
procedure in lecturing but rather the capacity of this audience to follow the
presentation—that an ordering principle for these various shifts will soon be
available, by means of which the process will be further facilitated. And so it
will not be difficult for this part of the audience to recognize in what is now
expressed in a particular way what was said in a different way before and vice
versa. In being liberated from my two different literal presentations {Buch-
staben} they may free themselves from any literal presentation, which would
mean nothing and would be better not existing if it were possible to hold a
lecture without one. And in freeing themselves, [they] build realization for
themselves in their own spirit, free from any formulas and with independent
control in every direction.
Let me add the following while we still have time, although it is not
essential and has relevance only to the smallest circle of those gathered here:Ninth Lecture
79
except for the fact that one [140–141] possesses as one’s own genuine prop-
erty only that which one possesses independently of the form in which one
received it, one can intentionally present it afresh and share it only under such
conditions. Only what is received living in the moment, or not far removed
from it, strikes living minds; not those forms which have been deadened by
being passed from hand to hand or by a long interval. If I had needed to hold
these lectures on the science of knowing immediately after the previous series
to the same audience, who had all known the science for a long time and had
no need for further instruction in it, and who simply wanted to prepare them-
selves further for their own oral presentation of philosophy, yet I still believe
that I would have been required to take almost as diverging a path as I have
taken this time, and I would have had to advise these future teachers of phi-
losophy about the utility of this divergence in just about the same way I have
advised you, for whom it is relevant.C H A P T E R
10
Tenth Lecture
Thursday, May 3, 1804
Honored Guests:
Our next task is now clearly determined: to see into L (light) as the genetic
principle of C (concept) and vice versa, and thus to find the oneness and dis-
junction of the two. (Let me add yet another parenthetical remark. Who
among you, prior to studying the science of knowing, has known L or C, not
in general and with confusion—since any sort of philosophy distinguishes an
immediate presentation of an actual object and a concept, which is usually an
abstract one—but true L and C in the purity and simplicity with which they
have been presented here? Our task concerns itself with doing this; and, with
the resolution of this task, the science of knowing is completed in its essentials.
Accordingly, the science answers a question that it itself must pose, and dis-
solves a doubt that it has first raised. It should not seem strange to anyone that
there is no bridge to it from the usual point of view and that one must first
learn everything about it from within it.) Something has happened for the res-
olution of this task on Monday, 1 which we will now briefly review to confirm
our grasp of it.
As factical manifestness makes clear, light plainly arises in a dual rela-
tion: in part as inwardly living—and through this inner life of its own it must
divide itself into concept and being—and in part in an external insight, 2 which
is freely created and which objectifies this light along with its inner life. 3 Let
us take up the first. What makes this inward life inward ? Obviously that it is
not external. But it becomes external in [being seen by] the insight. Thus,
what follows immediately and is synonymous: it is an [inward] life because in
this regard it is outside any insight, is inaccessible to it, and negates it. There-
fore light’s absolute, inner life is posited; it exists only in living itself and not
otherwise; therefore it can be encountered only immediately in living and
nowhere else. I said that the genuinely, truly real {Reale} in knowing rests here.
But we ourselves have just now spoken of this inner life and therefore in some-
This lecture begins at GA II, 8, pp. 142–153.
80Tenth Lecture
81
way conceived it. Yes, but how? As absolutely inaccessible to insight; so we
have conceived and determined it only negatively. It is not conceivable in any
other way. The concept of reality, of the inner material content of knowing,
etc., which we have introduced is only the negation of insight and arises only
from it; and this should not just be honestly admitted, rather, a philosophy
which truly understands its own advantage should carefully enjoin this idea.
In truth it is no negation, but rather the highest [144–145] affirmation {Posi-
tion}, which indeed is once again a concept; but in truth we in no way conceive
it, but rather we have it and are it. 4
Let this be completed and determined by us right now, and in this act
it will have its application uninterruptedly. And don’t let the truly crucial
point of the matter escape: [there are] two ways for the light to live
absolutely: internally and externally; externally in the insight, internally,
therefore, absolutely not in the insight, and not for it, but instead turning it
away.—By this means, our system is protected against the greatest offense
with which one can charge a philosophical system, and without exception
nearly always justly: namely vacuity. Reality, as genuine true reality, has been
deduced. No one will confuse this reality with being (objectivity); the latter
is subsistence-for-self and dependence-on-self which is closed in on itself and
therefore dead. The former exists only in living, and living exists only in it; it
can do nothing else than live. Therefore, because our system has taken life
itself as its root, it is secured against death, which in the end grasps every
[other] system without exception somewhere in its root. Finally, we have seen
{eingesehen} and enjoined that, since light and life too are absolutely one, this
reality [and the insight, through the negation of which it becomes reality,] 5
can altogether be only one and eternally self-same. [146–147] Thereby our
system has won enduring oneness and has secured itself from the charge that
there may still be some duality in its root.
Insight, I assert, is completely negated in living light. 6 But then we see,
and see into, the disjunction into C and B. Therefore, this disjunction, which
we previously ascribed to the inner light itself, should not be ascribed to this,
but instead to the insight that takes its place, or to the original concept of
light. The concept reaches higher, the true light withdraws itself. The
absolute negation of the concept may well remain a nothing 7 for the science
of knowing, which has its essence in concepts, and only become an affirma-
tion in living. 8 With this, two further comments which belong to philoso-
phy’s art and method.
1. Here we retract an error in which we have so far hovered. How did
we arrive at this error, or at the proposition that we now take back as erro-
neous? Let us recall the process. Driven by a mechanically applied law of rea-
son (therefore, factically), we realized immediately that it [i.e., the absolute]
cannot reside in A (the oneness of B and T) nor in the disjunction point, but82
The Science of Knowing
rather in the oneness of both; this was the first step. Then, as the second step,
we raised ourselves to the apprehension of the general law for this event, which
naturally we could apprehend only as follows: in an immediately self-present-
ing insight, a disjunction is negated as intrinsically valid and a oneness, which
cannot be any more exactly described, is absolutely posited. 9 What then did we
finally do? In fact we did nothing new, apart from the fact that we relin-
quished the specificity of the disjunctive terms “A” and “.” and likewise the
specificity of their unity, and posited disjunction, and also self-sufficient one-
ness, generally and unconditionally, in which case the possibility of the proce-
dure could arouse wonder and give rise to a question. Besides this, I say, we
simply grasped the rules of the event historically, always led by the event, and,
if that were removed from us, 10 [we] lacked all support for our assertion.
Therefore, although this, our second insight, appears to possess a certain
genetic character in the first mentioned part {Ingredienz}, in the second part
it is something merely factical; and so what we advanced yesterday as the
ground for the uninterrupted connection between the disjunctive terms is
confirmed here at a genuinely central point: [148–149] that our whole insight
might not yet be purely genetic, but is instead still partially factical. To get back
to the point 11 —this insight, arising in concrete cases and led to its general rule
in the second step, we now named pure, absolute light, simply in this respect,
that in terms of content it arises immediately, without any premises or condi-
tions. But in its form it remains factical and is dependent on the prior com-
pletion in concrete cases. We might have inferred from the following that it
cannot possibly have its application here: although dividing the concept into
“A” and “.” was given up as inadequate, there yet remained a new disjunction
in what was taken as absolute, since it was simultaneously negating and posit-
ing, the former through its formal being, and the latter through its essence.
But no disjunction can be absolute and merely factical, rather, as surely as it is
a disjunction, it must become genetic, since disjunction is genetic in its root.
(Remarks of the kind just made bring no progress in the subject matter, but
they elevate the freedom of self-possession and reproduction for everyone and
facilitate the comprehension of what follows.) 12 Result: since our initial sup-
position grounded itself partially on factical insight, we must give it up.
Further, how then have we arrived at this insight of giving something
up, as well as at the higher [term] for which we give it up? If you recall, we
did so by means (i) 13 of the distinction between two ways the light exists and
lives: inwardly and outwardly, a distinction admittedly given only with factical
manifestness; 14 (ii) by genetic insight into this distinction and by the question
how something absolutely inward might arise as inward; and (iii) by elevating
into a genetic perspective something that had been thought previously only in
faded, factical terms. Moreover, I admitted, as is indeed evident and as every-
one will remember, that this entire disjunction between inner and outer arisesTenth Lecture
83
{liegt} only in a factical point of view. Here, too, this observation: that in our
present investigation as well, the facticity, which is erased on one side, pops up
again on the other, and that we will not be entirely and purely relieved of it
until the present task is completed. 15 (For returning listeners this as well: the
distinction drawn here between the light’s inner and outer life is the same as
the distinction between immanent and emanent forms of existence that was so
important and meaningful in the previous series.)
[150–151] Second remark: 16 C and L are both only concepts: the first is
purely disjunction in general, a disjunction which can give no further account
of itself; i.e., from our current point of view, it has simply two terms which are
not further distinguishable. L, on the other hand, is not a disjunction in gen-
eral, but is rather the specific disjunction into being and concept. 17 The latter,
as the principle of disjunction in general, consequently has enduring inner con-
tent, as does B [Being], the principle of oneness. Therefore, the terms of this
second disjunction are not just two terms in general, they also have an inter-
nal difference. From our standpoint, therefore, L is still by no means negated,
nor can it be from that perspective. If nevertheless it must be negated, as it
evidently must be a priori, since otherwise it could not come to the zero state 18
in which no further disjunction truly remains, then completely different
means will have to be employed from the ones now available to us. 19
Now to characterize (in relation to our task) the point I have just
repeated, a point which fits our system in every aspect, of which I said last
time that it already belonged to our process, and which we need to apply in
solving our next and primary task: C and L are to be reduced to oneness, just
as they were before this point of ours; and this will have to be done so that C
is so rigorously mastered that we see into it as the genetic principle of light,
and vice versa. With which of the terms to begin is left either to our caprice or
our philosophical skill, which is unable to give any account of its maxims
before applying them. In the previously discussed point, L was taken up as a
starting point, as things stood then; [152–153] if so, the concept proceeds
genetically from this L, since L has transformed itself into the concept. Or
said more exactly: our own observing—which was not then visible, which we
lived, and into which we merged—divided itself beyond the then regnant L,
and in this division negated the L (light) into 0/C; 20 thus creating both out of
itself. Now note well that this shift in viewpoint is in no way merely a change
of the word and the sign, but rather that it truly is a real change; because what
stood here previously, whether called L or C, light or concept, was to be the
absolute (which is a real predicate) and should divide itself into C and B 21
(which is also a real predicate). Both predicates combine to form a synthetic
sentence that determines the absolute. In its essence—entirely apart from the
expressions and signs in which one realizes and presents this essence—this
sentence is contradicted by the really opposite sentence: the principle of the84
The Science of Knowing
disjunction into being and thinking is not absolute but something subordi-
nated (however one may more exactly name and signify this subordinate
something); in the absolute the two are not distinguished. From this,
another correction must follow first and foremost, not so much regarding
our viewpoint as our manner of expression. There were supposed to be two
different distinguishing grounds, which of course are to be reunited again,
but which would have been held sufficiently far apart from one another by
two basic principles that were as distinct as until now light in itself and its
representative concept have been. Now all disjunction collapses into one and
the same concept, and hence this latter could very easily provide the one
eternally self-same disjunctive factor, which does not appear in the original
appearance, but rather appears as doubled in the secondary appearance, in
the appearance as appearance.
But let me go back. As things stood previously, the spirit of our task was
to realize L as the genetic principle of C and vice versa. We have tried by
beginning with L: the attempt had its narrowly circumscribed result and the
matter does not stand as it did before but as the schema instructs. The spirit
of the task remains the same through all shifts in perspective, just because it
is spirit: L through C and vice versa. Our true L is now = 0, [154–155] and it is
clear that this cannot be approached more closely: it negates all insight.
Therefore, this first path is already completed on the initial attempt. Nothing
more remains besides taking up C and testing whether through it we can fur-
ther determine—not 0, since this remains purely unchangeable and indeter-
minable,—but rather it 22 as the truly highest term that we now are and live.
Thus—a new classificatory division, the determination on the basis of C,
becomes the second principal part of our present work.
Let me now give some preliminary hints and thereby prepare you for
tomorrow’s lecture, setting out a rough outline for you.
The concept’s inward and completely immutable essence has already
been acknowledged in an earlier lecture as a “through.” 23 Although in its con-
tent this insight is in no way factical but is instead a purely intellectual object,
still it has an factical support: the construction of the image and the thing
imaged, and the indifference of the inference between them. However, we
would be permitted to use this basic quality of the concept, if only, in this
application, we succeeded in negating its factical origin. If one embraces a
“through” just a little energetically, it can readily be seen that the same princi-
ple is a disjunction. Except the same question must always be repeated which
already arose previously on the same occasion: how should a dead “through,”
defined as we have defined it, come to life—despite all the capability with
which it is prepared to meet life, especially by means of the “throughness”
{Durchheit}, or the transition that it makes from one to another, if only it is
brought into play—because it has no basis in itself for coming to actualiza-Tenth Lecture
85
tion? How would it be if the internal life of the absolute light (= 0) were its 24
life, and therefore the “through” was itself first of all deducible from the light
by this syllogism: 25
i. If there is to be an expression—an outward existence of the immanent life
as such—then this is possible only with an absolutely existent “through.”
ii. But there must be such an expression.
iii. Hence, the absolute “through” (i.e., the original concept, or reason) exists
absolutely as everyone can easily see for himself.
Further, how would it be if just this living “through” (living to be sure by
an alien [156–157] life, but still living) as the oneness of the “through,” divided
itself into thinking and being, i.e., in itself, and in the origin of its life? This divi-
sion, as a division of the enduring “through” as such, would be comprehensive
for the same reason, and inseparable from it and its life. How would it be if it
did not remain trapped in this, its essence as “through,” but rather might itself
be objectifying and deducing the latter; it would hence certainly have to be
able to do just as we ourselves have done—the objectification and deduction
can themselves come about according to the law of the “through,” 26 since fun-
damentally and at base it 27 is nothing but a “through”: how would it be if in
this objectification and deduction it split itself again in the second way? [Fur-
ther, how would it be, since a “through” clearly can exist only by means of a
“through”—that is, its own being as a “through” can be only mediate—and the
first mode of being would scarcely be possible without at least a little of the
second, if the first division could not be without some of the second, and vice
versa? Since this “through” is our own inner essence, and a “through” dissolves
completely into another “through,” then absolutely everything based in these
terms must be completely conceivable and deducible.] 28 In all of these “how
would it be if ____” clauses, I have consistently regarded “0” as life; but it is
not merely this, rather it is something indivisibly joined with life, a thing we
grasp by the purely negative concept of reality. If it is indivisible from life, and
if life lives in a “through,” 29 then it lives as absolute reality, but since it is a
“through,” it lives it in a “through” and as a “through.” Now, consider what fol-
lows if the one absolute reality, which can only be lived immediately, occurs in
the form of an absolute “through.” I should think this: that it cannot be grasped
anywhere, unless an antecedent arises for what has been grasped, through
which it is to be; and, since it is grasped only as a “through,” it must also have
a consequence that is to follow from it. This must follow unavoidably by
absolutely every act of apprehending reality. In short, the infinite divisibility
in absolute continuity—in a word, what the science of knowing calls quantifi-
ability 30 as the inseparable form of reality’s 31 appearance—arises as the basic
phenomenon of all knowing.86
The Science of Knowing
In this last brief paragraph of my talk I have pulled together the entire
content of the science of knowing. Whoever has grasped it and who can see
it as necessary—the premises and conditions for this manifestness have
already been completely laid out—such a one can learn nothing more here,
and he can only [158–159] clarify analytically what has already been seen into.
Whoever has not yet seen into it has at least been well prepared for what is to
follow. For the one as for the other, we will move forward tomorrow.C H A P T E R
11
Eleventh Lecture
Friday, May 4, 1804
Honored Guests:
Yesterday, I succeeded in presenting the essence and entire contents of the sci-
ence of knowing in a few brief strokes. [Losing time at the right place means
gaining it; therefore, against my initial purpose, I will apply today’s hour to
presenting further observations about this brief sketch. The more certain we
are in advance about the form, the easier the actual working out of the con-
tents within this form will be.] 1
C 2 = “through,” in which resides disjunction. “If only this “through”
could be brought to life,” I said; it has nearly all the natural tendencies of life,
nevertheless in itself it is only death. It would be useful to reflect further about
this expression, since the “through” can be more clearly understood in it, than
it has yet been understood: 3 this “through,” which according to the preceding
represents the central point in our entire investigation. Indeed, it is immedi-
ately clear what it means to say: “ a “through” actually arises,” 4 “a ‘through’ has
taken place,” or “there is an existing ‘through.’” I further believe it will be clear
to anyone who considers the possibility of this existence that, considered for-
mally, something else belongs to it besides the pure “through.” In the “through”
we find only the bare formal duality of the terms; if this is to find completion,
then it needs a transition from one to the other, thus it needs a living oneness
for the duality. From this it is clear that life as life cannot lie in the “through,”
although the form which life assumes here, as a transition from one to the
other, does lie in the “through”:—so life generally comes entirely from itself
and cannot be derived from death.—
Result: the existence of a “through” presupposes an original life, 5
grounded not in the “through” but entirely in itself. 6
We see into this at once: 7 but what is contained in this insight? Evidently
the insight formed in positing the “through’s” existence (and the question
about the possibility of this existence) brings life with it, that is in the image
This lecture begins at GA II, 8, pp. 160–161.
8788
The Science of Knowing
and concept. Therefore, in this insight life is grasped in the form of a
“through,” i.e., only mediately. The explanation of the “through” is itself a
“through.” The first of these posits its terms at one stroke; and, in the result-
ing insight, it is itself posited as positing them in one stroke by the explana-
tory “through” (horizontally arranged): 8
a
——
a × b
So too, in the same way, with regard to the inner meaning, 9 sense, or content,
the first “through” does not posit its terms at one stroke; rather, life should be
the condition and the “through’s” existence [162–163] should be what is con-
ditioned; thus it [i.e., life] should be in the concept as a concept—in truth and
in itself—the antecedent, and the latter should be the consequent. [This is
the] perpendicular arrangement. Both obviously [exist] only in connection with
each other, and [are] only distinguishable in this context.
The concept remains the focus of everything. (To reconstruct: here in a
certain respect to preconstruct.) 10 [The concept] constructs a living “through,”
and, to be sure, does so hypothetically. Should this latter be, then the existence
of living follows from it. It is immediately clear that a hypothetical “should” is
not grounded on any existing thing, but is rather purely in the concept, and col-
lapses if the concept collapses; and that, therefore, the concept announces
itself in this “should” as pure, as existing in itself, and as creator and sustainer
from itself, of itself, and through itself. The “should” is just the immediate
expression of its independence; but if its inner form and essence are indepen-
dent, then so are its contents as well. Hence, the existence of a “through”
announces itself here as completely absolute and a priori, in no way grounded
again on another real existence which precedes it. Therefore, the concept is
here the antecedent and absolute prius to the hypothetical positing of the
“through’s” existence: the latter is only the concept’s expression, something
which depends on it and through which it, as concept, preserves itself as an
absolutely inward “through.” Which was the first [point]. 11
In this, its vivacity, the concept changes itself into an insight, 12 which,
unconditioned, produces itself—insight into a life in and of itself, which must
necessarily be presupposed. Ascending, I can therefore say: the absolute con-
cept is the principle of the insight, or intuition, into life in itself, that is to say
into [life] in intuition. 13 It seems quite possible, namely in a shallow and faded
way, to think the existence of a “through” without any insight arising into the
life that it absolutely presupposes. For the latter to happen, this existence must
be conceived with full energy and vivacity: Now, I say, (as is clear right away):
in this pale imitation the “through” is not really thought as it must be thought:Eleventh Lecture
89
that is, as a genetic principle. For if it were thought in that way, then what
ought to be manifest would be manifest. So the true focus, the genuine ideal
prius is no longer just the concept, but rather the inward life whose outcome
(posterius) is the concept. And the “should” is not, as I said before, the highest
exponent of the independence of reason, [164–165] rather the appearance of
inner energy is this highest exponent. (If I ask you to think energetically, I am
really demanding that you be fundamentally rational!) The hypothetical
“should” is again an exponent of this exponent, and the concept is not, as I said
initially, the principle of intuition, but rather the inner immediate life of rea-
son, merely existing and never appearing, which appears as energy (energy
which obviously is again the expression of a “through” immanent in itself ).
This inner life, I say, is the principle of concept and intuition at once and in
the same stroke:—thus it is the absolute principle of everything. This, I say,
would be the idealistic argument. 14
Once we have proceeded in this way, let us climb higher in order to get
to know the real spirit and root of this mode of argument. Without further
ado, it is obvious that we could have expressed our entire procedure thus: there
may be the intuition of an original and absolute life, but how and from what
does it come to be? Just construct this being as I have, or grasp it in its becom-
ing;—now this has happened actually and in fact, and the inner life of reason
as a living “through” has been deduced as the genetic principle for this being.
Thus, the basic character of the ideal perspective is that it originates from the
presupposition of a being which is only hypothetical and therefore based
wholly on itself; and it is very natural that it finds just this same being, which
it presupposes as absolute, to be absolute again in its genetic deduction, since
it certainly does not begin there in order to negate itself, but to produce itself
genetically. Thus, the maxim of the form of outward existence 15 is the principle
and characteristic spirit of the idealistic perspective. By its means, reason,
which we already know very well as a living “through,” becomes the absolute;
it becomes this in the genetic process because it already exists absolutely as the
constant presupposition. Absolute reason, as absolute, is therefore a “through” =
the form of outward existence. Prior and absolute being 16 shows itself to be
inwardly static, motionless, and dead at just this [point of going] “through,”
where it always remains; the inward life of reason, 17 which we have already
established, shows itself in this being’s hypothetical quality. 18 One need only
add now what is implicitly clear: that this idealistic perspective does not arise
purely in the genetic process, since it assumes a being as given, and that there-
fore it is not the science of knowing’s true standpoint. This is also clear for
another reason: in the idealistic perspective, reason exists, or lives, as absolute
reason. But it lives only as absolute (in the image of this “as”); hence it does
not live absolutely; its life or [166–167] its absoluteness is itself mediated by
a higher “through,” so that in this standpoint it is only derivative. So much by90
The Science of Knowing
way of a sharp, penetrating critique of the idealistic perspective. This is espe-
cially important, because beginners are easily tempted to remain one-sidedly
trapped in this point of view, since it is the perspective in which their specu-
lative power first develops.
Now let us turn things around and grasp them from the other side. If
the “through” is to come to existence, then an absolute life, grounded in itself,
is likewise presupposed. Therefore, this life is the true absolute and all being
originates in it. With this, intuition itself is obviously negated, though not
indeed as empirically given, since if we simply try to remain in an energy state
and to consider nothing further, then we will always find that we still grasp
[this state] in intuition. Let me mention in passing that this is idealism’s stub-
bornness: not to let one go further, once one has finally arrived at it. Faced
with this, since in any case idealism is something absolute, it does not allow
itself to be explained away by any machinations of reasoning, but rather yields
only to the arrival of what is primordially absolute. Among other things, this
idealistic stubbornness too has been attributed to the fantasy of the science of
knowing which circulates among the German public, disregarding the fact
that of course people cannot speak clearly about this charge because they do
not know the genuine science; e.g., Reinhold did this for his whole career. It
is like this: the non-philosopher or half-philosopher forgets himself, or the
absolute intuition, either because he never knew it or, if he knew it, he peri-
odically forgets it again. The one-sided idealist who knows it and holds it fast
does not let it develop, because he knows nothing else.—To return: by recog-
nizing the absolute immanent life we negate intuition as something that is
genetically explicable and [that plays a part] in a system of purely genetic
knowledge. [168–169] Because if the immanent life is self-enclosed, and all
reality whatsoever is encompassed in it; then not only can one not see how to
achieve an objectifying and expressive intuition of it, one can even see that
such an intuition could never arise—and, because of its facticity, just this last
insight cannot itself be comprehended anew, but simply directly carried out; it
is the absolute, self-originating insight. So, however stubbornly one might
hold on to his immediate consciousness of this intuition, it does not help
things; no one challenges this intuition in its facticity. What has been asserted
and demonstrated is just this: not merely that [this consciousness] is incon-
ceivable, but even that it can be conceived to be impossible. Thus, the truth of
what it asserts is denied, but in no way its bare appearance.
Let me note in passing that the place for denying ourselves at the root
is here, i.e., just in the intuition of the absolute, which of course might very
well be our root, and which up to now has played that role. Whoever perishes
here will not expect any restoration at some relative, finite, and limited place.
But we do not achieve this annulment by an absence of thought and energy,
as happens in other cases. Instead, [we do so] through the highest thinking,Eleventh Lecture
91
the thought of the absolute immanent life, and through devotion to the max-
ims of reason, of genesis or of the absolute “through,” which denies its applic-
ability here and thus denies itself through itself. Everything has dissolved into
the one = 0. 19
Reasoning which is conducted and characterized in these terms is real-
istic. There is no progression or multiplicity in it except for pure oneness.—
Let me relocate you in the context. 20 The two highest disjunctive terms stand
here absolutely opposed to one another, life’s inner and outer life, the forms of
immanent and emanent existence 21 as well, separated by an impassable gulf
and by truly realized contradiction. If one wishes to think of them as united,
then they are united exactly by this gulf and this contradiction.—
As we did before with the idealistic perspective, let us now discover 22 the
inner spirit and character of the realistic perspective just laid out. Obviously,
this whole perspective takes its departure from the maxim: [170–171] do not
reflect on the factical self-givenness of our thinking and insight, or on how this
occurs in mind, rather reckon only the content of this insight as valid. Thus, in
other words: do not pay attention to the external form of thought’s existence in
ourselves, but only the inner form of that thought. We posit an absolute truth,
which manifests itself as the content of thinking, and it alone can be true. As
before, it happens for us as we presupposed and wished; since the inner con-
tent alone should be valid, so in fact it alone really matters, and it negates what
it does not contain. Made genetic by us, it was just so. So much in general. 23
Now allow us to delineate realism’s presupposition of an inner absolute
truth more precisely. I believe that there is no way to give a better description,
such as is indeed needed, than this: this implicit truth appears as an image
{Bild}, living, completely determined, and immutable, which holds and bears
itself in this immutability. Now this implicit truth reveals itself in absolute life, 24
and it is immediately evident {einzusehen} that it can only reveal itself in the
latter. Because life is just as truth is: the self-grounded, held and sustained by
itself. Truth is, therefore, in and through itself only life’s image, and likewise
only an image of life gives truth, just as we have described it. By means of the
truth, as grounded by itself, only the image is added. So we stand at about the
same point as before, 25 between life itself and the image of life; as regards this,
we saw that they are fully identical in terms of content, which alone matters
in realism, and are different only in form, which realism leaves to one side.
Living
—————————
Concept Image Being 26
Now it is noteworthy that the image—which holds and sustains itself—
should exist only in the truth as truth, when the former seems, according to92
The Science of Knowing
its character, to be exactly the same as thought and the latter the same as
being; and that therefore in realism (and working from it, if we simply com-
pel it to clarify its fundamental assumption) we are led to a perspective
{Ansicht}, which is so similar to idealism that it could even be the same.
Without venturing further here into this last hint, which meanwhile
could be tossed in to direct our attention to what follows, I will simply con-
clude 27 today’s lecture and pull it together into a whole by means of the fol-
lowing observation. Both idealism and realism grounded themselves on
assumptions. Both of these assumptions—that is in their facticity and in the
circumstance that one actually arises here and the other there—grounded
themselves on an inner maxim [172–173] of the thinking subject. Hence, both
rest on an empirical root. This is less remarkable in the case of idealism, which
asserts facticity, than of realism, which, in its effects and contents, denies and
contradicts what it itself fundamentally is. As we have seen, both are equally
possible, and, if only one grants their premise, they are equally consistent in
their development. Each contradicts the other in the same way: absolute ide-
alism denies the possibility of realism, and realism denies the possibility of
being’s conceivability and derivability. It is clear that this conflict, as a conflict
in maxims, can be alleviated only by setting out a law of maxims, and that
therefore we need to search out such a law.
We can get a rough idea in advance about how a settlement to the con-
flict will work out. All the expressions of the science of knowing so far show
a predilection for the realistic perspective. The justifiability of this preference
follows from this, among other things: that idealism renders impossible even
the being of its opposite, and thus it is decidedly one-sided. On the other hand,
realism at least leaves the being of its opposite undisputed. It only makes it
into an inconceivable being, and thereby brings into the light of day its inade-
quacy as the principle for a science in which everything must be conceived
genetically. Perhaps a simple misunderstanding underlies the proof, given ear-
lier in the name of realism, that an expressive intuition of absolute life can in
no wise arise. In that case, what is proven and needs to be asserted is only that
such an intuition, as valid for itself and self-supporting, can never arise. This
assertion very conveniently leaves room for an interpolation: this intuition
might well arise, and must arise under certain conditions, simply as a phe-
nomenon not grounded in itself. Insight into this interpolation could thus
provide the standpoint for the science of knowing and the true unification of
idealism and realism;—so that the very intuition, purely as such, which we
previously called “our selves at root,” would be the first appearance and the
ground of all other appearances; and because this would not be any error but
instead genuine truth, it and all its modifications, which must also be intuited
as necessary, would be valid as appearances. On the other side, however, seem-
ing and error enter where appearance is taken for being itself. This seemingEleventh Lecture
93
and [174–175] error arise necessarily from [truth’s] absence, and hence can
themselves be derived as necessary in their basis and form, from the assump-
tion that this absence itself is necessary. Some have either discovered or
thought they discovered, I know not which, that in measuring the brow they
could measure people’s mental capacity on the basis of their skulls. 28 The sci-
ence of knowing could easily claim to possess a similar measure of inner men-
tal capacity, if only it could be applied. In every case the rule is this: tell me
exactly what you do not know and do not understand, and I will list with total
precision all errors and illusions in which you believe, and it will prove correct.C H A P T E R
12
Twelfth Lecture
Monday, May 7, 1804
Honored Guests:
In the last discussion period, 1 it was obvious from those who attended and
allowed themselves to converse about things not only that you have followed
me well even in the most recent deep investigations, but that, as is ever so
much more important, a comprehensive vision of the inner spirit and outer
method of the science we here pursue has grown in you. Consequently, I
assume that this is even truer with the rest of you who did not speak; and I
[will] abstract from everything that does not arise for me on this path, with
no misgivings about carrying the investigation forward in the strength and
depth with which we have begun.
A brief review in four parts: 1. production of an insight, which may have
many genetic aspects in its content, but which at its root can only be factical,
since otherwise we would not have been able to go higher. If there really is to
be a “through,” then—as a condition of its possibility—we must presuppose
an inner life, independent in itself from the “through,” and resting on itself.
2. We then made this insight, produced within ourselves, into an object, in
order to analyze it and consider it in its form. There we proposed initially (in
the whole of the second part) that we saw into our concept of an actual
“through,” 2 which appears freely created; or rather, since everything depends
on this, [we proposed] that this concept was energetic and living, that the
inward life of this concept was the principle of the energetic insight into a life
beyond, which grasped us, and which was intuited in this insight as self-suffi-
cient. Thus, it 3 was the principle of intuition and of life in the intuition. This
latter need not arise except in intuition, and its characterization as life in and
for itself is not intrinsically valid. Instead, it can be fully explained from the
mere form of intuition as projecting something self-sufficient in the form of
external existence. In case another perspective is also possible, and since it
This lecture begins at GA II, 8, pp. 176–177. It may have possible taken place on Tues-
day, May 8, 1804.
94Twelfth Lecture
95
begins with the energy of reflection and makes it a principle, this way of look-
ing at the insight can conveniently be labelled the idealistic perspective, in the
terminology already provisionally adopted and explained.
3. But this other view, posited as the basis for the insight, also proved to
be possible, and was realized as well. The presupposed life in itself [178–179]
should be entirely and unconditionally in-itself ; 4 and it is intuited as such.
Therefore, all being and life originates with it, and apart from it there can be
nothing. The reported subjective condition of this perspective and insight was
this: that one not stubbornly hold on to the principle of idealism, the energy
of reflection, but rather yield patiently to this opposite insight. The realistic
perspective. 5
With this a warning!—not as if I detected traces of this misunder-
standing in someone speaking about the matter, but rather because falling into
this error is very easy, as nearly all the philosophical public has done in regard
to the published science of knowing. 6 Do not think of “idealism” and “realism”
here as artificial philosophical systems which the science of knowing wants to
oppose: having arrived in the circle of science, we have nothing more to do
with the criticism of systems. Instead, it is the natural idealism and realism
that arise without any conscious effort on our part in common knowing, at
least in its derived expressions and appearances: and notwithstanding [the
fact] that both can certainly be understood {eingesehen} in this depth and so
on the basis of their principles, they nevertheless arise only in philosophy and
especially in the science of knowing. It is still the latter’s intention to derive
them as wholly natural disjunctions and partialities of common knowing, aris-
ing from themselves.
4. Both these perspectives were more closely specified in their inner
nature and character. Thus, just as at the start, we elevated ourselves above 7
both (we are not enclosed within them, since we can move from one to the
other), [and we moved] from their facticity to the genesis of both, out of their
relative and mutual principles. Hence, the insight which in this fourth part we
lived and were, was their genesis, just as they were the genesis for the previ-
ously created [terms], in which both came together. Thus, according to the
basic law of our science, we are constantly rising to a higher genesis until we
finally lose ourselves completely in it. 8
We characterized them in this way: through its mere being, the idealis-
tic mode of thinking locates itself in the standpoint of reflection, makes this
standpoint absolute by itself, and its further development is nothing more than
the genesis of that which it already was without any genesis other than its own
absolute origin. In its root, therefore, it is factical, not just in relation to some-
thing [180–181] outside itself (as is, e.g., Kant's highest principle 9 ) but rather
in relation to itself. It just posits itself unconditionally, and everything else fol-
lows of itself; and it frees itself from any further accounting for its absolute96
The Science of Knowing
positing. The realistic mode of thinking proceeds no differently. Abstracting
entirely from [the facticity of ] 10 its thinking, it presupposes the bare content
of its thought as solely valid and unconditionally true, and completely consis-
tently it denies any other truth which is not contained in this content, or, as
would actually be the case here, which contradicts it. But this residing in the
content is itself an absolutely given fact, 11 which makes itself absolute without
wanting to give any further account of itself, just like idealism. Therefore, both
are at their root factical ; 12 and even ignoring that they are presented one-sid-
edly, each annulling the other, in this facticity each bears in itself the mark of
its insufficiency as a highest 13 principle for the science of knowing. 14
Let me describe it again with this formula: at this highest point of con-
tradiction between the two terms absolutely demanding unification, we find
these: 0 and C, 15 or the form and the content, or the forms of outer and inner
existence, or [as] in the previous lectures, essence and existence. We appear to
have obtained the absolute disjunction; its unification promises to bring with
it absolute oneness and so to resolve our task fundamentally.
Today we will present considerations regarding this solution which still
remain preliminary 16 —preliminary because we must advance even further just
to get to the point:—in order to prepare ourselves soundly for the highest
oneness.
First, it must be clear that the problem cannot be resolved simply by
combining, rearranging, etc., what we know so far. In relation to our next aim,
everything up to now is only preparation and strengthening our spirit for the
highest insight; and if the preceding should have some further significance
beyond this, this significance can arise only by deduction from the highest
principle. Now we must bring up something entirely new, i.e., according to the
insight adduced just above, it is certain that something has remained to some
extent empirical and concrete for us. We must investigate this and master it
genetically. The rule, therefore, is 17 to investigate this facticity. We have
demonstrated the factical principles of the perspectives in which we recently
spent ourselves, one after the other (and which therefore undoubtedly contain
the highest [form of being] that we, the scientists of knowing, have ourselves
so far attained). One of the two must be developed. Which shall it be?
[182–183] If one grasps hold of it first, the principle of idealism is
admitted to be absolutely incontrovertible. Realism attacks this immediately
as idealistic stubbornness and a false maxim, which it repudiates. Thus,
[realism] denies the principle and cannot reason with idealism in any way.
On the other hand, idealism in turn makes even the beginnings of realism
impossible; it ignores realism completely, and hence can not have anything
against it, since, for idealism, [realism] does not exist. Now realism obviously
takes itself to be superior just because of its denial of idealism’s principle and
by its origin from this denial; thus in realism at least a negative relationshipTwelfth Lecture
97
to idealism remains, whereas even the possibility of such relatedness is extir-
pated in idealism. We must therefore attend to realism, temporarily
abstracting in every way from idealism; since, as said before, we cannot let
realism be absolutely valid, but rather want to correct it, and since we can-
not combat it from the perspective of idealism, we must fight against it on
its own grounds: catch it in self-contradiction. Through this very contradic-
tion, which indeed brings a disjunction with it, realism’s empirical principle
would become genetic, and in this genesis, perhaps it will become the prin-
ciple of a higher realism and idealism [united] into one. We have solved the
first task, finding which factical principle to develop. This is where we
grasp 18 realism in its strength. Its crucial point was life’s in-itself {Ansich} and
within-itself {Insich}. For now, we hold to this feature alone, and in the mean
time we could let go of life. From this “in-itself ” realism infers the negation
of everything outside it.
But how does it bring about this very in-itself? Let’s reconstruct the
process, thinking the in-itself energetically. I assert, and I invite you to con-
sider this yourselves, and see {einzusehen} it immediately as true—I assert: the
in-itself has meaning only to the extent that it [is] not what has been con-
structed, and completely denies everything constructed, all construction, and
all constructability. Consider well: when you say, “thus it is in itself, uncondi-
tionally in itself,” then you are saying [that] it exists thus entirely indepen-
dently from my asserting and thinking, and from all asserting, thinking, and
intuiting, and from whatever other things outside the in-itself may have a
name. This, you say, is how the in-itself must explain itself, if it wishes to
explain itself. No other explanation yields the in-itself. Result: the in-itself is
to be described purely as what negates thinking. 19
[This is] the first surprising observation: here for the first time real-
ism—the perspective which, during the last lecture, we made evident only
factically 20 on the basis of its consequences, [184–185]—is understood
genetically. Previously, that is, this insight arose and grasped us, that if this
life in itself is posited, nothing apart from it can exist. That is how it was;
we saw into it this way and could not do otherwise. Now we see {einsehen}
that realism, or we ourselves standing in its perspective, act like the in-
itself, 21 which negates everything outside itself. [We see] that realism there-
fore to some extent (at least in its effects) is itself the in-itself, and collapses
into it; and for this implicit reason, in the appearance of our insight, which
grasped us in the previous lecture, it annuls everything outside itself. There-
fore, we have comprehended genetically something about realism which
previously was only factical.
With this out of the way on one side, let’s reflect more closely about our
own insight, evoked before, and its principle. I call on you to think the in-itself
and its meaning exactly and energetically, whereupon you would then see into,98
The Science of Knowing
etc. You admit that without this exact thinking you would not have seen it;
perhaps you even admit that for your whole life you have thought the in-itself,
faintly to be sure, and yet this insight has not been produced in you. (It can be
shown that things have gone this way for all philosophy without exception:
since if this insight opened itself really vividly for anyone, then the discovery
of the science of knowing would not have taken so long!) Thus, your insight
into the negation of thinking in itself presupposes positive thought, and the
proposition is as follows: “In thought, thinking annuls itself in the face of the
in-itself.”
Now, to add even more consequences, with which I only wish to
make you acquainted: the negation of thinking over against the in-itself 22
is not thought in free reflection, as the in-itself ought 23 to be thought by us,
rather it is immediately evident. This is what we called intuition {Intuition},
and without doubt, since the absolute in itself is found here, this is the
absolute intuition. What the absolute intuition projects would therefore be
negation, absolute pure nothing—obviously in opposition to the absolute in
itself. And thus idealism, which posits an absolute intuition of life, is
refuted at its root by a still deeper founding of realism. It may well come up
again as an appearance; but, taken as absolute in the way it gave itself out
to be before, it is merely illusion. Hence, we do not get past {es bliebt daher
bei} the previously mentioned fundamental negation of ourselves over
against the absolute.
The negation is intuited ; the in-itself is thought. I ask how and in what
way it is thought; and I explain this implicitly obscure question by [186–187]
the answer itself. That is, we constructed this in-itself, assembling it from
parts, just as, for example, at the beginning of our enterprise we constructed
oneness in the background as not being empirically manifest identity, and not
multiplicity either, but rather as the union of these two. I ought not believe,
instead we posit it directly, together with its meaning, in pure simplicity as the
genuine construction: thinking’s negation is directly evident to us, it grasps us
as proceeding from it in its simplicity. We therefore—this is very important—
didn’t actually construct it, instead it constructed itself by means of itself.
Intuition, the absolute springing forth of light and insight, was bound
directly together with this construction. 24 We, 25 however, would certainly not
have produced this, since it obviously produces itself and draws us forward
with it.—Thus, the absolute’s absolute self-construction and the original light
are completely and entirely one and inseparable, and light arises from this self-
construction just as this self-construction comes from absolute light. Hence,
nothing at all remains here of a pregiven us: 26 and this is the higher 27 realistic
perspective. But now we still hold on to the requirement, as with right we can,
that we should be able to think the in-itself, and think it energetically, 28 thus
that the living self-construction of the in-itself within the light must haveTwelfth Lecture
99
yielded to us, and that once again this energy must be the first condition of
everything, which results in an idealism, 29 which lies even higher.
But on this subject there are once again two things to consider: first, we
are also aware of this thinking, of this energy, and apparently our claim that it
exists (without which general existence it could surely not be a principle) is
grounded solely on this awareness. This, however, presupposes the light. But
since the light itself, at least in this its objective form, does not exist in itself
apart from the absolute (as indeed it cannot, since nothing exists apart from
the absolute) but rather has its source in the in-itself, we cannot appeal to it:
something which itself bears witness against us, if we examine it more closely.
If one retains this higher presupposition of the light in all possible deliver-
ances of self-consciousness (as the source of all idealistic assertions), then the
constant spirit of idealism in its highest form and the fundamental error
which contradicts and destroys it fundamentally would be that it [188–189]
remains fixed on facticity, at the objectifying 30 light from which one can never
begin factically, but only intellectually.
So then (which is surely the same point from another angle), one must
reflect in opposition to the idealistic objection: You are not thinking the in-
itself, constructing it originally—you are not thinking it out ; indeed how could
you! Nor is it known to you through something else, which is not itself the in-
itself, rather it is merely known by you: thus your knowing in and of itself sets
it down; or, as the matter may be more accurately put, it sets itself down in
your knowing and as 31 your knowing. You have been doing this your whole life
without your will and without the least effort and in various forms, as often as
you expressed the judgment: so and so exists. And indeed philosophy has made
war with you and bound you in its circle, not because of this procedure itself
but because of the thoughtlessness of it. You will not give any credit to your
freedom and energy regarding the event itself. Only now that you are aware of
this action and its significance does your energy add anything; likewise again
with the declaration of that which gives itself to you without any effort: intu-
ition. Therefore, before we listen to you at all, we must inquire more exactly
how far the testimony of intuition is valid. 32
This, additionally, by way of conclusion. Whether it appears in its one-
ness as the concept of some philosophical system, either killed or never hav-
ing lived, as it was for us before our realization of its significance, or in a par-
ticular determination as the “is” of a particular thing, the faded in-itself is
always [an object of ] intuition and is therefore dead. For us it exists {ist} in the
concept, and is therefore living. Hence for us there is nothing in intuition, since
everything is in the concept. This is the most decisive distinguishing ground
between the science of knowing and every other possible standpoint for
knowing. It grasps the in-itself conceptually: every other mode of thought
does not conceive it, but only intuits it, and in that way kills it to some extent.100
The Science of Knowing
The science of knowing grasps each of these modes of thought from its own
perspective as negations of the in-itself. Not as absolute negations, but as pri-
vative ones. The things which, as we ascend, we find to be not absolutely valid
(for instance, one-sided realism and idealism), or which might be found to be
so, the science of knowing will take up again on its descent as similar possible
negations of the absolute insight.C H A P T E R
13
Thirteenth Lecture
Wednesday, May 9, 1804
Honored Guests:
Again today and even beyond I will climb on freely. I say “freely” for you,
because I cannot provide the foundation for the distinctions that will emerge
here before using them, rather I must first acquaint you with them in use, even
though a firm rule of ascent may well stand at the basis of what I am doing.
If one just grasps my lecture precisely in other respects, there is no danger of
confusion despite the former circumstance, because instead of the initially
given crucial points L (life) and C (concept) we have the two perspectives:
realism (= genesis of life) and idealism 1 (= genesis of the concept). This should
be known from Friday’s clear presentation and its review the day before yes-
terday. 2 (Not generally, [but rather] in relation to the point, as it was grasped
there, one had to adhere to life in itself, which was required to animate a
“through.” We will not move very far away from this now, and it will easily be
possible to reproduce from it everything that needs to be presented, or to trace
the latter back to it.) 3 In a word, these two perspectives are our present guide,
until we find their principle of oneness and can dispense with them directly.
Here, if anywhere, we need the capacity to hold tight to what is presented
firmly and immovably, and to separate from everything that may very well be
rationally {in der Vernunft} bound to it. Otherwise, one will leap ahead, antic-
ipate the inquiry, and will not grasp the genetic process linking what is taken
up first and its higher terms, which is what really matters; instead the two will
flow into one another factically. 4 “One 5 will anticipate,” I said; but it is not
actually “one,” 6 not the self to whom this happens, rather it is speculative rea-
son, running along automatically. (Then let me add this remark in passing:
speculation, once aroused and brought into play, as I partly know it has truly
been brought into play in you, is as active and vital as the empirical associa-
tion of ideas ever can be, because it is surrounded by a freer, lighter atmos-
phere: and once one has entered this world, one must be just as watchful
This lecture begins at GA II, 8, pp. 190–191.
101102
The Science of Knowing
against leaps of speculation as previously against the stubbornness of empiri-
cism. I want especially to warn those for whom the objects of the present
investigation appear very easily about this danger; I advise you to make them
a little harder for yourself, since this appearance of ease may well arouse the
suspicion that the subject can be grasped more easily by speculative fantasy
than by pure, ever serene, reason.)
[192–193] To work.
The in-itself reveals itself immediately as entirely independent from
knowledge or thought 7 of itself, and therefore as wholly denying the latter in
its own essential effects, in case one assigns such to it. We did not construct
the in-itself in this immediately true and clear concept, instead, as was imme-
diately evident to us, it constructed itself just as it was in the constructive
process, as denying thinking;—immediate insight, the absolute light, was
immediately united with this concept and was evident in the same way. Thus,
the absolute in-itself revealed itself as the source of the light; hence the light
was revealed as in no way primordial:—which is now the first point and which
obviously bears in itself the stamp of a higher realism. 8
Another idealism now tried to raise itself against this realism, proceed-
ing from this basis: 9 since we saw into the in-itself as negating vision, we must
ourselves have reflected energetically on it. Thus, although we cannot deny
that it constructed itself, and the light as well, all of this was nevertheless qual-
ified by own vigorous reflection, which therefore was the highest term of all.
As basing itself on absolute reflection, this is obviously idealism; and, since it
does not depend, as did the previous idealism, on reflection about something
conditioned (actually carrying out a “through”) as a means to realize the con-
dition—but instead 10 rests on reflection on the unconditioned in-itself—it is a
higher idealism.
We quickly struck down this idealism with the following observation.—
“You then,” we address it in personified form—“You think the in-itself: that is
your principle. But on what basis do you know this? 11 You cannot answer oth-
erwise, or bring up a different answer than this: ‘I just see it, am immediately
aware of it,’ and to be sure you do see it, unconditionally objective and intuit-
ing.”—(The last point is important and I will analyze it more closely. In real-
ism too we simply had an insight into the self-construction of the in-itself; but
we had an in-sight, i.e., we saw into 12 something living in itself, and this liv-
ing thing swept insight along with itself, the very same relationship which we
have already frequently found in every manifestness presented genetically.
Now, however, ignoring [the fact] that a purely objectifying intuition 13 seems
to hover above the origin, [194–195] still this intuition is drawn at once
toward, and along with, the genesis. In this insight, therefore, a unification of
the forms of outer and inner existence, of facticity and genetic development
seems to be hinted at. Things stand completely differently with the seeing ofThirteenth Lecture
103
its thought to which idealism appeals. That is, in that case we will certainly
not wish to report that we witness thinking as thinking (i.e., as producing the
in-itself ) in the act of producing, as we certainly actually and in fact witness
the in-itself producing its construction; rather intuition accommodates itself
immediately only to a thinking which is essentially opaque and can be pre-
sented only factically. Thus, it remains completely ambiguous whether think-
ing originates from this intuition, the intuition comes from the thinking, or
whether both might be only appearances of a deeper hidden oneness which
grounds them. 14
In case it is necessary to make this clearer. 15 Could you ever think really
clearly and energetically, which is what we are discussing here (because faded
thoughts and dreaming are completely to be ignored), 16 without being aware
of it; and conversely could you possibly be aware of such thinking without
assuming that you really and in fact were thinking? 17 Would the least doubt
remain for you about the truth of this testimony of your consciousness? I
think not. It is therefore admittedly clear, and immediately proved by the
facts that you cannot distinguish genuine thinking from consciousness of it,
and vice versa; and that in this facticity, thinking posits its intuition, and the
intuition posits the absolute truth and validity of its testimony. 18 We do not
quarrel with you about that. But you cannot provide the genetic middle term
for these two disjunctive terms. Hence, you remain stuck in a facticity. But,
on the other hand, the genesis which has arisen in the opposite, realistic per-
spective opposes you; that is to say, in that perspective we know nothing
about one term of your synthesis, your so-called thinking. But we recognize
that to which you appeal for verification of the latter. Although [we do not
do so] immediately, we still do in its principle. I say “not immediately”; just as
little do we recognize there a simple consciousness, expressing a fact
absolutely, as yours does according to our closer analysis. But I also said “in
its principle”: in any event 19 your consciousness presupposes light and is only
one of its determinations. But light has been realized as itself originating
from the in-itself and its absolute self-construction; however, if it originates
from the in-itself, then this latter cannot likewise originate from it, as you
wish. In your report that you are thinking actually and in fact because you are
conscious of yourself doing so, 20 you must posit your consciousness as
absolute, but the very source of this [196–197] consciousness, pure light, is
not looked at {angesehen} factically, which would bring us to the same level
you occupy; rather (which is more significant) it is seen into {eingesehen}
genetically as itself not absolute. And so this new idealism has been in part
determined further; it does not, as first appears, even posit as absolute a
reflection which, according to it, belongs simply to thinking, instead, it posits
the immediate intuition of this reflection as absolute and is therefore differ-
ent in kind from the first: in part it is refuted as in truth valid, although as an104
The Science of Knowing
appearance it is not yet derived. Assuming that it has become entirely clear
to you, hold onto this and let go of what is deeper.
By the way: in passing and while doing something else {aliud agendo}, I
have touched here on the very important distinction between a merely factical
regarding {faktischer Ansicht}, like our thinking of the in-itself, and genetic
insight {genetischer Ansicht}, like that into the in-itself ’s self-construction. By
means of the immediate testimony of our consciousness, we cannot witness
our thinking, as thinking, literally as production; we see it only so long as it
exists, or should exist, and it already is, or should be, while we see it; on the
other hand, we see the in-itself as existing and as self-constructing, simultane-
ously and reciprocally. 21 This point will have to come up again naturally as the
higher point of disjunction for a still higher oneness, and it will be very sig-
nificant. Meanwhile, let it be impressed [on you], and, as an explanation, let
the following historical commentary be added which may have whatever
worth it can for admirers. At the same time it may serve as an external test of
whether one has understood me.
Reinhold (or, as Reinhold claims, Bardili ) 22 wishes to make thinking qua
thinking the principle of being. Therefore, on the most charitable [198–199]
interpretation, his system would be situated within the idealism we have just
described, and one must assume that by this he means the thinking of the in-
itself which we carried out the day before yesterday. Now, first of all, he is very
far removed from explaining that this in-itself negates seeing, as we have
shown; but then, which is worse, in regard to the in-itself ’s real existence (with
which he generally has no dealings and which, to be sure, he would be able to
prove only factically from the existence of individual things) he does not
appeal to consciousness (a fact of which I reminded him, 23 but which escaped
him) because he sees very well that doing so would lead him to an idealism,
and he seems to have developed an unconquerable horror of any idealism.
Hence, in the first place, his principle stands entirely in the air, and he works
to build a realism on absolutely nothing; and he could only be driven to this
by despair, following the rule: “since it wouldn’t work with anything I tried so
far, then it must work with the one thing still left in my field of vision.” Sec-
ond (and I have made this observation especially to make this point), since,
according to an absolute law of reason, thinking does not let itself be seen into
as producing itself as thinking, then naturally Reinhold cannot realize it this
way either, nor can he genetically deduce the least thing from it. Hence he can
only say, like Spinoza: since everything that is lies in it, and since now so-and-
so exists, it too must lie in it. Because he was educated in the Kantian school,
and later by the science of knowing, he may not now do this. Hence he labors
to deduce; but since, if one only has a clear concept, this appears to be com-
pletely impossible, [200–201] total darkness and obscurity arise in his system,
so that nobody grasps what he really wants. If one considers this system fromThirteenth Lecture
105
the point of view of the science of knowing, and indeed from the very point
from which we have just seen it, then the obscurity of its principles is clear.)
Let’s get back to business and draw the conclusion, because in passing
we have once again gained a very clear insight into the true essence of the sci-
ence of knowing, that is, of the principle which we are still obliged to present.
The refuted idealism makes immediate consciousness into the absolute, the
primordial source, the protector of truth; and indeed absolute consciousness
reveals itself in it as the oneness of all other possible consciousness, as reflec-
tion’s self-consciousness. In the first place, then, this stands firm as one of our
basic foundations. Wherever we say “I am conscious of that,” our testimony
bears the basic formal character of an absolute intuition, which we have just
described, and makes a claim for the intrinsic validity of its contents. This
consciousness is now realized as self-consciousness and reflection in its root: all
possible disjunctions and modes of consciousness must be deduced from self-
consciousness; and we would therewith already have achieved a comprehen-
sive study.
It is clear that this consciousness is completely one in itself and capable
of no inner disjunction; because the thinking that arose in it was that of the
in-itself which, as in-itself, is entirely one and self-same; thus it too was one,
and the consciousness of it was only this one consciousness; and therefore it
was also one. The self, or I, which arises here is consequently the absolute I—
pure eternally self identical, and unchanging—but not the absolute as will soon
be more precisely evident. A specific disjunctive principle will also have to be
identified if:
a. in the course of thinking “the one in-itself ” it should appear in a multi-
faceted perspective, even though in the background it is to remain perpet-
ually the same single in-itself, or categorical is, and
b. as a result there also arises a manifold view of thinking, or reflection, and
hence of the reflecting subject, or I (all of which, just like the in-itself, are
also to remain the same one in the background). 24
It could well be, and indeed will turn out, that, if we remain trapped in the
absolute I and never raise ourselves beyond it, we may never discover this dis-
junctive principle in proper genetic fashion, and it will have to be disclosed
factically.
(An historical note in passing: [202–203] even where things go the best
for the science of knowing, it has been taken for this idealistic system we have
just described, which presupposes what we have exactly characterized as the
absolute I to be the absolute and derives everything else from it; and no author
known to me, friend or foe, has risen to a higher conception of it. That most
have remained at an even lower level than this conception goes without saying.106
The Science of Knowing
If a higher understanding is to arise for anyone besides the originator of this
science, it could only be among the present listeners, but it would not be writ-
ten; because what can be understood in writing stands under the previous rule;
or it will be hit upon by you. This remark has the consequence that nobody will
get a report about the essence of the science of knowing from anyone other
than the present one from the originator. It will be immediately evident even
from just the clear and decisive letters of what has been published on the topic
how false this interpretation is.) 25
The intrinsic validity of this idealism is refuted; nevertheless, it may pre-
serve its existence as an appearance, and indeed as the foundation of all
appearance, which we have to expect: it has been refuted on the grounds that
it is factical and that a higher development points to its origin. One calls a fact
{Thatsache} “factical,” {faktisch} and since here we are speaking of conscious-
ness, this fact would be a “fact of consciousness”; or, to put it more strongly:
according to this idealistic system, 26 consciousness itself is a fact, and since
consciousness is for it the absolute, the absolute 27 would be a fact. Now, from
the first moment of its arising the science of knowing has declared that the
primary error 28 of all previous systems has been that they began with some-
thing factical and posited the absolute in this. It, on the other hand, lays as its
ground and has given evidence for, an enactment, 29 which in these lectures I
have called by the Greek term genesis, since these are often more easily under-
stood rightly than German terms. Therefore from its first arising, the science
of knowing has gone beyond the idealism we have described. It has shown this
in another equally unambiguous way: particularly regarding its basic point, the
I. It has never admitted that this [204–205] I—as found and perceived 30 —is its
principle. “As found” it is never the pure I, but rather the individual personal
being of each one, and whoever claims to have found it as pure finds himself
in a psychological illusion of the kind with which we have been charged by
those ignorant of this science’s true principle.—So then, the science of know-
ing has always testified that it recognizes the I as pure only as produced, and
that, as a science, it never places the I at the pinnacle of its deductions, because
the productive process will always stand higher than what is produced. This
production of the I, and with it the whole of consciousness, is now our task.
The idealism which has been rejected as intrinsically valid is the same as
absolute immediate consciousness. Therefore, so that we now express as force-
fully as possible what it comes down to, the science of knowing denies the
validity of immediate consciousness’s testimony absolutely as such and for this
exact reason: that it is this, 31 and it proves this denial. Solely in this way does
the science of knowing bring reason in itself to peace and oneness. Only pure
reason, which is to be grasped merely by intellect, remains as solely valid. So
that no one is confused even for a moment by a fancy which might easily arise
here, I immediately add a hint which must be discussed further in what follows.Thirteenth Lecture
107
Namely, someone may say: “But how can I grasp something in intellect 32 without
being conscious in this intellectualizing?” I answer: “Of course you cannot.” But
the ground of truth as truth does not rest in consciousness, but only and
entirely in truth itself. You must always separate consciousness from truth, as
in no way making a difference to it. Consciousness remains only an outer
appearance of truth, from which you can never escape and whose grounds are
to be given to you. But if you believe that the grounds why truth is truth are
found in this consciousness, you lapse into illusion; and every time something
seems true to you because you are conscious of it, you become at the [206–207]
root idle illusion and error. Here now it is obvious: 1. how in fact the science
of knowing keeps its promise and, as a doctrine of truth and reason, expunges
all facticity from itself. The primordial fact and source of everything factical is
consciousness. It can verify nothing, in light of the science of knowing’s proof
that whenever truth is at issue, one should turn it aside and abstract from it. To
the extent that this science in its second part is a phenomenology, a doctrine of
appearance and illusion, which is possible only out of the first and on the
ground and basis of the latter, to that extent it surely deduces both as existing,
but simply as they indeed exist, as factical. 33 2. It has become completely clear
why nothing external could be brought up against the science of knowing, but
that one always must begin with it to gain entrance into it. The beginning
point for a fight against it is either grasped in intellect or not. If it is grasped in
intellect, then it is either grasped intellectually immediately—and that is the
principle of the science—or it is grasped mediately and these must be either
deductions of the fundamental phenomenon or phenomena derived from it.
One can come to the latter only through the former. In this case, therefore, one
would in every circumstance be at one with the science of knowing, be the sci-
ence of knowing itself, and in no case be in conflict with it. If it were not
grasped intellectually but should nevertheless be true, then one must appeal for
verification to one’s immediate consciousness, since there is no third way to get
to the absolute itself or even to an appearance of it. But making this appeal one
is straightaway turned aside with the instruction that precisely because you are
immediately aware of it and appeal to that fact, it must be false. Of course
thoughtlessness and drivel have created a fancy title for themselves, that of
“skepticism,” and they believe that nothing is so high that it cannot be forced
under this rubric. It must stay away from the science of knowing. In pure rea-
son doubt can no longer arise; the former bears and holds itself—and everyone
who enters this region—firmly and undisturbedly. But if skepticism wishes to
doubt the implicit validity of consciousness—and this is approximately what it
wants in some of its representatives—and it does this provisionally in this or
that corner, although without having been able to bring about a properly basic
general doubt;—[208–209] if it wants this, then with its general doubt it has
arrived too late for the science of knowing, since the latter does not just doubt108
The Science of Knowing
this implicit validity 34 in a provisional way, but rather asserts and proves the
invalidity of what the general doubt only puts in question. Just the possessor of
this science (who surveys all disjunctions in consciousness, disjunctions which,
if one assumes the validity of consciousness in itself, become contradictions)
could present a skepticism which totally negated everything assumed so far; a
skepticism to which those who have been playing with all kinds of skeptical
doubts as a pastime might blanch and cry out: “Now the joke goes too far!” Per-
haps in this way one might even contribute to arousing the presently stagnat-
ing philosophical interests. 35C H A P T E R
14
Fourteenth Lecture
Thursday, May 10, 1804
Honored Guests:
[Yesterday, an idealism, which made absolute consciousness (in its actuality that
is) into its principle was presented, characterized, and refuted;—“in its actual-
ity” I say, since today we will uncover still a different [idealism] in a place where
we do not expect it, one which makes the same thing its principle, only merely
in its possibility. Now this absolute consciousness was self-consciousness in the
energy (of reflective thinking, as it later turned out). On this point I will add
another remark relating to the outer history of the science of knowing, which
naturally is not intended to parade my conflict with this unphilosophical age in
front of you, but rather only to provide hints to those who are following this
science in my published writings and who want to rediscover it in the form
employed there—telling them where they should direct their attention.] 1
[Let me add just this, it is clear that just as the form of outer existence
as such perishes, so its opposite as such perishes as well; therefore, realism, or
more accurately objectivism perishes along with that idealism which, because
of language’s ambiguity, we might better call subjectivism. Reality remains, as
inner being—as we must express ourselves just in order to talk; but in no way
does it remain a term of any relation, since a second term for the relation, and
indeed all relations in general, have been given up. Therefore, it is not “objec-
tive,” since this word has meaning only over against subjectivity, which has no
meaning from our standpoint. Only one recent philosophical writer (I mean
Schelling) has had a suspicion about this truth, with his so-called system of
identity; not, certainly, that he had seen into the absolute negation 2 of subject
and object, but that with his system he aimed at a synthesis post factum; and
with this operation he believed he had gone beyond the range of the science
of knowing. Here is how things stand about that: he had perceived this syn-
thesizing in the science of knowing, which carried it out, and he believed him-
self to be something more when he said what it did. This is the first unlucky
This lecture begins at GA II, 8, pp. 214–215.
109110
The Science of Knowing
blow that befell him: saying, which always stems from subjectivity and by its
nature presents a dead object, is not more but less excellent than doing, which
stands between both in the midpoint of inwardly living being. Moreover, he
does not prove this claim, but lets the science of knowing do it for him, 3 which
again seems odd: that a [216–217] system which admittedly contains the
grounds for proving our own system’s basic principle should be placed below
it. Now he begins and asserts: reason is the absolute indifference between sub-
ject and object. But here must first also be added that it cannot be an absolute
point of indifference without also being an absolute point of differentiation.
That it is neither one nor the other absolutely but both relatively; and that
therefore, however one may begin it, no spark of absoluteness may be brought
into this reason. So then he says: reason exists; in this way he externalizes rea-
son from the start and sets himself apart from it; thus one must congratulate
him that with his definition he has not hit the right reason. This objectifica-
tion of reason is completely the wrong path. The business of philosophy is not
to talk around reason from the outside, but really and in all seriousness to con-
duct rational existence. Nevertheless, this author is the hero of all passionate,
and therefore empty and confused heads; and especially of those who do not
disavow defects like those reproved above, to which, when possible, they come
even more extremely because they think either that the inferences are good,
although the principles are false, but that the whole is still excellent, admit-
tedly overlooking that all the individual parts are good for nothing; or finally
that although it is neither true, nor good, nor beautiful, it still remains very
interesting. For my own person, I have said all this only in the interest of his-
tory and to elucidate my own views, but in no way to weaken anyone’s respect
for their hero or to lead them to myself. Because if anyone wishes to be con-
demned to error, I have nothing against it.] 4
Further, and to the point; [here is the] chief result: 5
Consciousness has been rejected in its intrinsic 6 validity, despite the fact
that we have admitted we cannot escape it. We absolutely, i.e., even here in the
science of knowing. 7 Therefore, 1. if we have once seen into this fact, although
factically we could never negate consciousness, we will not really believe it
when judging truth; instead, when judging, we will abstract from it; indeed,
on the condition that we want to get to truth, we must do this, but not uncon-
ditionally, since it is not necessary that we see into 8 the truth. Here for the first
time, we ourselves have become entwined mediately in science and the circle of
its manifestness, and we became the topic without any effort on our part in a
new way, 9 because we were speaking of consciousness and we ourselves occur
empirically as consciousness, which might serve very well for the genetic
deduction of the I, which we seek. Further, we should cultivate here a maxim
for ourselves, a rule of judging which can be appropriated only 10 through free-
dom; and this maxim should become the absolute principle [218–219] in andFourteenth Lecture
111
for us: if never of truth itself, then of this truth’s factical appearance. In one
respect, this is generally meaningful and may lead to a new idealism, in a
region where it alone can have value as the principle of appearance; and in
another respect it confirms the opinion expressed previously when we were
characterizing and refuting the more deeply placed idealism and realism: that,
as grounded on conflicting maxims, both could only be reunited by means of
a higher maxim. And so notice this also at the same time: that the present
maxim is quite different from that of the previous realism (which allowed
truth alone to be unconditionally valid) in that the present one has a condition
{bedingt ist}: “if truth is to {soll} be valid, then it must, etc.”; at the same time
certainly acknowledging that truth need not necessarily be valid. Finally, free-
dom shows itself here in its most original form, in respect to its actual opera-
tion as we have always described it, not as affirmative, creating truth, but rather
merely as negative, averting illusion. 11 All of which, although very significant,
are only expositions of this 12 insight: if consciousness in itself has no validity
and relation to truth, then we must abstract from all effects of this conscious-
ness in the investigations which lie before us, and whose task is to deliver truth
and the absolute purely to the light of day.
2. From what, then, do we actually need to abstract, and what is its
unavoidable effect? This is evident from that salient point and nerve, for
whose sake consciousness was rejected as insufficient. But, in virtue of our last
investigation, the nerve was this, that it projected something factically, that is
in its highest potency; and in our case [it projected] the energy, which then
would become thinking, whose genetic connection with it it could in no way
give; thus a thing that it projected purely and through an absolute gap {per
absolutum hiatum}. Grasp this character exactly, just as it has been given, and
to that end remember what was said last time: e.g., you would not presume
that you could actually think without being conscious of it, and vice versa; nor
that you could be conscious of your thinking, without in fact actually think-
ing, or that this consciousness was deceiving you; but if you were asked to pro-
vide an explicable and explanatory ground for the connection of these two
terms, you would never be able to provide such a ground. Thus, changing
places with you at the site from which you conducted your proof: your con-
sciousness of thinking should contain a thinking process {Denken}, actual, true
and really present, without you being able to give an accounting of it;
[220–221] therefore, this consciousness projects a true reality outward, dis-
continuously: an absolute inconceivability and inexplicability.
This discontinuous projection is evidently the same one that we have pre-
viously called, and presently call, the form of outer existence, which shows itself
in every categorical is. For what this means, as a projection, concerning which
no further account can be given and which thus is discontinuous, is the same as
what we called “death at the root.” The gap, the rupture of intellectual activity112
The Science of Knowing
in it, 13 is just death’s lair. Now we should not admit the validity of this projec-
tion, or form of outer existence, although we can never free ourselves from it fac-
tically; and we should know that it means nothing; we should know, wherever
it arises, that it is indeed only the result and effect of mere consciousness (ignor-
ing that this consciousness remains hidden in its roots) and therefore not let
ourselves be led astray by it. This is the sense of our discovered maxim, which is
to be ours from now on in every case where we 14 need it. This very is is the orig-
inal appearance: which is closely related to, and may well be the same thing as,
the I which we presented previously as the original appearance.
3. Thus it is decreed against the highest idealism, and this maxim
imposes the highest realism yet on us. Before we go further under its leader-
ship, though, it may be advisable to test it against the law which it itself has
brought forward, thus to draw it directly before its own seat of judgment in
order to discover whether it itself is indeed pure realism.—It proceeds from
the in-itself and proposes this as the absolute. But what is this in-itself as such
and in its own self? You are invited here to a very deep reflection and abstrac-
tion. Although the foregoing brings to an end the thinking 15 of this in-itself,
reflected first by consciousness, although, we have likewise already had to
admit before that we did not construct this in-itself but that it already is found
in advance as completely constructed and finished and comprehensible in
itself, and thus as constructed by itself, so that we in any case have nothing to
do with this; we may still investigate this original construction more closely in
terms of its content.
—I have said that I invite you to a very deep reflection and abstraction.
What this abstraction is—an abstraction which may be described in words as
well as possible—will barely be made clear from what follows; but this will not
harm anything, and it is certain in any case that it is already clear, and I impose
on myself the task [222–223] of grasping the highest in words, and on you the
task of understanding it in a pure form. Thus—once again, as already previ-
ously, the topic is the in-itself, and we are called, here as before, to a consider-
ation of its inner meaning and to its re-construction. We will not complete
again what has been done before, an act which, holding us in a circle, would
not advance us from the spot. But, if we wish [to achieve] something else, how
is this to be distinguished from the preceding? Thus, above we presupposed the
in-itself and considered its meaning while we supplemented it with life, or a
primal fantasy, dissolved ourself in the latter, and had our root in it. 16 Of course
this life is not supposed to be our life, but rather the very life and self-con-
struction of the in-itself: it was then an inner determination (one which arose
immediately in this context) of the original life itself, which still remained
dominant in this case. So it was previously.—Now, however, we elevate our-
selves for the first time to the in-itself that is presupposed in this procedure,
[knowing it] as presupposed and unconditionally immediate, independent ofFourteenth Lecture
113
this living reconstruction, determinate, and comprehensible: and without this
original significance the reconstruction, as a reconstruction and clarification,
would have no basis and no guide. For this reason, I said previously that the
originally completed construction, the enduring content, should be demon-
strated. Therefore, things must proceed in this work so that what is absolutely
presupposed remains presupposed, 17 so that as a result the living quality which
we bring with us will mean nothing at all—neither as our vivacity, nor as vivac-
ity in general—likewise the very validity of primal fantasy, although it cannot
be withheld factically, yet is denied as real, in which denial the true essence of
reason may well consist. (Or, more succinctly, if only one will understand it;
what has been made perceptible resides in this latter construction, and the mean-
ing of the in-itself grasped purely intellectually should be found there. 18 )
So much in the way of preliminary formal description of this new
reconstruction of the in-itself. Now to the solution.—However one may
wish to take up the in-itself, it is still always 19 qualified by the negation of
something opposed to it, thereby as in-itself it is itself something relative,
the oneness of a duality, and vice versa. Certainly, it is genuinely at once a
synthetic and analytic principle, as we have all along looked for: but still it
is no true self-sufficient oneness; since the oneness lets itself be grasped only
through duality: although admittedly duality also lets itself be completely
grasped and explained through oneness. In a word, the in-itself, 20 grasped
more profoundly, is no in-itself, no absolute, because it is not a true oneness,
and even our realism has not [224–225] pushed through to the absolute.
Viewed still more rigorously, 21 in the oneness 22 there is in the background a
projection of in-itself and not-in-itself, which posit one another reciprocally
for explanation and comprehensibility, and which negate one another in
reality; and in return, the oneness is a projection of both terms. Further, this
projection happens completely immediately, through a gap, without being
able to provide the requisite accounting of itself. Because how an in-itself
and a not-in-itself follow from oneness as simple, pure oneness cannot be
explained. Of course, it can be done if the oneness is already assumed to be
the oneness of the in-itself and not-in-itself; but then the inconceivability 23
and inexplicability is in this determinateness of oneness, and it itself is only a
projection through an irrational gap. This determinateness would have no
warrant other than immediate awareness; and actually, if we will think back
to how we have arrived at everything so far, it has no other ground. “Think
an in-itself ” it began, and this thinking, or consciousness was possible. And
this possibility has shaped our entire investigation to date; thus we have sup-
ported ourselves on consciousness, if not quite on its actuality, then certainly
on its possibility, and in this quality we have had it for our principle. Hence,
our highest realism, i.e., the highest standpoint of our own speculation, is
itself revealed here as an idealism, which so far has just remained hidden in114
The Science of Knowing
its roots; it is fundamentally factical and a discontinuous projection, does
not stand up to its own criteria, and, according to the rules it itself estab-
lished, it is to be given up.
4. Why should it be given up? What was the true source of the error
which we discovered in it? Being in-itself [was discovered] 24 as a negation and
a relational term. Hence we must unconditionally let that go, if it, or our entire
system, is to survive. But something is still left over for us. I affirm this and
instruct you to find it with me: being 25 and existence 26 and resting, 27 taken as
absolute, 28 remains—and of course I add: “being and resting on itself,” but I
already clearly knew that the latter would be a mere supplement for clarifica-
tion and [226–227] illustration, but would mean nothing at all in and for itself
and would add no supplement to the completeness and self-sufficiency of
being’s inner essence.—If I wish to look back to the previous, already dis-
carded expression, “being in-itself ” means a being which indeed needs no
other being for its existence. Precisely through this not-needing it becomes
intrinsically more and more real than it was before; and not-needing this not-
needing does not also belong to its absolute not-needing, and so too with the
not-needing of this not-needing of not-needing, so that this supplement in its
endless repeatability remains always the same and always meaning nothing in
relation to the essence taken seriously and inwardly. Thus, I see into [the fact]
that (generally, in its core, thus indeed as the point of oneness, which was pre-
viously tested and discarded) the entire relation and comparison with the not-
in-itself—from which the form of the in-itself as such arises—is completely
null in comparison to the essence. It is without meaning or effect. Because I
see into this, and thereby handle the addition in almost as negative a way as
does the essence itself, then I must, as an insight, participate in the essence in
some manner still to be developed.
Now, to be sure, if I pay attention to myself, I can always become aware
that I objectify and project this pure being: but I already certainly know that
this means nothing, alters nothing about being, and adds nothing to it. To be
sure, in another shape this projection is the in-itself ’s supplement, whose noth-
ingness has already been realized: therefore I will never be deceived by it. In
brief, the entire outer existential form has perished in this shape, since it is the
latter in the highest [element] in which it occurs, the in-itself; we have only the
inner essence left with which to deal, in order to work it through: but we truly
work it through if we see into it as the genesis for its appearance in the outer
existential form; and nothing else can lead us to that except not allowing our-
selves to be deceived by this form. If it does deceive us, then we just are it, dis-
solved and lost in it, and we will never arrive at its origin {Genesis}.
I wished at least to attach this last, fourth, point here, in order not to
end the lecture with death and destruction, as it does on first appearance.
Tomorrow its further development. 29C H A P T E R
15
Fifteenth Lecture
Friday, May 11, 1804
N.B. Which Contains the Basic Proposition
Honored Guests:
My task for today is this: first of all to work out fully and completely the main
point discovered last time; then to present a general review of the new mate-
rial added this week, and thus, as it were, to balance the accounts, since with
this lecture we conclude the week’s 1 work and a discussion period intervenes
{dazwischen fällt}. 2
On the first matter. This much as a reminder in advance. To begin with,
the point I must present is the clearest and simultaneously the most hidden of
all, in the place where there is no clarity. 3 Not much can be said about it, rather
it must be conceived at one stroke; even less can anything be said about it or
words used to assist comprehension, since objectivity, the first basic twist of all
language, 4 has already long been abandoned in our maxim, and is to be
annulled here within absolute insight. At this point, therefore, I can rely only
on the clarity and rapidity of spirit which you have achieved in the previous
investigations. So then, on a particular occasion I divided the science of know-
ing into two main parts; one, that it is a doctrine of reason and truth, and sec-
ond, that it is a doctrine of appearance and illusion, 5 but one that is indeed true
and is grounded in truth. The first part consists of a single insight and is begun
and completed with the single point which I will now present. To work! After
the problem of absolute relation, which appeared in the original in-itself, which
itself pointed to a not-in-itself, 6 nothing remained for us except the pure, bare
being by which, following the maxim, our objectivizing intuition must be
rejected as inadequate.—What then is this pure being in its abstraction from
relatedness? Could we make it even clearer to ourselves and reconstruct it? I say
This lecture begins at GAII, 8, pp. 228–229. The subtitle here occurs only in the Copia.
The GA editors observe that “N.B.” abbreviates either nota bene or NebenBemerkung. I
prefer the former.
115116
The Science of Knowing
yes: the very abstraction imposed on us helps. Being is entirely of itself, in itself,
and through itself ; this self is not to be taken as an antithesis, but grasped with
the requisite abstraction purely inwardly, as it very well can be grasped, and as
I for example am most fervently conscious of grasping it. Therefore, to express
ourselves scholastically, it has been constructed as a being in pure act {esse in
mero actu}, so that both being and living, and living and being completely inter-
penetrate, dissolve into one another, and are the same, and this self-same
inwardness is the one completely unified being, which was the first point.
[230–231] This sole being and life can not exist, or be looked for, out-
side itself, and nothing at all can exist outside it. Briefly and in a word: dual-
ity or multiplicity does not occur at all or under any conditions, only oneness;
because by its own agency being itself carries self-enclosed oneness with itself,
and its essence consists in this. Being—understood by language as a noun—
cannot literally be 7 actively 8 except immediately in living; but it [“being”
understood as a noun] is only verbally; because completely noun-like being is
objectivity, which in no way suffices: and it is only by surrendering this sub-
stantiality and objectivity, not merely in pretense but in the fact and truth of
insight, that one arrives at reason.
On the other side, 9 what lives immediately is the “esse,” 10 since only the
“to be” lives, and then it, as an indivisible oneness, which cannot exist outside
itself and cannot go out of itself into duality, is something to which the things
we have demonstrated immediately apply.
But we live 11 immediately in the act of living itself, therefore we are the
one undivided being itself, in itself, of itself, through itself, which can never go
outside itself to duality.
“We,” I say— 12 to be sure, we are immediately conscious that, insofar as
we speak of it, we again objectify this “We” itself with its inward life: but we
already know that this objectification means as little in this case as in any
other. And surely we know that we are not talking about this We-in-itself,
separated by an irrational discontinuity from the other We which ought to be
conscious; 13 rather [we are talking] purely about the one We-in-itself, 14 living
purely in itself, which we conceive merely through our own energetic negation
of the conceiving which obtrudes on us empirically here.—This We, in imme-
diate 15 living itself; this We, 16 not qualified or characterizable by anything that
might occur to someone, but rather characterizable purely by immediately
actual life itself. 17
This was the surprising insight to which I wished to elevate you, in
which reason and truth emerge purely. If anyone should need it, I will point
to it briefly from yet another aspect. If being is occupied with its own absolute
living, and can never emerge out from it, then it is a self-enclosed I, and can
be nothing else besides this; and likewise a self-enclosed I is being: 18 which “I”
we could now [232–233] call “We” in anticipation of a division in it. Hence,Fifteenth Lecture
117
we in no way 19 depend here on an empirical perception of our life, which would
need to be completely rejected as a modification of consciousness; rather we
are depending on the genetic insight into life and the I, which emerges from
the construction of the one being, and vice versa. We already know, and
abstract completely from, the fact, that this very insight as such, together with
its reversal, is irrelevant and vanishes; we will need to look back to it again only
in deducing phenomena. 20
As it has been presented now, this intrinsically cannot be made clearer
by anything else; since it itself is the original source and ground of all other
clarity. Still, the subjective eye can get clearer, and become more fit for this
clarity, through deeper explanation of the immediately surrounding terms;
therefore I add in this regard another consideration which lies on the system’s
path anyway. Yesterday, and again today too at the beginning 21 of our medita-
tion, although we constructed being according to its inner essence, if we only
remember, we placed being objectively before ourselves, despite the fact that,
though only as a result of following our maxim, we grant no validity to this
objectivity:—factically, 22 even though surely not intelligibly or in reason, being
remains separated from itself. But, just as in this reflective process we are
grasped by the insight that being itself is an absolute I, or we; thus the first
remaining disjunction—between being and the we—is completely annulled,
even in facticity, and the first version of the form of existence is also factically
negated. Previously [we knew that] at the very least we emerged factically out
of ourselves toward being, and [in this process] it could very well happen that
being would not come out of itself, especially if we did not wish to be being.
We did not accept this being as valid, merely on the basis of a maxim that had
its proof via derived terms, and thus might well need a new proof here. As we
become being itself in the insight we have produced, we can—as a result of
this insight—no longer come out of ourselves toward being, since we are it;
and really we absolutely cannot come out of ourselves at all, because being
cannot come out of itself. Here the preceding maxim has received its proof,
[234–235] its law, and its immediate realization in the insight: because this
insight in fact no longer objectifies being. Now to be sure, I say, no other
objectifying consciousness arises together with this insight— 23 because in
order for that to occur, a self-reflecting would be required to stand in
between—however the possibility does arise of an objectifying consciousness:
[namely,] our own. 24 Now, as regards the content of this new objectification,
it is already clear that it does not bring with itself any disjunction in our sub-
ject matter, as does the first—between real being 25 and absolute non-being—
rather this content brings only the mere repetition and repeated supposition
of one and the same I, or We, which is entirely self-enclosed, which encom-
passes all reality in itself, and which is therefore entirely unalterable; therefore,
it does not contradict the original law of not going out of oneself in essence.118
The Science of Knowing
But, as the first stage of our descent into phenomenology, we will have to
explore whence this empty repetition and doubling may arise; today it is
merely a matter of establishing the insight which is expressive of pure rea-
son—that being or the absolute is a self-enclosed I—in its unalterability.
Now on to the second part of the general review. As must arise in any
presentation of the science of knowing, we have proceeded factically, inwardly
completing something useful for our purpose, and paying attention to how we
have done it, compelled always, as is evident, by an unconscious law of reason
working in us. On this path, which I will not repeat now, we had elevated our-
selves to a pure “through,” 26 as the essence of the concept; and had understood
{eingesehen} that the latter’s realization presupposed a self-subsisting being.
Having presented this insight as a fact and reflected further about its princi-
ple, it turned out either that one could posit the energy of thinking a
“through” needing completion as the absolute and thus as the source of intu-
ition and of life in itself within intuition, which was an idealism; or that, con-
sidering that life ought to exist in itself, one could take the latter as the prin-
ciple, with the result that everything else perishes; this latter [was] a realism.
Both were supported by maxims. The former on this: the fact of reflection is
to be taken as valid, and nothing else; the latter on this: the content of the evi-
dent proposition is to be taken as valid, and nothing else; and, for that very
reason, both are at bottom factical, since indeed even the contents of what is
manifest, which alone should be valid for realism, is only a fact.
Along with the necessity that arose from this to ascend higher and to
master the facts genetically, we turned our attention to what promised to be
most significant here, [namely] to the in-itself, bound to the realistic princi-
ple, [236–237] to life in itself: and this further deliberation was the first step
we made this week. 27 It turned out that the in-itself manifests as an absolute
negation of the validity of all seeing directed toward itself: 28 that it constructs
itself in immediate manifestness, and with its own self-construction even gives
off immediate manifestness or light: yielding a higher realism which deduces
insight and the light themselves, items that the first realism was content to
ignore. A new idealism attempts to establish itself against this new realism.
That is, we had to take command of ourselves and struggle energetically to
contemplate the in-itself in its significance. So, [we] believed we realized that
this in-itself first appeared as a result of this reflection as simultaneously con-
structing itself with immediate manifestness in the light; and that conse-
quently this energy of ours would be the basic principle and first link in the
whole matter.
Realism, or we ourselves, since we are nothing else than this realism,
fights very boldly against this as follows: “If you really actually think . . . ,” and
to what will you appeal for confirmation of this assertion: you can adduce
nothing more than that you are aware of yourself, but you cannot derive think-Fifteenth Lecture
119
ing genetically in its reality and truthfulness, as you should, from your con-
sciousness in which you report it; but, by contrast, we can derive the very con-
sciousness to which you appeal, and which you make your principle, geneti-
cally, since this can surely be only a modification of insight and light, but light
proceeds directly out of the in-itself, manifestness in unmediated manifest-
ness. The higher maxim presented in this reasoning would just be this: to give
no credence to the assertions of simple, immediate consciousness, even if one
cannot factically free oneself from them, but rather to abstract from them.
What is this consciousness’s effect, for the sake of which it 29 is discarded; and
therefore what is that which must always be removed from the truth? Answer:
the absolute projection of an object whose origin is inexplicable, so that
between the projective act and the projected object everything is dark and
bare; as I think I can express very accurately, if a little scholastically, a proiec-
tio per hiatum irrationale (projection through an irrational gap).
Let me again draw your attention to this point, both for now and for
all your future studies and opinions in philosophy: if [238–239] my current
presentation of the science of knowing has been clearer than all my previous
versions of the same science, and can maintain itself in this clarity, and if
clear understanding of the system is to make a new advance by these means,
the ground for this must lie simply in the impartial establishment of the
maxim that immediate consciousness is in no way sufficient and that hence
it does not suffice in its basic law of projection per hiatum. Of course the
essence of this truth has ruled in every possible presentation of the science
of knowing from the very first hint which I gave in a “Review of Aeneside-
mus” in the Allgemeinen Literatur Zeitung; because this maxim is identical
with the principle of absolute genesis. If nothing that has not been realized
genetically is permitted, then projection per hiatum will not be permitted,
since its essence consists precisely in non-genesis. If one has not made him-
self explicitly aware that this non-genesis, which is to be restrained in think-
ing, 30 remains a factical element of the consciousness which is unavoidable
on every path of all our investigations and of the science of knowing itself,
then one torments and exhausts himself trying to eliminate this illusion, as
if that were possible. And the sole remaining way of breaking through to
truth is to divide the illusion, and intellectually to destroy each part one at
a time, while during this procedure one actually defers the illusion to
another piece at which annihilation will arrive later, when the first piece
could once again serve as the bearer of the illusion. This was the science of
knowing’s previous path—and it is clear that it too leads to the goal,
although with greater difficulty. However, if one knows the origin of non-
genesis in advance, and that it always comes to nothing, although it is
unavoidable, then one no longer fights against it, but rather allows it to work
peacefully: one simply ignores it and abstains from its results. It is possible120
The Science of Knowing
in this way alone to gain access to insight immediately as we have done deci-
sively, and not just by inference from the nonexistence of the two halves.
Let me continue the recapitulation: the realism which presented the
recently analyzed maxim was itself questionable and was brought before the
judgment seat [240–241] of its own maxims. There, on closer consideration,
the in-itself 31 (inasmuch as it is assumed to be something original and inde-
pendent of all living construction and to possess this same guiding meaning)
turned out to be incomprehensible without a not-in-itself. Therefore, in the
understanding it turned out to be no in-itself 32 (i.e., something comprehensi-
ble by itself alone) at all, and instead it [turned out to be] understandable only
through its correlative term. Therefore, the unity of understanding, 33 which
reason presupposes here, cannot merely be a simple self-determined 34 oneness;
instead it must be a unity-in-relation, meaningless without two terms which
arise within it in two different connections: in part as positing one another
and in part as negating one another, thus the well-known “through” and the
five-foldness recognized in it. If one must also now concede that once oneness
has been admitted, the terms incontestably posit themselves genetically, still
oneness itself is not thereby explained genetically; hence it is present simply
by means of a proiectio per hiatum irrationalem, which this system, [presenting
itself ] 35 as realism, has made against its own principles.
With this disclosed, absolutely everything in the in-itself which
pointed to relations was to be abandoned, and so nothing else remained
behind except simple, pure being as absolute, self-enclosed oneness, which
can only arise in itself, and in its own immediate arising or life: which there-
fore always arises from a place where an arising, a living, simply occurs, and
it does not arise except in such an arising, and therefore occurs as absolute I,
as today’s disclosure could be put briefly. Generally, one can think this sim-
plest of all insights in indefinitely many forms, if it has once become clear. Its
spirit is that being exists immediately only in being, or life, and that it exists
only as a whole, undivided oneness.C H A P T E R
16
Sixteenth Lecture
Tuesday, May 15, 1804
Honorable Guests:
The fundamental principle presented now: 1 “Being is entirely a self-enclosed sin-
gularity {Singulum} of immediately living being 2 that can never get outside itself ”
is in part immediately clear in itself, and in part it has been shown in the dis-
cussion that it is clear to this assembly in particular. 3 Hence, we do not need
to linger with it any further.
I said that it contains and completes what one could present as the first
part of the science of knowing: the pure theory of truth or reason. We proceed
now to the second part; in order to deduce from the first part, as necessary and
true appearances, everything which up to now we have let go as merely empir-
ical and not intrinsically valid. In advance of this undertaking, I must remind
you of only one thing. Resolving this task in absolute oneness of principle 4 is not
without difficulties; especially since, according to a remark 5 about method I
made at the end last time (on the occasion of a general review), this task is
entirely new and has not even arisen in the earlier presentations of the science
of knowing. For that reason, it happens that this resolution cannot remain
without some complications.
However, in order to be completely clear with you about this point, I
will employ here the method that I have generally used before of giving you
an initial factical acquaintance with the terms which come next so as to pre-
pare you adequately for their subsequent combination and connection. This
preparation is the next purpose of today’s lecture. 6
1. In the insight we 7 had produced into inner being, we began—after
fully abstracting from that objectification, about which we already know that
it intrinsically lacks validity—from [244–245] this being’s construction, to
which we expressly challenged ourselves. (You see that I revert to this, partly
as it has always been done up to now, and partly because thereby some kind
of idealistic outlook enters in again. I refrain from giving an account of this
This lecture begins at GA II, 8, pp. 242–243.
121122
The Science of Knowing
here.) 8 Now mark this well, since this point can bring you great clarity, I will
not myself reason {räsonnieren} here as all previous idealisms have reasoned:
“consequently being depends on its being constructed, and this is its princi-
ple”; because this claim could have truth and meaning,—but only in relation
to being’s factical existence in the form of external, objectifying existence,
which existence then [is] absolutely presupposed and so—according to our
basic maxim—the projection through a gap would not be abstracted from. 9
This factical existence in general is to be put into question in its basic princi-
ple and deduced for the first time; but, trusting in the truth of the insight’s
content and so in our principle, we will thus conclude entirely realistically: if
being cannot ever get outside of itself and nothing can be apart from it, then
it must be being itself which thus constructs itself, to the extent that this con-
struction is to occur. 10 Or, as is completely synonymous: We certainly are the
agents who carry out this construction, but we do it insofar as we are being
itself, as has been seen {eingesehen}, and we coincide with it; but by no means
as a “we” which is free and independent from being, as could possible seem to
be the case, and as it actually appears to be, if we give ourselves over to appear-
ance. 11 In short, if being is constructed, as in fact it seems to us to be, then it
is constructed entirely through itself. The basis for this construction, as is
immediately apparent and understandable to us, can not be located external to
being, but only within itself as being, entirely and absolutely; and indeed
absolutely and necessarily apart from all contingency.
At this point notice a) 12 I have said “If being is constructed,” expressing
myself hypothetically, thereby perhaps reserving a future division of the state-
ment into a true part and a false part. So if someone were to insist that 13 it was
actually constructed, you might wonder how on the present standpoint such a
one might conduct the proof? I know of no other way than by means of his
consciousness. 14 However, we have already given up on [accepting] such a
proof as valid by itself; but here for the first time [the question of ] to what
extent and in what sense consciousness and its statements can be [accepted as]
valid is to be decided. In particular, we must decide on the extent to which
consciousness suffices in the highest things it asserts, of which the fact that
being is constructed can serve here as an example. 15 Therefore, we should not
reach ahead anticipating the results of the investigation, which we will make
possible only by means of the hypothetical assertion. 16
[246–247] b) To be sure, it has already become immediately clear in our
earlier investigation that mere being is of itself, from itself, through itself
immediately an esse, that it therefore constructs itself, and that it is only in this
self-construction. This comprised the whole content of our insight. But the
self-construction that we talk about here, which we present only hypotheti-
cally as a declaration of pure consciousness, and which we append to being in
itself only mediately as an inference, this is—as I ask you to grasp {einzusehen}Sixteenth Lecture
123
immediately—something entirely different, merely idealistic, imaginal. This is
the only way I can describe it in words. In contrast, the first one alone would
be real, clearly having the predicate “real” only by contrast with the other, 17 and
thus negating the absoluteness of the previous insight with this predicate,
which is understandable only relationally and through its opposite. Now the
task is just to find to what extent—I say not so much being’s ideal or real self-
construction—as rather the analytic/synthetic principle, which grounds it, is to
be [accepted as] valid. This question of validity can only be resolved by deduc-
ing the principle genetically. 18 Therefore, we take care not to anticipate, and
we grant the entire distinction only hypothetical validity.
2. Let us go back. If being is constructed ideally, as we assume, then this
happens directly as a result of its own immanent essence. Be sure not to over-
look [the fact] that we have thereby 19 actually won something new and great. 20
That is, the ideal 21 is posited in this absolute insight organically and absolutely
in essential being itself, completely, without any real hiatus in essence, and so
without any disjunction in essence. This insight is also genetic, positing an
absolute origin as unconditionally necessary on the condition that it be the
ground and be assumed. 22 Now this insight brings along an absolute that, but
by no means a how; we cannot see how the absolute essence ideally constructs
itself, nor can the inner ground of this construction be constructed further. We
must not be put off by this, since only thereby does this insight secure itself as
the absolute, beyond which there is no other, and this construction as the
absolute one, beyond which no other can be placed. To be sure, it must come
down to such an absolute insight and construction, and it is clear that, only at
this point, with an insight and construction proceeding directly from essence,
could we arrive there. The gap {hiatus}, which as a result of the absolute
insight is in essence nothing at all, exists only in respect of the We; 23 and,
indeed, in case [248–249] the essence of consciousness, properly so-called, is
to consist just in this, no longer in the absolute and pure genesis but rather in
the genesis of the genesis, as it appears here, then, if this We 24 (or this re-gen-
eration of the absolute genesis) were to be deduced, it would be in conscious-
ness that it [the gap] could well remain. [We can then very easily see] that
here, quite probably, we have untied in passing the genuine knot at its root;
and that the new difficulty, which has not concealed itself, has fallen further
down, where it can let itself be easily resolved by closer consideration of the
basic point already discovered. 25 In the meantime, since the point has not yet
been put as clearly as possible, we will continue our exposition without stay-
ing here longer.
3. We have now grasped {eingesehen} 26 this. Following our consistent
method, let us make this insight itself genetic. Under what condition did it
arise? Evidently this, that an ideal self-construction of being be assumed, at
least hypothetically. “It is assumed,” obviously means and can be explained ipso124
The Science of Knowing
facto like this: it is projected absolutely into the form of outer existence, in a pro-
visional way, without any ground or principle for this projection, thus through
an irrational gap.
Now, a major portion of our task is to demonstrate the genetic princi-
ple for this irrational gap, which so far we have presented only factically, whose
validity we have denied, but without our being able to dispense with it. 27
(Observe. A philosophical lecture can frequently count on the unno-
ticed assistance of the understanding, without always providing the distin-
guishing grounds for the distinctions that it makes; the fact usually explains
its true meaning through itself and its results. However, in such cases one
always counts on a happy accident that is just as likely not to occur. It is
always more exact not to leave any distinguishing grounds unexplained; and
especially we should not allow ourselves to be led astray by the fact that often,
and on many subjects, the explanation makes more obscure what was clearer
with the unnoticed assistance of the understanding, because it should not be
so, and satisfying ourselves with understanding’s unnoticed aide is not the
genuinely philosophical disposition. In the hour before last, we had looked at
the case of a distinction between two ways of thinking about the in-itself,
whose distinguishing ground I specifically stated, although the distinction
might have been clear enough simply as a matter of fact. The present case is
similar. The principle for the irrational gap as such, i.e., for the absolute
absence of principle, [250–251] as such, should be demonstrated. Obviously
not insofar as it is an absence of principle, because then it would negate and
destroy itself, a very different thing from it being provided with a principle.
So, in what respect is it to be and in what respect not?—Let us now make the
meaning clear. Being’s ideal self-construction is projected through an
absolute gap, and is thereby made into an absolutely factical and external
existence. Now this existence, 28 as absolute existence, can have no higher
principle at all in the sphere of existence, and in this sense is precisely lack-
ing a principle. Its “principle” in this unprincipledness is just the projection
itself. Hence too it is not claimed, and cannot be claimed, that being in itself
constructs itself ideally, but rather only that it is projected as constructing
itself so; this is important and breaks the doubt aroused by the first remark
at 1. above. Therefore, nothing remains—once this factical being has been
annulled as absolute by the demonstration that the projection is its princi-
ple—except the projection itself, and this as an act, 29 as everyone is requested
to become aware. To say that a principle must be provided for it means there-
fore that a principle must be provided for it solely as an act in general, and as
this act that in itself posits something unprincipled.) 30
What could this principle be? The absolute insight, which forces itself
on us, that the ideal self-construction must itself be grounded in absolute
essence, is conditioned by the presupposition of this ideal self-constructionSixteenth Lecture
125
without any ground, and thus by this projection we ourselves have made to
complete the science of knowing. And so, the principle has been found in
what is conditioned by it, and the newer, higher insight that is thereby created
can be encompassed in the following sentence: If {Soll} 31 the absolute insight
is to arise, that, etc., then such an ideal self-construction must be posited
entirely factically. The explanation through immediate insight is conditioned
by the absolutely factical presupposition of what is to be explained.
4. Now do not forget that everything here remains only hypothetical. If
it should be seen into, then must ____ , 32 etc. Should the consequent be posited
absolutely and categorically? Undoubtedly, if the antecedent is, and without
doubt not, if the antecedent is not; because the latter has no principle except
the former. But if the first should {soll} be posited absolutely, then it is not
apparent as absolute, 33 because it has been posited as absolute hypothetically
{problematisch}. 34 As I add now only to arouse attention, [252–253] in this
hypothetical “should” as our highest point so far, everything comes together
whose derivation is now our task: the ideal construction of being as a self-con-
struction, as well as the projection through a gap. Just so, it is clear that this
hypothetical quality of the “should” must remain as it has been presented.
However, it is equally clear that something categorical must arise too, since
otherwise our science would be baseless and without principle through its
whole range 35 as well as in its starting point. However, this categorical quality
must now just manifest itself hypothetically in the “should” qua “should,” so
that henceforth the chief principle of the process of appearance {Erscheinung}
(and, if 36 it were believed, of what appears {des Scheins}) should consist in this:
that the absolutely categorical “should” appear as hypothetical in relation to
the insight, the true and the certain. That is, [it should appear] as able to be
or not to be, as able to be thus or be otherwise.
5. In order to prepare the way for this point, to the extent time allows,
I urge you to reflect maturely with me on the essence of the “should.” 37 Obvi-
ously, an inner self-construction is expressed in the “should”: 38 an inner,
absolute, pure, qualitative self-making 39 and resting-on-itself. One can assist
the intuition of this truth, which in any case also makes itself. It is, I say, an
“inner self-construction,” completely as such: nothing else 40 supports the hypo-
thetical “should,” except its inner postulation entirely by itself and without any
other ground; because if it had some other ground, it would no longer be a
hypothetical “should,” but rather a categorical “must.” “Inner postulation
entirely by itself ” 41 I have said; hence a creation from nothing, producing itself
entirely as such. A “resting-on-itself ” I have said, because (letting myself take
it up in a sensory form, which harms nothing here) it falls back into nothing
without this continuing pursuit of inward, living postulation and creation
from nothing. Hence it is the self-creator of its own being and the self-sup-
port of its duration.126
The Science of Knowing
This, as we have described it, is simply then the “should,” and, accord-
ing to the presupposition, it is grasped intuitively by all of you in this way.
Therefore, with all its initially apparent hypothetical character, 42 just for that
very reason there is something categorical and absolute here, the absolute
determinateness of its essence. Before we now show further what follows from
this, let me add today two further comments in conclusion.
a. The “should” bears every criterion of the intrinsic being intuited in the
basic principle: [it is] an inner, living from itself, through itself, in itself,
creating and bearing itself, pure I, and so forth; and [it is] certainly
[254–255] organized and coherent internally, entirely as such. As regards
the latter, in case it needs further explanation after the clarity with which
it must have already presented itself previously in intuition: we then always
objectified the fundamental principle’s “inner being” factically, although
this objectivity was not valid. We also have previously objectified the
“should.” 43 Finally, however, we have been lost in it factically, in its inner
description and insight, and now for the first time we free ourselves from
it, and it from us, in reflecting about it, a process which, according to our
previous method, can be explained as a projection through a gap taking as
its principle the “should” itself. Accordingly, this “should”—purely and
simply in its oneness, and without any supplement—can easily be being’s
immediate ideal self-construction, i.e., that is in no way to be further recon-
structed, 44 but rather that provides the subject matter {die Sache } directly in
the construction itself. On the other hand, being’s previous, hypothetically
posited, construction from the “should” has finally found a principle in this
“should,” just as inner being finally has too for its projection through a gap,
which we had proposed accordingly. [This principle] in itself is construction
and subject matter, ideality and reality, 45 and it cannot be one without the
other. This duality may reside in our objectifying consideration of the sci-
ence of knowing, which therefore abandons its claim to intrinsic validity.
b. This “should” has constantly, but without notice, played the principal role in
all our previous investigations. “Should it come to this or that, to a realiza-
tion of the through, etc., then must . . .”; our insights have always gone along
in this fashion. Therefore, no wonder that after letting go of everything else,
what remains for us is only the thing that is truly first in all these cases.C H A P T E R
17
Seventeenth Lecture
Wednesday, May 16, 1804
Honored Guests:
I have stated in the last hour how the part of our science on which we are now
working might be different from the part completed first, and what our fol-
lowing lectures intend: namely to introduce for the first time the materials for
resolving our second task, and to make you familiar with them. At the same
time, I admitted that the next lectures might not be without difficulty and
confusion. It is easier to take in and to grasp {einsehen} that which rests in rea-
son completely and simply as oneness, as did the earlier fundamental principle,
since only abstraction is needed for this task.—“Easier,” I say, than to trace
what in itself and originally is never 1 a oneness back to a 2 oneness in order to
produce a completely new and unheard of concept in oneself, for which other
arts are undoubtedly required. Now, we first lay out multiplicity in an order in
which is most convenient to us for insight. These terms 3 can first be correctly
ordered and understood {eingesehen} on the basis of their principle, 4 which itself
is first to be discovered from 5 them. At this point in the course of the external
lectures, there is an unavoidable circle that can be annulled only by its own
completion. It is possible, and indeed expected, however, for one to grasp the
process—that, to be sure, has its proper order—and terms, and to give them
what clarity they can have under the circumstances. 6 I have said that a new,
heretofore entirely unknown, principle must be presented; and also simulta-
neously I would add this remark: that (thinking of the previous division of the
science of knowing into two parts) we are concerned not just with presenting
the second part, but also with uniting the latter with the first part. 7
The course of the previous sessions was this: we constructed 8 the pure
being, which we had grasped {eingesehen}, as an entirely self-enclosed singu-
larity. In this way, I assumed, we could become immediately conscious of our-
selves, and, as required, we were actually conscious of ourselves. This therefore
was a completely simple, factically objectifying projection of an act that we
This lecture begins at GA II, 8, pp. 256–257.
127128
The Science of Knowing
ascribe to ourselves as likewise independently existing entities; and in this
manner we could have been tempted to deduce being itself from this act of
construction in a one-sidedly idealistic fashion. However, [258–259] we
wisely refrained from doing this, well understanding that by this we would
return whence we had first arisen and consequently would not have advanced.
But we proceeded in this manner, and necessarily had to do so, if we wanted
to come to something more than the one being, for example to the latter’s way
of appearing.—“As concerns the truth in itself of this construction, this can
appeal to nothing else besides the bare assertion of consciousness.” 9 We can-
not now discard this statement unconditionally, as just previously we uncon-
ditionally rejected it, thanks to our present, entirely altered, aim; because pre-
viously we sought pure being in itself, and it has been shown that
consciousness is entirely insufficient for this purpose. Now we no longer seek
this pure being in itself, since we already have it and so our search for it is over.
Instead, we want to grasp it in its primordial appearance; and so conscious-
ness, and here in particular the construction, could be the first term of this
appearance for us to grasp. We cannot allow this term to be unconditionally
valid any more than before; since the extent to which and the conditions under
which it is valid are exactly the issue. Therefore, we must present this claim
hypothetically, without prejudging future inquiry: if, and to the extent that, a
construction of this kind is actual, 10 that is, takes part in being and not merely
seems to be, but has being actually appearing in it, then ____ . Through this
then, we are asked to point out in immediate manifestness the condition for
the real and true being of such a construction, in case and to the extent that
being could come to it. This condition has now been found and has become
evident without any difficulty:—If this construction, which appears to us, 11 is
actually and in fact connected with true being 12 in reason—but in no way con-
nected with factical existence in consciousness, which counts for nothing until
it is better grounded (this detail is not to be overlooked) 13 —if the construction
which appears to us is in this sense, then it is not in any way based in the vain
“I” of consciousness that emptily objectifies being; rather it is grounded in
being 14 itself. For being is one, and where it is, it is whole; in being qua being; 15
therefore entirely and absolutely necessary.
(On the condition that you do not allow yourselves to be distracted, let
me add here an additional remark that can spread much illumination. Posit
pure immanent being as the absolute, substance, God, as indeed it really is,
and posit appearance, that is grasped here in its highest point as the absolute’s
internal genetic construction, as [260–261] the revelation and manifestation
of God, then the latter is understood as absolutely essential and grounded in
the essence of the absolute itself. I assert that this insight into absolute inward
necessity is a distinguishing mark of the science of knowing as against all
other systems. I cannot emphasize it enough, because the absolute absenceSeventeenth Lecture
129
{Dunkelheit} of insight strives against it with all its might, since freedom is
always the last thing [this darkness] will surrender. If it cannot save [freedom]
for itself, then at least it tries to secure it in God. 16 In everyone without excep-
tion, an absolute contingency exists next to absolute substance. Here some-
thing is seen from the beginning as absolutely necessary in reason and in itself,
which afterwards will appear not in reason and not in itself but as contingent
in another connection that still is to be worked out. Only on this condition
can the science of knowing hope to deduce the phenomenon in a genuine and
grounded manner, and not merely as a pretense; because a genuine derivation
must have a reliable principle. Otherwise, as has often actually happened, one
deduces from the intrinsically contingent something else which is also con-
tingent, and obtains other contingent things from these, which themselves
stand firm only on condition of the reliability of the previous thing, whose
reliability likewise depends on the first. As if a good, proper, and reliable
standpoint could arise when one had two terms, neither of which could stand
by itself, each relying reciprocally on the other.) 17
This remark as well: it is evident that in our present investigation it still
seems as if, as I freely admitted at the outset, this investigation is still search-
ing for its principle but has not yet got it, something I have described as erro-
neous, since its first term—the construction of inward being—still remains
hypothetical in connection to that about which alone we are inquiring, true
being in reason. So, the thing which can first be ascertained under this condi-
tion, being’s necessary self-construction, can itself not be otherwise than
hypothetical to the same degree. Therefore, from here on [262–263] you
should direct your attention to the question whether and when a self-sustain-
ing principle emerges.
If there is a construction of being, then it is grounded absolutely in
being itself; we grasped {sahen . . . ein} this directly and reflected further on the
insight and its inner, law-governed form. Then it was immediately clear that
we began with the presupposition of inner being’s construction, which we
incorrectly attributed to the “I” of consciousness, but we have already learned
better than this and let go of the attribution. But this much remains indu-
bitable, that being’s construction is projected as an absolute fact.—Have we
now brought this implicitly simple, factical projection into connection with
other terms by the use we have made of it? Evidently so; we saw {sahen wir
ein} that if such a construction exists, then it must be grounded in being. Now,
we have undertaken this entire speculative venture freely; the resulting insight
(which might very well not have been engendered) is conditioned by our pro-
cedure (which we might very well have omitted), and therefore it is in no way
a firm standpoint. All the same, in order to achieve such a firm standpoint we
applied a procedure that, to the extent that it needed to, proved its legitimacy
by its bare possibility. We said: assume that the insight, engendered by us, is130
The Science of Knowing
to arise, then you will see {einsehen} that under these conditions the projection
of factical being, previously only possible, becomes necessary.
In this way, we would first of all have made good progress beyond every-
thing achieved so far toward a proof that, although to be sure we do not feel
firm ground beneath us, we might be on a good path. The absolute projection
through a gap, and thereby the form of outer existence, that could not be
understood conceptually in all our previous investigations, is explained as nec-
essary under the assumption that a higher term (the insight) should be, an
assumption that itself was previously hypothetical. Thus, the hypothetical sta-
tus departs from the lower term, though only by transferring itself to the
higher; but at least with this it is simplified and its proper location is revealed
to us, where we can hope to grasp it at the root.
After what I have said about the necessity of a self-supporting principle
for this investigation as well, we will next eliminate this hypothetical status
completely; and here the most secure means is to look it straight in the eye. It
is entirely compressed into the hypothetical “should”; this is sufficient by itself
for our next purpose; 18 therefore, we let go of the site where this “should”
appears, insight, etc. Quite apart from our current procedure, it could be obvi-
ous from the entire previous investigation that one now needs to keep this
“should” as one of the deepest foundation points of all appearance, [264–265]
as I will observe in passing. All our preceding investigations and engendered
insights have started with the hypothetical “should” and have proceeded from
it as a terminus a quo: “If there is really to be {soll} a “through,” then there must
____”; “ Should the achieved insight arise, then there must ____”: idealism;
“should this life be life in itself, then there must ____”: realism; all the way up
to the highest relation: “Should an in-itself be comprehensible, then a not-in-
itself must be thought” and so on. 19
This “should” loses itself entirely only in the insight into pure being and
into the way in which we evaporate into it, so that an absolute categorical
character enters, without any hypothetical presupposition. As soon as we
reflect again on this insight, the process which yielded the historical origin of
our second part and our entire present investigation, it 20 reinstates itself with
a “should,” thus as something contingent, seeking the basic condition for this
contingent quality, a necessary self-construction of being. Now, so far in the
ascent we have clung to the content of the generated insight without reflect-
ing on the hypothetical form in which, as a whole, it appeared. This was
entirely correct because we wanted to arrive at the original content of the truth
as such. (Here in passing the question that some have asked me concerning
the true grounds for our first part’s preference for realism and for the maxim
that ruled there always to orient ourselves realistically answers itself decisively
and fundamentally.) 21 But now as we descend we have to hold on to just this
neglected “should,” which indeed provides the enduring inner soul of all theSeventeenth Lecture
131
idealisms, which consistently excluded themselves during the ascent, and
which were struck down by an opposed higher term only in respect to their
content; but still persist in their form, as we see. Now this form cannot be dis-
turbed directly by the original content, since everything that the latter can
attain has already been achieved in the ascent. Rather, it must be explained
and justified inwardly on its own terms. It must refute its own ungrounded
claims, to the extent that they are ungrounded; roughly just the way we refuted
the content’s highest idealism, that at first presented itself as realism, by means
of the law that it itself presented, and revealed it as idealism.
In a word, and in order to lead you even deeper into the systematic con-
nection between
a. the term that earlier was the highest in appearance—the distinguishing
and the unification of the in-itself and [266–267] not-in-itself in the whole
five-fold 22 synthesis, and
b. absolute inner being, as the absolutely realistic element, the “should” 23
enters here as a new middle term, in which the self-differentiating and
likewise synthetic relation of the two indicated relational terms must show
itself. To find this is the proper content of our task: to find it as a firm prin-
ciple is its form.
First, however, the connection to inward being. The form of being is
categoricalness. Therefore, something categorical must be found in the
“should” 24 itself, however hypothetical it might appear. In order to uncover this,
I have demanded that the inner essence of a hypothetical “should” be carefully
considered (following our consistent method of raising into clarity something
that was at first dimly projected).—We have already done this last time;
because of the subject’s importance, I will repeat the entire operation today.
If you say forcefully and deliberately: “should so and so be,” then it is clear
that thereby an inner assumption is expressed, without any foundation, simply
of itself and from itself, thus an inwardly pure creation, and to be sure stand-
ing there completely pure entirely as such, 25 because the “should,” if it is taken
only as purely hypothetical—as is required here and without achieving which
the required insight will not arise—expresses complete external groundless-
ness, simple internal self-grounding, and nothing else. Further, (in this way, I
tried to grasp the same thing and make it clear from the other side): the
absolute assumption 26 is expressed in the should, an assumption that is uncondi-
tionally allowed to drop, just as it is unconditionally presupposed. Should it
(and with it probably the entire “If . . . should, then ____ must” 27 that depends
on it) not drop away, (with which dropping away all knowledge and insight
probably drop away as well) then it must hold and sustain itself.—As surely as
we have now seen into this, just that surely has the should 28 been illuminated132
The Science of Knowing
for us as an absolute that holds and sustains itself out of itself, of itself, and
through itself as such, on the condition that it exists. This, I say, is a “should”;
and were it not precisely so, then it would not be a “should”; therefore we have
a categorical insight into the unchangeable, unalterable 29 nature of the
“should,” 30 an insight in 31 which we can completely abstract from the outward
existence of such a “should.” “Can abstract,” I say since, with adequate deliber-
ation, I refrain from drawing a conclusion here which easily presents itself, but
which is not yet sufficiently ripe, given the context. To the extent that our task
simply consisted in discovering something categorical in the “should,” it has
been fulfilled by what has happened.
In explaining the should, I have not warned you about the illusion that
it is we who assume there what is hypothetical and who hold [268–269] and
sustain it; since the rule is to leave this “we” of mere consciousness entirely out
of action until it is deduced, and being able to do so is the art without which
no entry into the domain of the science of knowing is possible. If, in the
meantime, this “I” has forced itself on anyone, then let it immediately remove
itself at this point. Namely, whether or not you have created and carried the
assumption, it is still always completely clear that you have a “should ” 32 only
on this condition of self-creation and carrying forward. Therefore, even if you
are the creator, the “should ” 33 always contains the rule and law of proceeding
in that manner, otherwise it is not a “should,” and we have not wished to say
any more than that here; abstracting completely from the question that you
raised and that we will work out in another place.
And now, in conclusion, a very 34 sharp distinction, that will become
decisive in what follows, and that cannot be made clear too soon. The strong
similarity between inner being as something self-enclosed and self-sufficient
in-itself, of-itself, and through-itself, and the should as just the same has
already been pointed out earlier. There is nonetheless a distinction between
the two that I have named and made dimly 35 recognizable in the stated for-
mula: “the should” is something in-itself, etc., as such. 36 I urge you now to clar-
ify this distinction for yourselves along with me. Being was constructed as
something absolute in itself, etc. I ask: should there now exist, or is their actu-
ally in our insight, if it is of the right kind, another persisting being or sub-
stantive, besides this absolute, self-constructing esse? Not at all. Instead both
merge into each other 37 and into the pure self-enclosed singularity, and the
doubled repetition is entirely superfluous, insufficient, and neglected. This is
not at all the case with the “should,” if you will look into it quite acutely. 38 The
latter stands out as a fixed, substantial middle point and bearer of absolute self-
production and continuation. 39 The latter is not just immediate, as was the case
with being, but rather only mediate through presupposing and positing a
“should”—in brief on the assumption that the “should” itself again should be,
and thus should be seen {eingesehen} through its own doubling. Here there isSeventeenth Lecture
133
not, as there was before, an immediate rational insight, but rather only a medi-
ate one, conditioned again by a higher projection through a gap, precisely of
the “should ” ; 40 just in the way we have actually proceeded. We have wished to
indicate this relation by the added phrase “as such,” 41 i.e., itself in objective,
factical oneness of essence.
To what further things this new discovery might lead must emerge on
its own. Before [270–271] hand, this much arises in regard to method: that,
just as a projection through a gap (the projection of being’s construction) is
deduced as necessary from the fact that a particular insight “should be” {sein-
sollen}, 42 another projection, just that of the “should” itself, presents itself [on
the one hand] as a condition for this insight and on the other side again as
conditioned by it. We now need to venture further into this; that therefore our
present investigation, just like the previous one, advances upward only in this
precisely delineated circle, because it is still looking for the latter’s 43 principle. 44C H A P T E R
18
Eighteenth Lecture
Thursday, May 17, 1804
Honored Guests:
[Here is] what has been presented so far: we presuppose a construction of
being. On the principle that nothing can be except being, the construction is
seen {eingesehen} as arising necessarily from being—of course, with the same
certainly it has generally; supposing therefore “If it should be . . . , then ____
must be.” 1 But the “should” is something in itself, of itself and from itself as
such.” This, and in particular the “as” most recently added, is now firmly fixed
for you as another new middle point and bearer for the self-producing and
self-sustaining “should.”
Today I add another basic observation concerning the true inner spirit
of the reasoning processes presented so far, and [we] will then work on our
remaining task from another angle.
1. As concerns the first, our higher insight from the standpoint of the
hypothetical “should” took the following form: should an insight into this or
that occur (in this case in particular the insight that the ideal self-construction
is to be grounded in being itself ), 2 then ____ must. “Since you now,” I would
say, “actually provide 3 the content of this insight, which, according to your
account hasn’t yet occurred but whose condition you are seeking, you already
without doubt have it in sight and in your concept; you are constructing it
really and in fact”; (as is the case here with being’s ideal self-construction). 4
This remark permeates all consciousness and can be illustrated in every
case. I cannot reflect how and according to what law anything (e.g., a body in
space, space, a line, etc.) is conceived or constructed, unless I have already
grasped it apart from all reflection and according to a universal law. In the pre-
sent case, the law is constructed in one of the most general cases, which con-
tains others within itself. “Therefore,” I continue, “you seek either that which
you already have, or you seek the same vision and the same concept 5 (the same
in regard to contents 6 ) only in another qualitative form {Bestimmung}. That
This lecture begins at GA II, 8, pp. 272–273.
134Eighteenth Lecture
135
the latter is the case becomes clear through a more exact consideration of the
proposition laid down. 7 The content of your vision, which, as the content of
mere seeing, is separated and existing for itself, 8 should be brought into con-
nection 9 with something else in seeing, both as its condition and as condi-
tioned by it, initially through what you call insight. 10 Thus, in order to state the
true result of your desire definitely and exactly: just to arrive at your demand,
a seeing that is already completely determinate in and of itself (and that you
must presuppose as determinate) should [274–275] be further qualitatively
determined in this persistent, objective determinateness as seeing, since the
objective determinateness remains. Therefore, to put it briefly, you demand a
new genesis in the seeing that has already been presupposed as existing and as
remaining the same objectively.
A new inner genesis of seeing, as formal seeing itself, without any alter-
ation in the [seen] content ; 11 (what we have already called objectivity 12 ). Now,
the material of this formal genesis, its result, is itself again a genesis: the con-
stant content should be brought into a genetic relation with another term,
that creates, and is again created by, it; thus the entire familiar “through,” or
the relation in its synthetic five-foldness, should come in. As things stand, it
can well be that this external material genesis with and out of the content,
which is nonetheless not changed in its inner nature, is itself grounded in
mere seeing’s formal genesis, and resides not so much in the subject matter
as in the altered eye, through which the entire present multiplicity is traced
back to the oneness of the same principle, of the formal further 13 determina-
tion. This formal further determination, or new genesis, is called for through
a “should,” which has itself been recognized as a genesis in its inner nature
unconditionally as such. And so this genesis could have its ground in the
“should” 14 itself as the relation and five-fold synthesis within the formal gen-
esis, so that the “should” is the basic principle for everything, as we have pre-
viously already taken it to be. In brief, the spirit of our whole reasoning
process, conducted since the beginning of the second part, is the demand for
an inner genesis 15 in the seeing presupposed for genesis 16 itself. This process
adds nothing to the seeing in its true meaning, and so it must be inoperative
in relation to this meaning, just as we have always wished. Likewise, this very
inner formal genesis, as wholly concerning only the way of viewing, 17 may be
the principle of absolute idealism = of appearance; and we ourselves have
entered into a new and higher idealism through the principle presupposed in
our entire reasoning process: that being is constructed ideally, i.e., separately
from its real self-construction.
That just this insight, now characteristically distinguished from the pre-
supposed original seeing, presents itself alone as certain, compared to which
the original seeing is to be only hypothetical in relation to its content (it is
clear on immediate reflection that the matter stands so, and our certainty136
The Science of Knowing
appears finished and closed)—this circumstance probably lies in the partiality
of idealism itself, which here gives testimony for itself, knowing nothing else.
Now we have to investigate this claim for the first time.
[276–277] 2. A recognized basic rule: nothing can be accomplished in
any way against an idealism except from the standpoint of realism. Therefore,
as soon as our reasoning has been traced back to its spiritual oneness and
understood to be idealism, we cannot stand by it any longer without being dri-
ven around in circles. We must turn instead to the corresponding realism and
consider this more deeply in its origins.
a. As we remember, we entered this realism after the last consideration
of the in-itself, and of the insight that, in our knowing, this in-itself is relation
and multiplicity; therefore, that it is not absolute oneness, thinkable without
any composition or division, but is rather, as we said: a oneness of understand-
ing {Verstandes-Einheit}. 18 We discarded this knowing entirely, and yet know-
ing still remained, which thus was absolute inner oneness, without any com-
bination or separation: oneness in itself. We also refrain from saying for
example that we have produced it in this oneness; since we truly would not
have wished that something should remain behind after abstracting from
everything, or [have wished] to encompass what remains with our will, had we
willed or been able to will this, so that it would indeed have been left over for
us: instead it was just unconditionally left over: oneness of itself. Everything
depends on this last point; it is what has been overlooked in every system and
what becomes clear only to the deepest deliberation. What we are naming the
We, that is our freedom, which is derived here for the first time from the pre-
viously mentioned, new formal genesis of the absolutely presupposed seeing =
re-construction, can only abstract from its own creation of the act of recon-
struction, but it cannot creatively construct primordial reason; although after
complete abstraction primordial reason enters without delay. So then anyone
who—in inseparable 19 awareness of the simultaneity of his completed abstrac-
tion and the arrival of primordial reason, and in the equally inseparable 20
awareness that he is the one freely abstracting—immediately transfers his own
freedom to reason’s emergence, such a person deceives himself and remains
trapped in an idealism. This final illusion is negated here in immediate man-
ifestness by means of deep reflection. After abstraction from the highest one-
ness of understanding, a knowing 21 remains, just because it remains, without
any possible assistance from us, pure light or pure reason in itself. 22
b. This pure reason is equally immediately inner being 23 and completely
one with it. Previously 24 we called what remained after all abstraction “inner
being”; here we have called it “pure light,” or “reason.” [278–279] But what-
ever we may wish to name it, it is what remains unconditionally by itself after
all abstraction, an entirely indivisible singularity; and I would very much like
to know whether any disjunction can be made in the presented concept, andEighteenth Lecture
137
whether the insight that it is a completely self-contained singularity does not
clearly show that, whatever variation in the words used to name it, one and
the same nature could be meant.
c. Previously as well as now, we have named it a real self-construction in
itself, of itself, and through itself, and we could not describe it differently. Now,
abstracting completely from the facticity of this description, which to be sure
can only be a reconstruction, and through which we happen into the first
named idealism, [let us] reflect [instead] on its inner truth, and—with this I
ask for your complete attention—on the surprising result that I intend to
bring out. I ask: does it not now depend entirely on the pure thing itself
remaining after all abstraction that it exists entirely from itself—whether you
call it being or reason? 25 For example, is it arbitrarily posited as existing on its
own? How could it be? For this would be a genuine contradiction, since in that
case it would not be from itself but would exist through an arbitrary act of
positing. If it is posited as something left over after abstracting from every-
thing outside itself, then it is necessarily posited as of itself. For if it were not
of itself, then it would be of another, 26 so that in its absolute positing—i.e., in
the original creation of its being—it would not be possible to abstract from
this other. (That because of babble and thoughtlessness this other might not
be considered could be factically true and still should be explained ; 27 it is not
true in the one absolute, self-consuming oneness.) Once again, it is posited
absolutely, creatively, as something of itself 28 —it is evident that this of itself is
actually manifested and is not just thought up;—so it is posited as existing
absolutely and remaining behind after abstraction from everything.—Hence
it is clear that light, or reason, or absolute being, which are all the same, can-
not posit itself as such without constructing itself, and vice versa: that both
coincide in their essence and are entirely one. 29 —Notice here: 1. the insight
that being must construct itself unconditionally has arisen here through the
mere consideration of its inner nature entirely immediately and without any
factical presupposition, an insight that, according to idealism’s pretensions,
should only be producible mediately from the factical presupposition that a
constructive act is present. By this means idealism is first of all fully refuted,
insofar as it grounds itself in the necessity of a presupposition for a [280–281]
particular insight, although merely a possible one, since the insight has actu-
ally been produced without the presupposition. Idealism must therefore look
around for higher support, if we are still to come to it. Further, the proposi-
tion alluded to in passing has therefore come up, that this same insight is pos-
sible in two different ways: mediately, from presuppositions, and completely
immediately. How would it be, if the entire distinction that we have sought
between philosophical and common knowledge, between the standpoint of
the science of knowing and that of ordinary knowing (and in case within the
latter there should be degrees of mediatedness, the distinction between the138
The Science of Knowing
various standpoints of this common knowing) were to lie in just this distinc-
tion between these differing ways. Philosophical systems are always closest for
us:—the presupposition that idealism wants as the principle of mediate 30
insight, is factical. How would it be if, for example, the proof of absolute being
from the factical existence of finite entities, which is conducted in nearly every
system, and according to them in ordinary consciousness as well, were just this
idealistic path of mediate insight, with which one remains satisfied, for the
lack of the immediate path. In itself this is correct and is applicable in its place
within the gradual process of cultivation, i.e., in rising up to the highest; but
it generally fails the test against criticisms that strive ahead to the highest! 2.
The distinction between being’s real and ideal 31 self-construction that we
made earlier, and on which idealism built, is now completely annulled. Being,
or reason, and light are one; and this one cannot posit itself, or be, without
constructing 32 itself; this is therefore grounded 33 in its nature, and is entirely uni-
tary, as is its nature. Therefore, if we are to return later to such a distinction, 34
then it must first be derived. 35
3. We saw {einsahen} that, in reason per se, its self-positing and its self-
construction 36 as “from itself,” etc., coalesce entirely into one. 37 And as cer-
tainly as we saw into this, in this insight 38 we were the oneness of reason itself.
Now a duality still remains here, not however as in the oneness of under-
standing, whose parts are to be integrated—since parts within the oneness
are rather completely denied and negated here; and the oneness does not
understand itself through parts but rather posits itself unconditionally and
absolutely—but rather as a means for achieving oneness. Therefore, it may
perhaps turn out that a re-construction is already present here, one that
would be posited backwards toward the idealistic side by an absolute
“should,” and which we could not avoid merely factically, even though its
intrinsic validity is not admitted; therefore that we stand at the precise place
from which our task could be completed. How things may stand with this I
reserve for further investigation.
[282–283] Now I add a supplementary remark, with which I did not pre-
viously wish to interrupt the course of the inquiry. As the opportunity has
arisen, I have tested those recent philosophical systems, which have made the
greatest impression in regards to their principles, in order thereby to bring
greater clarity to the science of knowing; thus Reinhold’s system and thus
Schelling’s system. Next to these, and perhaps even more than they, Jacobi’s
system recommends itself, because with great philosophical talent it tries to jet-
tison philosophy itself, and thus it flatters the prevailing spiritual indolence and
denial toward philosophy. The scene for testing this system’s principles was just
above. It proceeds from the following principles: a. We can only re-construct
what originally exists.—We ourselves have precisely presented and precisely
defined this claim, which for Jacobi is almost a postulate: the seeing, deter-Eighteenth Lecture
139
mined primordially in its content, is formally genetic 39 in relation to an unal-
tered content, and therefore it is the insight into a connection; and we ascribe
this genesis to ourselves, a genesis that is only re-construction in relation to the
truly original content and that is truly original construction and creation from
nothing in relation to the terms added factically.—Regarding the last point, the
absolute creation of everything factical from the I, he very clearly took this over
from us; and it is very plausible that he granted being to the factical, i.e., to
what is sensible outside the one rational being, and thereby left us only recon-
struction. b. Philosophy should reveal and discover being in and of itself.—Cor-
rect, and exactly our purpose.—Through the persistent assertion of these two
principles this author has earned the age’s great thanks, and has favorably dis-
tinguished himself from all of the philosophers who just reconstruct impar-
tially, [284–285] or even just fool around with nature and reason. c. Therefore,
we cannot philosophize, and there can be no philosophy. 40 This latter claim, just
as I have stated it, is his true opinion, and must be his true opinion, if he is to
have any opinion at all. For he contributes nothing by his usual addition: phi-
losophy as a whole. 41 Because, if there is no philosophy as a whole, then there is
no philosophy at all, but rather only edifying remarks for every day of the year.
I grant him everything as it is presented, only taking it more seriously than its
original proponent does. We, the we who can only reconstruct, cannot do phi-
losophy: equally there is no philosophy individually and personally; 42 instead
philosophy must just be, but this is possible only to the extent that we perish,
along with all reconstruction, and pure reason emerges pure and alone; since
this latter in its purity is philosophy itself. From the perspective of “we” or “I”
there is no philosophy; there is one only [once one has gone] beyond the I.
Therefore, the question about the possibility of philosophy depends on
whether the I can perish and reason can come purely to manifestation. This
author could demonstrate that this must indeed be possible from his own
words. Because when he says: we can only reconstruct, he achieves ipso facto in
that very moment something more than re-construction, and has at least drawn
himself happily out of the “We” of which we have spoken. For if he could [do]
this, then for his whole lifetime he would enact {thun}, but without speaking
about it, just as by his previous statement he enacted elevating himself to
reconstructing the reconstruction. Of, if we will free him from this, he [can] tell
us how he came to the universal statement by which he prescribed an absolute
law for his “We,” and thereby pre-constructed the “We’s” essence for them, and
did not merely re-construct it. In which case he would have to resign himself
to express himself like this: “I and everyone I know, as many as I can remem-
ber to the present could only re-construct; whether perhaps tomorrow
[286–287] something else will happen, we will have to see.” Finally he will
have to tell us whether he understands this concept of “reconstruction” without
presupposing something original, independent prior to all construction. As140
The Science of Knowing
surely as he understands himself, he must become aware of such a thing beyond
all reconstruction. Grasping this original something and reconstruction as fol-
lowing from it as an absolutely essential law of the “We,” just as we have artic-
ulated it here is the task of a philosophical system, which we have presented
entirely according to its sense, but have only partially solved.C H A P T E R
19
Nineteenth Lecture
Friday, May 18, 1804
Honored Guests:
Since today we finish the week, I do not wish to let you go without having
equipped you with some definite result. This resolution compels me for now
to pass by certain middle terms that still remain for deeper consideration
between that with which I ended yesterday and that which I will attach to it
today, in order to reserve them for the descent.
1. As an introduction to our essential business for today, [here is] a clar-
ifying remark that should direct your subsequent attention and that at the
same time also briefly and concisely repeats the first major part of yesterday’s
lecture! I say: in all derivative knowing, or in appearance, a pure absolute con-
tradiction exists between enactment {Thun} and saying {Sagen}: propositio facto
contraria. (Let me add here by the way, as I thought previously on an appro-
priate occasion, a thoroughgoing skepticism must base itself on just this and
give voice to this ineradicable contradiction in mere consciousness. The very
simple refutation of all systems that do not elevate themselves to pure reason,
i.e., their dismissal and the presentation of their insufficiency even though
their originator is not thereby improved, is based on just the fact that one
points out 1 the contradiction between what they assert in their principles and
what they actually do [in asserting them]: as has been done with every system
that we have tested so far, and yesterday with Jacobi’s as well.) In the first half
of yesterday’s talk, this contradiction showed up in what we had identified so
far as the highest principle of appearance, that is in the “should,” 2 immediately
after we had conceived it in its firm and completely determinate nature as
something from itself, etc., as such; namely, a particular insight (which in our
case was this: that being constructs itself ) is posited through the “should” as
not present but rather merely as possible, and as possible only under a certain
condition that is still sought. If we are even to arrive at the consideration of
its conditioned possibility, this [condition] finally must be presupposed as a
This lecture begins at GA II, 8, pp. 288–289.
141142
The Science of Knowing
seeing that is fixed in its content and to that extent unchangeable,. Hence,
[two things] stand in complete contradiction in this “should,” its enactment
{Thun}—its true inner effect, to presuppose a seeing that is unalterable in its
contents—and its saying—a different action {Thun} on its part, according to
which the insight is supposed to be not actual but only possible under a con-
dition yet to be added. I add only for the sake of recapitulation that the true
external nature of this “should” is found as the [290–291] demand for a fur-
ther inner and merely formal determination of a presupposed seeing that is
unalterable in its content, through which further determination this presup-
posed seeing comes into a genetic connection with another term that is cre-
ated purely by this further determination. And I immediately formulate the
following conclusion: absolute reason is distinguished from this relative know-
ing by the fact that, in the case of absolute reason, what exists, or what it does,
is expressly said in it; and that it does what is expressed in it in absolute qual-
itative sameness.
2. In the second part of yesterday’s investigation, we tried to represent
pure reason in ourselves. I noted at the end of this presentation that, because
of the duality 3 that to be sure was annulled intellectually but that remained fac-
tically inextinguishable in you, it became evident that pure reason could not
display itself immediately in you and could rather only be reconstructed.—The
same qualitative determination of a presupposed seeing that is unalterable in its
contents, a determination pointed out within the “should,” we also called
reconstruction; therefore, the contradiction between saying and doing just dis-
covered in all derivative knowing is contained within reconstruction itself, a
fact that can itself be made clear immediately. To be sure, reconstruction explic-
itly puts itself forward as reconstruction and therefore in its own concept quite
properly posits the point of origin {das Ursprüngliche}, and in that there is no
contradiction. But since it leaves the content unchanged and can actually cre-
ate nothing new without completely negating the relation between itself and
the absolute—its construction is therefore groundless and the fact itself con-
tradicts the postulate of the absolute necessity in the pure, positive in-itself. 4
I should now immediately climb past this contradiction [we have] dis-
covered and relieve it (that is, past the groundlessness of the concept of a
reconstruction). However, in accordance with my initially stated resolve, I am
retaining it [in order] to annul it mediately in the descent; and so [let’s turn]
directly to yesterday’s reasoning to indicate the location of absolute recon-
struction, 5 and to remove this circumstance. 6
We brought the already established absolute insight to life in this way:
a. it arose for us after we abstracted completely from all relations, and it
remained behind as a oneness, not just because we wished it, but 7 simply by
itself. Pure light, or reason.Nineteenth Lecture
143
b. Previously we named it inner being, here light or reason; but it is clear that
no distinction whatsoever occurs in the one singularity that remains behind
by itself as one, and that consequently both designations are only two dif-
ferent names for the one that is grasped {eingesehene} as completely indi-
visible and inseparable. 8 [292–293]
c. We saw into this “one,” and still see into it now as something from-itself,
etc. = [as] self constructing. 9
I asked: should not then this from-itself 10 reside completely in its nature
as absolute truth? 11 And I discussed it even further in the following consider-
ation: out of its self-positing as this, 12 self-construction follows, and vice versa;
because if it is posited as this, 13 as remaining after abstraction from everything
else, then it is posited as remaining and persisting because of itself; since if it
were not because-of-itself, it would be because-of-another, from which it
would not then be possible to abstract in its true original creation, or which
could not be absent for this creation. Conversely, if it is a true, actual, energetic
from-itself, then it is not from another; 14 since then it would not be truly from
itself. Therefore, it is necessary to posit it, as it has been posited.—But let us
take a keener look at this reasoning itself and the procedure within it.—(And
I remind you that this is the most difficult and significant thing that has so far
come before us.—) 15 First of all, without exception in our whole argument
process and in the entire conduct of our lectures up to now, the absolute has
been treated as what is left over after abstraction from everything manifold;
and if equally we have expressed specifically enough the absolute from-itself
and pure oneness in-itself, then with these words which we have added as
clarifications, we have surely again made use of this same relation; as more
certain 16 proof that even we ourselves, the scientists of knowing and what we
actually did and pursued, found ourselves in the previously uncovered contra-
diction between saying (of the from-itself ) and doing (explaining by means of
the not from-itself ). 17 Thus the first premise of our proof here reads: “If it is
posited as this, i.e., as left over after abstraction from everything else . . .”; 18
which is a sure demonstration of reconstruction. Second, in the center of our
entire proof we have absolutely presupposed both genesis and the absolute
validity of the Law of Principles. 19 The center of the proof was, “If it is not
from another, then it is from itself; and if it is not from itself, then it is from
another.” If someone now were to say to us: “Quite right: one of the two—
from itself or from another—and in case one, then not the other, if of course
I grant you the use of your “from” at all. But if I say instead: in brief it is, and
that’s all, who will then ask about a “from” (Von)?” To be sure we can answer
such a one as follows: “You are reflecting; so in addition to this “is” you also
have consciousness; you therefore have not one but two, that you can never
make into one [294–295] and an irrational gap lies between them; you are in144
The Science of Knowing
the familiar death of reason”;—so the loophole always remains open to him
that is taken by every non-philosopher: 20 “I must just stay in this ‘from’ and it
is impossible to escape it”; so everything finally comes down to this that we
justify ourselves in the use of the “from.” 21 Therefore this would be our next
task, to justify the “from” in general as such, entirely abstracting from its appli-
cation. So far, as I ask that you recall, it has not arisen in any other way than
in factical necessity.
This justification will disclose itself if only we rigorously pursue the
analysis of the preceding argument. In the first half are to be found the
remarkable words that without doubt became immediately evident and clear
to you: “if it were not because-of-itself, it would be because-of-another from
which it would not then be possible to abstract in the true primordial act of
creating {Creiiren}, or which could not be absent for this creation {Creation}”; 22
and yesterday I also added: “even for truly primordial creation,” since through
thoughtlessness and foolishness one could easily forget the other through 23
which alone the first can be. What then is understood by this primordial cre-
ation which likewise in total tranquility provides the center of the proof? Evi-
dently that our thinking, or the light, if it should be of the right kind, must
accompany the genuine real creation of things and originate along with it:
hence if the one were to be through another, it would have to take the
“through another” up into itself and express it; contrariwise, a thinking that
omitted this “through” would be mere thinking and not absolute, and would
set down a true creation only factically as bare, dead existence. This was the
first point. 24
Now it seems here as if the real creation, as real, could exist on its own
and go its own way; and some assert it. The basis for this illusion has in fact
been grasped here. That is, it rests in the possibility of viewing primordial cre-
ation too in a pale and factical way, as a result of which it seems to be capable
of existing independent of, and separated from, its appearing {vom Blicke}. But
we have already seen earlier that light and inward being (by no means the
external existence created by faded thought) are entirely one and the same; or,
in case we had not yet realized this, then this is the place to prove it immedi-
ately; because if absolutely unchanging and unchangeable self identical light
must accompany creation, then there is no light without creation and creation
is likewise inseparable from light: since it is only because of the light and in the
light. 25 Creation = “from,” “through,” etc., so absolute light is itself an absolute
“from.” This was the second point. 26
[296–297] Now we, the scientists of knowing, have tacitly presupposed
this as the inner principle of the possibility for the entire subordinate proof
procedure, which we are now dropping, and indeed—this is the important
thing—we have done it without any design or plan before the deed and imme-
diately through the deed itself. But I claim that the bare possibility of this pre-Nineteenth Lecture
145
supposition shows its truth and correctness. Let me prove this first indirectly.
We ourselves in our doing and pursuing are knowing, thinking, light, or what-
ever we wish to call it. If knowing were now absolutely limited, i.e., to the
faded thought of an existence separated from thinking, then we could never
have been able to get out of it to this presupposition of an absolute creation.
Since we have really posited it, and the light as absolutely one with it, since
we ourselves are immediate light, we have quite certainly validated the truth
of our claim in immediate being, in action, since we enacted in that place the
very thing we said, and said what we enacted; and the one could not be with-
out the other. Results: 1. The contradiction we have noted up to now in what
we ourselves are and pursue, between doing and saying = the real and the ideal,
is now annulled, as it alone can be, ipso facto in us ourselves, and since this is
the criterion of pure reason, we are ipso facto pure reason. 2. Light has a pri-
mordial conception of its own nature that ipso facto preserves itself in imme-
diate visible completion of itself. (Note well: here we are holding simply to the
immediately evident content of our sentences. It is obvious that questions can
still be raised about their form. These questions will raise themselves, and the
basic principles for deriving relation from the absolute may well lie just in
answering them.) 3. On the very grounds given, let us leave our factical con-
ception of the nature of the light, which may well give rise to the entire we 27
whose origin we are seeking, and let us hold simply to the content. In light,
absolute genesis. 28 Obviously, the light, as light, is qualitative oneness (which in
fact enters as just plain seeing that cannot further be seen), which permeates
the entire inner genesis 29 as bare pure genesis (I am relying here on your [pow-
ers of ] penetration; since language can in no way bring us to our goal 30 —). I
can now construct this for you further as follows: this oneness 31 permeates the
duality in the “from a–b” ; which duality exists only in the absolute “from”—but
not at all outside it in some independence and [in an] independent differen-
tiation of terms—so that [the duality’s terms] may be reversed with complete
indifference.—These all are [298–299] constructions in sensory terms,
through which I anticipate myself. The ground of their possibility must lie in
and be derived from me myself, insofar as I am the factical concept. In the
strictest sense, nothing matters more than this: light is the qualitative oneness
that penetrates the “from.” This was the first point. 32
Now likewise, following our concept, this “from” and (just for that reason
and consequently) the light’s permeation of it, and also therefore the entire
qualitative oneness of the light, that indeed can only be thought in relation to
a “from” and its duality in order to annul it,—all this, I say, has its ground in
the light itself, 33 no longer as qualitative but rather as an inscrutible oneness.
Therefore, there exists between the light in itself and the entire preceding rela-
tionship a new, only entirely one-sided “from”; and this latter denotes the
absolute effect of the light; to the contrary, the entire first relationship simply146
The Science of Knowing
shows the appearance of this effect, of the qualitative immediately self-effect-
ing light. This was the second point. 34
As a genesis, every “from” posits light;—just as previously light posited
genesis: and indeed, since the absolute “from” of the pure, inaccessible princi-
ple rests here, it posits absolute light, without the genesis ever becoming visi-
ble, and [posits] itself only in this 35 absolutely factical light and from this facti-
cal light.—If you have seen into this, then reflect now on yourself. We have seen
into the “from” in just this way, and by means of it have seen into the 0 whose
inaccessibility we have previously admitted, and have seen into it as uncondi-
tionally existing, objective 36 and so as having to exist, if appearance is to arise.
This is the fact. How have we explained 37 it? Thus: there is an absolute, imme-
diate “from,” which as such must appear in a seeing, itself moreover invisible. 38
Hence we ourselves, with the whole content of our immediate seeing, are the
primordial appearance 39 of the inaccessible light in its primordial effect, and a–b
is mere appearance of appearance. And so the primordial facticity, the original
objectification of reason, as existing and genetic, is thereby clarified from the
original law of light, and our task has been completed in its highest principle.
I have no reservations in letting you go for the week with these provi-
sions.—Monday [21 May]—a discussion.C H A P T E R
20
Twentieth Lecture
Wednesday, May 23, 1804
Honored Guests:
Being is an unconditionally self-enclosed, living oneness. Being and light are
one. Since in the light’s existence 1 (= in ordinary consciousness) a manifold is
encountered—we have initially expressed our problem empirically 2 and we
must continue to speak this way until it has been solved—a ground for this
manifold must let itself appear in the light itself as absolute oneness and in
its manifestation, a ground that will explain this entire manifold as it occurs
empirically. “In the light and its manifestation” 3 I have said: therefore we must
first of all derive the appearance of the light from the light, [and] the mani-
fold will arise in the former. 4 This is roughly the main content of what has
been achieved so far and of what remains to be done. This is to be noted
especially:—[the task is] 5 to present appearance in general and as such.
(Obviously, as soon as appearance has been explained and the principle of the
manifold has been explained from it a priori and in principle, all appeal to
empirical experience falls away, and what was previously held factically will
be conceived genetically.)
At present we have already pushed ourselves quite near our highest
principle. If transcendental insight has been opened for us, then having in
mind the most recent link in the chain is sufficient for understanding the lec-
tures; if the earlier links by means of which we ascended to the later ones are
not equally present, nothing is lost; we will rediscover everything anew on the
descent. I must bring you back to this last link by repeating the previous lec-
ture, at the same time {bei welcher} I will also expand and add. It has already
been proven earlier that the absolute, simply as absolute, 6 must be from itself,
whatever else it may be (earlier it was “being,” “light,” and “reason,” none of
which mattered to this argument and did not belong to it); and this proof fur-
ther coheres with the postulate that inner being could not be constructed from
outside but must rather construct itself, with which postulate we opened the
This lecture begins at GA II, 8, pp. 300–301.
147148
The Science of Knowing
entire so-called “second half ” of our investigation. (In this fashion, everything
achieved so far toward bringing out the second part, and with it the whole,
could also again be reproduced.) In the last hour, this completed proof was
itself investigated in its central nerve and points of manifestness, and what
turned out to be its foundation was the simple presupposition that genuine
true seeing, or light, must have accompanied the actual Creation; 7 and since
being and seeing had already been grasped {eingesehen} earlier as being the
same, that genuine true light must itself be an immanent creation, or an
absolute “from” {Von}. This, I say, [302–303] turned up as a mere presupposi-
tion, 8 grounding the process of our proof regarding the essence of the absolute,
but itself based on nothing. Still, a little reflection 9 shows us that this presup-
position proves its correctness 10 simply by its mere possibility and facticity;
because we ourselves were the knowing, insofar as we conducted the proof and
made the fundamental assumption concerning the essence of knowing: that it
is a “from” ; 11 and, note well, we certainly could be, and indeed are satisfied, that
knowing cannot be both something in and for itself apart from any view
{Ansicht} into itself and also a “from,” but rather that it can be both only within
[such a] view. By the actuality of this view within ourselves, we have proven
directly and factically that in this respect it is so. It is, and it is this; because it
quite certainly is, and quite certainly is this, and we ourselves, the scientists of
knowing, are it as such. This is an immediate demonstration of the essence of
knowing, conducted through the fact itself and its possibility. At this point let
yourself take in even more fully what was established last time, although only
in passing: we did not make this presupposition because we wished to, or with
any sort of freedom; and if only this free element, which is to be summoned
in response to some particular reflection, is to be called We, 12 then we 13 actu-
ally did not make it; rather it made itself directly through itself. All our pre-
ceding investigations have started from the fact that we were requested to
think energetically about something we were aware of internally and also were
able to ignore; so both took place only in consciousness; this provided our
premises, and, to be sure, this energetically considered object was always
accompanied by the explicit supplement of a “should”—“If this should be so,
then ____ .” From our thinking this premise energetically, manifestness grips
us without any assistance on our part, and carries us away, attaching to the
controversial premise that conditions it and is conditioned by it. Therefore,
the knowing, which we pursued in this way, instantiated the basic character-
istic mentioned before of being merely reconstructive and was, in this recon-
struction, a secondary and merely apparent knowing, transferring knowing’s
implicitly unconditional content into a conditioned relationship. All systems
without exception remain fixed in this knowing; their premises therefore are
hypothetical for them only (but not absolutely in reason, by which even they
themselves are driven, although to be sure without their knowing it) and theTwentieth Lecture
149
relationship alone is evident, which however gives no final or fixed manifest-
ness, since the relationship itself depends on the reality of its terms. They sup-
plement the lack of this strength only by an arbitrary [304–305] reliance on
the premises and by averting their eyes from their difficulties; without this
reliance they could release themselves to absolute skepticism at any moment.
Up to this point [it has been] this way. Now, absolute manifestness has
extended itself to the premise, to the absolute presupposition itself; and
thereby it has annulled both all freedom and every “We” that was presup-
posed 14 as a premise in relation to the secondary manifestness of the context.
Hence, we are transposed into a completely different region of knowing, not
simply as something purely self-grounding {Vonsich}, but rather immediately,
and ipso facto from itself {von sich}. But as for what relates to the premise as a
premise, 15 undoubtedly in this quality a consequence is posited through it, and
these two in turn posit a relation; therefore, in this quality, it serves admirably
to explain secondary knowing; and, since it is absolute, to explain the latter
from the absolute, which is exactly our task. As a premise, it is undoubtedly
the principle of appearance we have been seeking. But since appearance is not
itself the purely self-contained absolute—as becomes evident in the premise:
the simple fact that it requires a consequence and a context shows it to be
insufficient by itself—there must therefore be a higher notion of knowing.
This remark can cast a great deal of light over what we still have left to
achieve, so I want to analyze it further.
Now we let go of the point that we presuppose it, or more accurately,
that it posits itself as a presupposition, and thus let go of the proposition’s
form for reasons having to do with method, and simply hold on to the con-
tent of the proposition: “the light is an absolute “from””—analyzing what we
actually mean by this.
1. All along, and obviously in this proposition, light 16 is posited first and
foremost as an essential, qualitative, and material oneness, not further conceiv-
able, but instead only to be carried out at once, just the way we carry it out in
all of our knowing, from which we cannot escape. I want to be understood on
this point that is easy in itself, that requires simple, strong attention. Question:
what then is knowing? If you know, then you just know. 17 [306–307] You can-
not know knowing again in its qualitative absoluteness; since if you did know
it, and even now were knowing it, then for you the absolute would not stay in
the knowing that you knew about, but rather in the knowing by which you
knew it; and it would go on this way for you even if you repeated the procedure
a thousand times. It remains forever the same, that in absolute knowing you
recapitulate knowing as essential qualitative oneness. Initially this insight needs
only to be carried out; reflecting on the law of its completion still remains before
us. This light is now absolutely presupposed as a “from” without prejudice to its
qualitative oneness; since [if it did not preserve this oneness] 18 the light would150
The Science of Knowing
not be “from,” and so [it is presupposed] as permeating this “from.” Notice first
what is new and important here: it is presupposed 19 to be like this, uncondition-
ally. Thus it has presupposed itself in a particular act; and this presupposition
is now proven 20 by its facticity and possibility, and further by the possibility of
a deeper determination of knowing, which opposes the simple existential form
in its attachment to the mere dead “is.” 21
In our insight, it in no way follows 22 from our insight 23 into the essential
light as such (which we once again should have been grasping energetically
and freely), by means of which manifestness descended into the connection
between light as such and the “from.” And we again lapsed into secondary and
merely apparent knowing with which there must, to be sure, eventually be an
end, an end we have sought so avidly from the beginning. “It does not follow,”
I said, since, as we have seen, there is generally no such insight into the light
in itself, hence, I said: the light posits itself as a “from” in a particular and
absolute act or genesis; an act that cannot be mastered immediately in this
genesis, as the genesis of genesis, because otherwise the genesis would not be
an absolute genesis. (What this latter means, and does not mean, since here as
well is yet another disjunction, will show itself it what follows.) I say: accord-
ing to the preceding observation, it thus posits 24 itself absolutely; the act is a
self-contained, self-sufficient act; it is posited by us merely in our inferential
chain, which we now entirely let go of, as the mere means by which we have
ascended to our present insight, until we find it again on the descent. This is
the first, and significant, point.
That the light in its changeless qualitative oneness is a “from” therefore
means: it is a qualitatively changeless permeation of the “from.” 25 In the pre-
vious hour, we made the following application of this point: disjunction is
found everywhere in the “from”; absolutely out of, and from, the “from”; 26 by
no means [308–309] presupposing terms that were primordially different
independently of the “from,” instead [they were] produced absolutely as terms,
absolutely distinguished as such only through the “from” and otherwise
through nothing at all. The one, eternally qualitatively self-identical light, by
virtue of its identity with the “from,” must, in this qualitative oneness, spread
itself over these terms, whatever their distinction from each other.
Let me now apply and animate this insight right away, and thereby make
it unforgettable for you. A “from” is posited immediately through the light:
L
a — b
Hence, if light exists, then necessarily there is also a “from.” 27 Now if the
light is identical with the “from,” then, as surely as it itself exists, it spreads
itself in unchanged qualitative oneness across every “from,” and comprehendsTwentieth Lecture
151
every “from.” And if one again posits within the first “from” another one that
is deducible and conceivable on the basis of the original synthesis of light and
the “from” 28
a — b
|
|
a–b a.b
it is completely clear that the same original light, qualitatively unchanged, by
means of its identity with the original “from,” must at the same stroke accom-
pany all subordinate divisions of the original “from” 29 into further “from’s.”
And it is also clear that whatever possesses the principle of this secondary
splitting of the original “from” accompanies this progression of the light as
entirely necessary and at one stroke, and can reconstruct it purely a priori and
without any empirical presuppositions; which indeed is the second and subor-
dinate task of the science of knowing, since we now are pursuing the much
higher task of presenting the principle of this principle. This “from,” in pure,
absolute, immediate oneness 30 and without any disjunction, as the pure self-
positing of the original light, is the light’s first and absolute creation; the
ground and original source even of the is, and of everything that exists; and the
disjunction within this “from,” in which true living perishes and is reduced to
the mere intuition of a dead being, is the second re-creation in intuition, that
is, in the already divided original light. And thus the science of knowing jus-
tifiably presents itself as the complete resolution of the puzzle of the world
and of consciousness.
This, I asserted, was the next application that I made last hour of the
proposition, “The light is a ‘from,’” attending to the disjunction in the “from.”
But it is even more important to look at the essential and qualitative oneness
of this “from” and at the words that were said previously about the original
creation. Recall them. In its pure qualitative oneness, “from” is genesis: that the
light is identical with it and permeates it in this its essence, means: in this its
second power (namely its appearance), it is itself genesis; [310–311] genesis
and seeing converge together completely and unconditionally.—The words
are easy to understand; it is not so easy to give them the deep meaning
intended here in living insight; and it is nearly true that the only way I am able
to guide you forward is with an example. The subject matter that I wish to
present to your intuition appears in every transition from lassitude to energy,
and for our purposes, the example cited just above will serve best: the one in
which we had tacitly presupposed absolute knowing to be a “from”; when
interrogated about the justification of this presupposition, we recalled that
indeed we knew ouselves in this presupposition and were the knowing. I ask:
does not this new awareness, that was not yet there prior to our presupposing,152
The Science of Knowing
seem as if it were a popping up, a new production? Now at this point you are
certainly able to abstract purely, as is my present demand: that this is con-
sciousness, is a consciousness of knowing, and what is more of knowing as a
“from.” What remains for you after this [act of ] abstraction? Evidently just a
knowing/seeing/light, exactly absolute, qualitative, as it has already been
described, and [it was] this therefore because you abstracted completely from
all content, which you could do according to the presupposition; consequently,
[it], as itself light, conducted the proof of legitimacy empirically; further, [it is]
a consciousness of absolute genesis. Now—note this addition, the proof becomes
more rigorous and the insight purer through it—you can more suitably posit
this genesis or freedom in the act of abstraction from all content of the pre-
sented consciousness, which is thus required of you. As things stand, it is
immediately clear to you that this pure light, as it has been described, could not
arise without abstracting from all content, 31 nor can the latter appear without
arriving at pure light; that therefore the appearances of both terms are indivis-
ible, and permeate one another; and that hence pure light appears as permeat-
ing genesis, or as producing itself. By means of this proof more is nearly proved
than should be proved, and future research is anticipated, as I note in prepara-
tion; the light’s positing of the “from,” and the fact that it posits itself as a
“from,” has already become immediately visible. What we have to be concerned
with next here can be shown with a little preparation in two examples. Because
you were instructed to reflect energetically and a new consciousness emerged
for you, this new consciousness is not to exist {soll . . . nicht sein} as something
new without the energy; this consciousness [312–313] and the energy should
open up together indivisibly. Now you certainly posit genesis here partly in
yourself, in the energy of your reflection, and partly in the essence of reason
itself, since the manifestness is to emerge without any further action on your
part; but this entire distinction ought to have no validity in itself, and it should
be abstracted out, and so—leaving undecided whether genesis’s true principle
lies in me or in reason itself—there always remains an absolute, self-producing
knowing that does not possibly occur without the genesis.
Now this means, as was said before, that light permeates the “from” in
the qualitative oneness of its (the “from’s”) essence: the presented intuitions of
this penetration were only explanatory means. But, independent of all factic-
ity, we have seen a priori that if light is to be, such a permeation is necessary.
This is the one side of the previous proof for the content of the sentence:
the light = “from,” that we have repeated and enriched today. There is still
another, and of this more tomorrow, equipped with today’s new discoveries!
In conclusion another comment about the whole of the science of
knowing, one that I share with you not so much for your own guidance, since
I hope you do not need it, but rather as a weapon of defense against the igno-
rant. Already earlier, and again today in passing, the proof of knowing’s essen-Twentieth Lecture
153
tial criteria is conducted on the basis of our capacity to see it thus. The nerve
of the proof is clear: we ourselves are knowing; since we can know only in this
way, and presently actually know thus, then knowing is constituted so. It is
equally clear that the failure to discover this principle of proof, or not paying
attention to it once it has been found, grounds itself on the truly foolish
maxim of searching for knowing outside of knowing. Concerning this, noth-
ing more needs to be said. I would only bring this to your attention: the proof
simply does not succeed for anyone who is really not able to make clear and
intuitable what can only be made so by his own capacities; through his incom-
petence he is barred from the subject itself and from any judgment about this
world that is entirely concealed from him. It is the same for those who could
but will not, that is, who will not submit themselves to preliminary conditions
of sharp thinking and strong attention; because everyone who can, will do the
thing itself; and everyone who will, can do it. This is true when the science of
knowing does not yet stand at the apex. 32 One should not therefore wonder
how that which has in itself the highest clarity and manifestness, cannot in
any way be made clear and true for very many people {Subjekten}; one can
rather himself lay out the grounds for this impossibility, if they will just come
to understand the premise that there might be something they do not now
know; and that they are not able to know directly and without much prepara-
tion and strong discipline, as things are with them now.C H A P T E R
21
Twenty-first Lecture
Thursday, May 24, 1804
Honored Guests:
(We will make use at once of what we have already understood, and take a
shortcut without further repetition and closer definition of subordinate
terms. 1 You know that such a thing is possible in the science of knowing, and
why. That is, [it is possible] because the subordinate terms will recur in their
full developmental clarity during the descent; and the ascent is undertaken not
for the subject matter itself, but for clearing our vision and opening it to the
absolute by abstracting from all relations. 2 )
I connect this with what has gone before: the light has been presupposed
as an absolute “from.” Then we immediately proved the legitimacy of this pre-
supposition by means of its bare possibility and facticity, because we ourselves
were light and knowing. Based on this last key step in the proof, 3 the presup-
position is true and legitimate in the “We” ; not, of course, in the previous “We”
that freely posited premises, (since in this case knowing posits itself, as was
clearly explained yesterday). Instead, [it is true and legitimate] in the We that
merges into the light, and is identical with it. Moreover, it truly is just as it
factically occurs, but it occurs as a presupposition. Hence, taken strictly (as we
have not so far taken it, and for good reasons) 4 it has been truly and factically
proven that the light can presuppose itself as a “from,” and that in us it actu-
ally does so. In us, to the extent that we have merged, and disappear identi-
cally, into the light itself = [we] are the science of knowing. Unnoticed, this
presupposition has made itself, and we will build on that. But [the We] 5 on the
occasion of which it made itself, has in that sense not even made itself, instead
we, who are freely abstracting and reflecting, have made it. Consequently, by
this “We” one may well mean that light makes itself into a “from” only in the
science of knowing, as a higher, absolute knowing; and so we provisionally
indicate a distinguishing ground (for which we have been searching) between
lower, ordinary, empirical knowing and higher, scientific, genetic knowing. 6
This lecture begins at GA II, 8, pp. 316–317.
154Twenty-first Lecture
155
We said, “It is presupposed.” 7 However, all presuppositions bring along a
hypothetical “should”; and let themselves be expressed through it. In fact, we
have not argued differently than this in the two previous sessions when ana-
lyzing the contents of this “from”: “Is there light” = “if light is to be”—and “Is
there an absolute ‘from’” = “If there should be an absolute ‘from,’ then must
____ ,” etc. However, we have not only presupposed the absolute oneness of
the light hypothetically; [318–319] instead we have also realized it uncondi-
tionally. To be sure, we have done so only in its qualitative character 8 —which
as you will remember, was itself a result of the “from”—as was the case with
the absolute origin in knowing as permeating the “from” in its qualitative one-
ness, as we discovered yesterday. Hence, both are a result of the hypothetical-
ness, so that only pure, bare oneness, henceforth presented as inconceivable
and understood as categorical, remains left over. I wished to undertake the
delineation of this very boundary in passing, and it is commended to you.—
Now back. Our reasoning has proceeded in the hypothetical form of a
“should”; and this to be sure unconditionally as itself knowing, and as primor-
dial knowing, since knowing itself has posited this “from,” then transcended
this posited and objectified “from,” which we analyze from below 9 and derive
from it.
(This as well about method. Obviously we are once again reflecting
about what we were and did in the previous presupposition and analysis, in the
same way we have proceeded in our entire ascent; and I could have proclaimed
our activities in just this form. Purely because we have left the realm of arbi-
trary freedom behind and have arrived with our own effort in the realm of
organic law, I preferred to compel you to the present reflection through the
reminder that indeed everything grounds itself simply on the presupposition,
rather than appealing to your freedom.)
2. 10 In its innermost essence, 11 a “should” is itself genesis and demands a
genesis. This is easily understood; you ask, If such and such should be, then is
it or isn’t it? 12 The “should” tells you nothing about this. What then does it say?
It sees 13 a principle; therefore, it explains categorically that being can be
admitted only on the condition of the principle. Thus, only genetic being, or
being’s genesis, can be admitted. Thus, it is the absolute postulation of gene-
sis; and since everyone whose transcendental sense has been awakened will
allow no genesis to be valid in and for itself without such a postulation, even 14
immediately absolute genesis, and the genesis of objective genesis only medi-
ately, according to a law that we have yet to exhibit. Or this as reinforcement:
I have said [the “should”] is the postulation of genesis. Now it is immediately
clear that the “should” is a postulation, and that a postulation is a genesis, at
least an ideal one; otherwise, it is, as such, completely incomprehensible and
accordingly the addition “of genesis” would not be worth while in any way. So
it is evident that, in our hypothetical “should,” being’s 15 genesis is demanded,156
The Science of Knowing
which, as a genesis of being, the hypothetical “should” is content not to be able
to provide. Instead, it waits for it from a principle outside itself. The demand,
as itself a genesis (ideal [320–321], as we have called it, in order to name it
only provisionally with this partially clear term), however, lies in the “should,”
and the “should” is it. Thus, there can be a disjunction within the absolute
genesis itself, through which it would be real and ideal. This entire disjunction,
discovering the basis of which may well be our most important task, can now 16
follow from genesis, or the “should.” Through resolving this disjunction, those
words, which we have used so far only provisionally and according to a dim
instinct in the hope of an eventual clarification, will themselves become clear. 17
This is merely a hint at a part of our system that necessarily must remain
obscure here.
However, the following is completely clear in what has been said: in
virtue of yesterday’s demonstration, genesis = the “from” in its qualitative one-
ness. We ourselves, or knowing and light as such, which are entirely the same
as us on the level of our present speculations, are this “from” immediately, in
that which we pursue and live. So there is no further need for the “from” that
is posited and presupposed through some specific act of ours or of the light,
nor for anything that we have derived from it in our analysis. Therefore, we
let it all go as just a means of ascent, until it shows up again on our descent. I
said, “in that which we pursue and live”; and this very pursuing and living, as
pursuing and living, follow directly from [our] dissolution into genesis.
3. By virtue of the hypothetical “should,” the we, or knowing, is
absolutely genetic in relation to itself. Because we ourselves were knowing, we
pursued it in the following way: “should knowing be (that is, should we our-
selves be, since we ourselves are knowing), then ____ must” and so forth.
Thus, [it is the] genesis of nothing else, but rather of itself, of the simultane-
ously productive [one]. Thus, with this it is absolute genesis, which carries in
itself the already sufficiently seen {eingesehen} character of being or light: that
it is completely self-enclosed and can never go outside itself.
4. This absolute self-enclosure of genesis in its fundamental point (in
which it should be a genesis of genesis) does not prevent two points of origin
{Genesen} or two knowings from appearing impermanently. We ourselves con-
duct one when we say, “If knowing (or we ourselves) should be ____”; and the
other one should be, if its principle is fulfilled. 18
I regard insight into this distinction of two aspects of knowing, a dis-
tinction that is still only factical, as simple. Yet, it is so important that I can-
not very well leave it to mere luck, and so a bit more by way of elucidation.
We ourselves are the absolute light, the absolute light is us, and this is gene-
sis itself. Nothing can depart from this; therefore, a distinction [322–323]
cannot be admitted within the subject matter itself, without contradicting
our first fundamental insight. Hence, the disjunction that remains is not aTwenty-first Lecture
157
disjunction between two fundamentally distinct terms, instead it is a disjunc-
tion within one, which remains one throughout all disjunctions. Something
of this sort has already presented itself to us earlier. Stated popularly: it is not
a disjunction of two things, but rather just different aspects of one and the
same thing.
5. Letting this disjunction stand provisionally, just as it has appeared to
us factically, with the intention of working further on the basis of it, the ques-
tion arises, which of the two aspects is to be considered provisionally as
absolute in order to explain the other from it? [It is] obviously the first term,
in what we ourselves live and pursue, given the insight aroused in us yester-
day, that seeing and light reside always only in immediate seeing itself and
never in the seeing that is seen. By no means could it be in the objectified is,
that waits [to receive] its being from a principle and is therefore truly dead
within. This choice can be shown to be completely necessary through another
circumstance as well. For if we wish to work further, then we wish to pursue
and live knowing further as well. Therefore, in fact and absolutely, we must
remain in life and cannot abstract from it, as is evident. To do so would exactly
be not to live and search further, but instead to remain here, which would con-
tradict our intention not to stand still and instead to go further. 19
(In passing, this aspect is the one that we have always called idealistic.
Thus, our science, standing between idealistic and realistic principles, would
at last become idealistic, and indeed, as we have seen, be forced to do so by
necessity, and contrary to its persistent preference for realism. We will not
promise that the matter will come to rest with this principle, as it now stands
and as it will at once be explained more clearly. We can promise more confi-
dently that we will never again use objectivity as a principle. From this it will
follow that, if the idealistic principle too should prove inadequate, we will
need to find a third, higher principle that unites the two.) 20
6. Just as is demanded in this principle {Satze} of absolute idealism, the
inner self-genesis 21 is presupposed as a living inward oneness—what this is, on
that point you understand me—as oneness, thus as light, qualitatively
absolute, only as something to be enacted and by no means as something to
be understood. This latter [oneness is to be presupposed] as genesis (i.e., as was
made obvious yesterday in each transition from dull to energetic thinking), as
disappearing into the arising of an absolute “from” 22 that in turn merges into
it, so that seeing and this [324–325] 23 arising are entirely inseparable—that is,
as genesis of self, or I, so that accordingly what emerges in immediate light
may be an “I”—as a result of which, light and this very We, or I, would merge
purely into one another. The principle demands this: inner self-genesis is to
be presupposed as intrinsically living oneness, then also knowing’s objective
aspect is allowed to stand and to be united with the previous [aspect], as it can
be united in knowing alone. From the genetic principle, it would then follow158
The Science of Knowing
that a principle must be assumed for the absolute, inner, and living self-gene-
sis, and that this latter [event must occur] in a higher knowing that united
both, as is obvious. This latter knowing is then the highest, and the two sub-
ordinate terms {Seitenglieder} are merely what is mediated through it.
That in the higher knowing a principle is presupposed for absolute self-
genesis means that inwardly and materially this higher knowing is non-self-
genesis. 24 Yet it does not not exist, rather it exists actually and in fact; thus [it
is] positive non-self-genesis; and yet it is immanent 25 and is itself an I, because
this is its imperishable character, as absolute. What else is negated besides
genesis? Nothing, and to be sure this is negated positively; but the positive
negation of genesis is an enduring 26 being. Thus, knowing’s absolute, objective,
and presupposed being becomes evident in this higher knowing, hence directly
genetic, as it has previously appeared merely factically.—Once again, in order
to review the terms of the proof: as a result of positing a principle for self-gen-
esis within knowing itself, [we derive] the explanation of genesis as not
absolute; consequently, [we infer] its positive negation within knowing; and
consequently [we also infer] the positing within knowing of knowing’s
absolute being.
If you have just grasped this rigorously, then I can add something else
in clarification. 27 This knowing, that only ought to exist, is of course a self-
genesis of knowing, its self-projection beyond itself, as we, who are standing
over it reconstructing the process and its laws, very well understand {einsehen}.
However, the question still always remains as to just how we arrive at this
insight and so apparently get outside of knowing. Yet, in contrast to an
absolute self-genesis, which is itself annulled as absolute by the addition of a
principle, the immanent knowing that never can get outside itself for just that
reason can never appear as self-genesis but only as the negation of all genesis.
Here, therefore, there is a necessary gap in continuity of genesis, and a pro-
jection per hiatum—but here presumably not [326–327] an irrational one.
Rather, it is [a projection] which separates reason in its pure oneness from all
appearance, and annuls the reality of appearance in comparison with it.
“Reason,” I say, in order to clarify this for us; in this case we were only
concerned with deriving the form of pure being and persistence. In our case,
this persistence is now, and certainly always and eternally, genesis. This exist-
ing knowing—which to that extent, is not genetic as regards its external
form—is enclosed in itself in unchangeable oneness, and so indeed [is] also
genesis, just as it seems to be above. Thereby the absolute, inward awareness
declares itself, without any external perceiving, knowing, or intuiting, all of
which fall out in self-genesis—[an awareness] of an original principle and an
original principled thing {Urprincipiat} in a one-sided, and certainly not reci-
procal, order; or pure reason, a priori, independent of all genesis, and negating
it 28 as something absolute.Twenty-first Lecture
159
Let us go further: in what we most recently lived and pursued, we our-
selves have not become pure reason itself, nor dissolved into it, instead we
have merely deduced it from 29 an insight into it. However, this was possible
only to the extent that we presupposed self-construction as absolute, as we
did; because what should follow, follows only on the condition that it [that is,
self-construction] is annulled as absolute in itself, from itself, and through
itself; and this was the center {Nerv} of our proof. Since it was mentioned pre-
viously that the higher knowing, which projects reason, is also at bottom self-
genesis and consequently does not just appear to be, we can very appropriately
call this self-genesis the reconstruction of the non-appearing original genesis,
thus the clarification of the terms of the original genesis, hence [we can call
it] the understanding.—Accordingly, it follows for us that there is no insight
into the essence of reason without presupposing understanding as absolute;
conversely, [there is] no insight into the essence of understanding except by
means of its absolute negation through reason. However, the highest, in which
we [328–329] remain, is the insight into both, and this necessarily posits both,
although [it posits] the one in order to negate it. From this standpoint, we are
the understanding of reason, and the reason of understanding, and thus both
in oneness. Now the disjunction stands forth in its clearest definition. Just one
more principle and the matter will be completely explained. [We will talk]
about this, next Monday.
This besides; I regard what I have just presented to you as not at all easy.
However, that lies in the subject matter, and we have to go through it some-
time, if we want to see solid ground. I can promise you a bit more illumina-
tion on this from an insight into the principle we are still seeking, but then the
difficulty will lie in the principle itself.
One cannot speak properly in front of others about speculation in these
heights freely and without preparation, since one has enough work speaking
of it in formal, prepared lectures. For this reason, and in order to escape our
mutual impulse nevertheless to handle this matter freely, [we will now have]
a special discussion period.C H A P T E R
22
Twenty-second Lecture
Monday, May 28, 1804
Honored Guests:
Although I can justifiably report that our speculations now already hover at a
height not reached previously, and have introduced insights that fundamen-
tally change the view of all being and knowing (and I hope that all of you who
have historical knowledge of philosophy’s condition up to now will agree),
nevertheless, all of this is still only preparation for really resolving speculation’s
task. We intend to complete this resolution during the current week. Hence,
your entire attention is claimed again anew. Whoever has completely under-
stood everything so far, and seen into [it] to the level of eternally ineradica-
ble, forever immovable conviction, but who has not yet seen into, and achieved
conviction about what is to be presented now, such a person has achieved at
least this protection from all false philosophy: he can set each of them straight
fundamentally. He also possesses some significant truths, disconnected and
separated from one another; but he has not yet become able to construct
within himself the system of truths as a whole and out of a single piece. I now
intend to impart this capacity to you, and after that the main purpose of these
lectures on the pure science of knowing will have been achieved.
Whether one names the absolute “being” or “light,” it has already been
completely familiar for several weeks. Since attaining this familiarity, we are
working on deriving not, as is obvious, the thing itself, but rather its appear-
ance. The request for this derivation can mean nothing else than that some-
thing still undiscovered remains in the absolute itself, through which it
coheres with its appearance.
We know from the foregoing (which, to be sure has been discovered
only factically, but which nonetheless would have its application in a purely
genetic derivation) that the principle of appearance is a principle of disjunc-
tion within the aforementioned undivided oneness and at the same time, obvi-
ously, within appearance. However, as regards the absolute disjunction, I urge
This lecture begins at GA II, 8, pp. 330–331.
160Twenty-second Lecture
161
you to recall an analysis conducted right at the beginning of this lecture series,
in which the following became evident. If disjunction were to be found
directly within absolute oneness, as is unavoidably required by the final form
of the science of knowing, and is what we intend here, it must not be grasped
as a simple disjunction, but rather as the disjunction of two different disjunc-
tive foundations. [It must be] not just a division, but rather the self-intersect-
ing division of a presupposed division that again presupposes itself. Or
[332–333], using the expression with which we have designated it in our most
recent mention of it, [it is] no simple “from,” but rather a “from” in a “from,”
or a “from” of a “from.” The most difficult part of the philosophical art is
avoiding confusion about this intersection, and distinguishing that which is
endlessly similar and is distinguishable only through the subtlest 1 mental dis-
tinguishing. I remind you of this so that you will not become mistrustful if, in
what follows, we enter regions in which you no longer understand the method
and [in which] it should even seem miraculous. Afterwards we will give an
account of it; but beforehand we actually cannot.
This much as a general introduction for the week:—Now back to the
point at which we stood at the end of the last hour.—Absolute self–genesis
posited and given a principle, obviously within knowing, which is thus a “prin-
ciple-providing”[occurrence]. 2 Thus within this knowing there follows the
absolute, positive negation of genesis = completed and enduring being; and
indeed, because this entire investigation concerns light’s pure immanence, our
investigation has long ago brought in {über die Seite} a presumed being exter-
nal to knowing, knowing’s [own] completed and enduring being.
Now We saw into this connection during the last hour and see into it
again here; as is evident, we see one of the two terms in and through the
insight into this connection, determined as such. Thus, this latter is itself
mediated, and just we (= the insight we have now achieved) are therefore the
unconditionally immediate [term].
Two remarks about this. 1. I have just recalled again, that here the inner
being and persevering is knowing’s being, the very thing already recognized as
absolute genesis and [the thing] we have also already validated in the last hour
as a priori rational knowledge of an absolute principle. At this point, we must
hold on to the fact that it is knowing’s being, even in case we should let go of
the addition as a shortcut in speaking; because otherwise we will fall back into
where we were before, far removed from further progress. Therefore, our
entire chain of reasoning must always be present to us, now more than ever. 3
2. It is said that every philosophical system remains stuck somewhere in dead
being and enduring. If now a system derived this being itself in its inner
essence, as ours has done, by positing a higher principle for absolute genesis,
whereby it then necessarily becomes the positive negation of genesis, and
therefore [becomes] being; if too this being is not the being of an object and162
The Science of Knowing
so doubly dead, but rather the being of knowing, 4 and so of inner life; then
such a system seems [334–335] to have already accomplished something
unheard-of. However, we are required to see {einzusehen} clearly here that by
this we have not yet achieved anything, since even this saturated being of liv-
ing is again something mediate and is derived from that which alone now
remains for us, the insight into the connection. Now let us apply this directly
for our true purpose, which we have long recognized. Knowing’s derived being
will now yield ordinary, non-transcendental, knowing. Through our present
insight into the genesis of the former’s principle (that is, the principle of the
recently derived knowing), and through reflecting on this insight we elevate
ourselves to genuine transcendental knowing or the science of knowing; and
[we have done so] not merely factically, with 5 our factical selves, so that we are
the factical root. We have already been this since the time when we dissolved
into pure light; instead [we have done so] objectively and intelligibly, so that
we, achieving insight factically, at the same time penetrate the law of this
insight. Henceforth we have to work in the higher region that has now been
opened. Only here will the principle of appearance and disjunction that we
seek show itself to us; which then should only be applied to existing (= ordi-
nary) actual knowing. Now also this addition: since the beginning of what we
have provisionally called the second half, a hypothetical “should” has been evi-
dent as simply creating a connection, and as linking a conditioning and con-
ditioned term (which are both produced absolutely from it) to a knowing,
which must be originally present independently of the “should” and its entire
operation, if one just understands it correctly. This could be called the first
sub-section 6 of the second part. Since we concerned ourselves with an absolute
presupposition about the essence of knowing as an absolute “from,” we wished
to know nothing more about this entire hypothetical “should” and its power
of uniting and joining, as a merely apparent knowing. We said at this point
that so far we (the we that still has not been grasped so far) have concurred
about the premise’s arbitrary positing, pointed out by energetic reflection, and
only the connection 7 has made itself manifest without our assistance. Now the
premise too presents itself without our assistance; therefore in the premise too
we coincide with the absolutely self-active light. Let’s hold on to this. We have
been doing this for a long time in our discussions about the “from,” until I
thought you were sufficiently prepared for the higher flight that we began in
the last hour. You may take this as the second sub-section of the second
part.—In the last hour, the bare connection 8 presented itself again, and (as we
might suspect, but will more exactly show and demonstrate) [336–337] with
it the hypothetical “should,” from which we had already hoped to be free. This
should surprise us. If this “should” has reappeared with the same significance
in which it was previously struck down, then we have not advanced and are
just sailing on in the seas of speculation without a compass. Through the hintTwenty-second Lecture
163
given recently as to the difference between ordinary knowing (based on the
principle of knowing’s being) and transcendental [knowing] (based in the
genetic insight into this same principle), it is probable that [the “should”] does
not arise in the same way. Instead, the “should” that we have dropped operates
in ordinary knowing with tacitly assumed premises; in contrast, the “should”
arising now operates in transcendental 9 knowing, which grounds its premises
genetically—as emanating from a “should.” Therefore, in the preceding lec-
ture, we have begun a third sub-section of the second part, and the two
extreme sections come together in the middle (with the premise) once again
distinguished by a duality in the premise. By this means, the two outermost
parts (= transcendental and actually existent knowing) would be the two dif-
ferent distinguishing grounds, proceeding from the middle ground of the
premise, which both unites and separates them. This is the compass that I
would share with you for the journey we have already begun.
2. 10 I said that a hypothetical “should” appears again in our completed
insight. To begin with, this is obvious. “Given [that there is] a principle for
self-genesis, then it follows that ____ .” 11 Previously, we have found the two
terms factically, and to that extent separately; but in the last hour, we have
united them genetically, according to our basic rules and maxims. Because we
comprehend them mediately in the insight into their connection, then, given
this insight, we no longer need to assume them factically. They reside a priori
in the insight, and we can drop the empirical construction, until perhaps it
arises again in some deduction.
3. Now let us grasp this hypothetical character at its core. According to
the maxims and rules that we arbitrarily adopted at the beginning of our entire
science, and hence arbitrarily, we appear to ourselves—so it has been, and so
it now is openly admitted—as genetically uniting both terms. Here is the
inner root of hypotheticalness (now fully abstracted from the hypothetical
character of the subordinate terms), precisely the “should’s” admitted inner
production, containment, and support of itself as identical with the [338–339]
free We, i.e., the science of knowing. This very inward hypotheticalness could
be what first shows itself and breaks through in the subordinate terms’ hypo-
thetical character. Hence, it is only a matter of negating this inner hypotheti-
calness, so that thereby the quality of being categorical will be manifest in it,
and thereby we will justify our insight in its truth, necessity, and absolute pri-
ority. In that case, the inference, made 12 here only provisionally, first achieves
categorical validity; namely, the inference that both terms (i.e., knowing’s self-
genesis and its being) occur only mediately in a genetic insight into the one-
ness of both, and in no wise immediately.
(A remark 13 belonging to method: One must not allow oneself to be dis-
tracted [as to] how I legitimately assume this. This is more necessary than
ever, since here the method itself becomes absolutely creative; further, nothing164
The Science of Knowing
besides a remark like this can be adduced in explanation of what happens here.
Our insight has come about through the application of the basic maxims of
our science, to apply the principle of genesis thoroughly and without excep-
tion. This has been required, and it should prove itself. Likewise, the maxims
of the science of knowing itself, and with it all of science, have been required,
and they should prove themselves. Science itself should justify and prove itself
before it truly begins. Thereby, the science of knowing would be liberated
from freedom, arbitrariness, and accident; as it must be, or else one could
never come to it.)
4. Without digression, we conduct the required proof, according to a law
that has already been applied, thus. We could produce this insight, and we
actually did so; we are knowing; 14 thus this insight is possible in knowing and
is actual in our current knowing.—Just a few remarks about this proof. a. The
genesis first accomplished by us is an absolute, self-enclosed, genesis; by no
means is it a genesis of a genesis, 15 because it negates itself inwardly within know-
ing, as we have shown in the last hour. To us, however, who are contemplating
further and constructing the process in its laws, it manifested as genesis. In
immediate knowing, however, it was merely a persisting intuition, as external
in its result: which explicitly was non-genesis 16 = being. b. The proof of absolute
genesis was conducted purely through its possibility and facticity, and thus [is]
itself only immediately factical. In this case, therefore, facticity and genesis
entirely coincide. Knowing’s immediate facticity is absolute genesis; and the
absolute genesis is—exists as a mere fact—without any possible further ground.
[340–341] To be sure, it must happen so, if we are ever actually to arrive at the
ground. c. This is a cogent example of how much in the science of knowing
depends on one always having the whole context present, since the distinctions
can be drawn provisionally only through this context. Knowing, as genesis, is
proved factically in this way. What then happened several sessions ago when
we proved knowing as a “from” factically in the same way? Is the “from” some-
thing other than genesis, and have we not conducted the identical proof? Yet
the present, factically demonstrated, genesis is something entirely different
from the one proved earlier. You could grasp this distinction {Character} only
by noticing that in this case it is a question of the genesis of absolute knowing
in its fundamental construction, whereas previously [it was a question of ] the
genesis of its absolute self-genesis. In that lecture, to be sure, I had to let this
criterion go and abstract from it, looking simply at the core element of the new
proof (since otherwise we would never arrive at this new proof ) and relying on
the solid insight already engendered in you. However, you can add this crite-
rion now, and can use it to rebuilt and reinforce the insight, in case it begins to
waver. Later of course I will add inner distinguishing criteria, e.g., for both
these points of origin {Genesen}, by which they can be distinguished in them-
selves, independent of their relation. However, they are not even possible, orTwenty-second Lecture
165
comprehensible, before the distinction in the thread of relatedness has been
completed factically, because these inner distinctions are nothing but the
genetic law of the factical difference, which arises only in the fact. For just this
reason, the science of knowing is not a lesson to be learned by heart, but rather
an art. Presenting it, too, is not without art. 17
d. I also want to make you aware of the following. Even in the recently
conducted proof, within whose content facticity and genesis should merge
purely into each other, there still remains in the minor premise of the syllo-
gism the same term that arose above factically and still has not been geneti-
cally mastered. 18 “We know, or are knowing.” [This is] of course immediately
clear and intelligible, but its principle is by no means clear. Investigations
remain to be made here, and herein lies perhaps the most important part of
our remaining solution.
Let me announce the process that will follow. For good reasons lying in
my art, I will not proceed at once with this point, but instead I will wait until
it arises of itself. On the contrary, I will add this also: in the insight we have
completed, an insight has now indeed arisen for us that is objective, immedi-
ately compelling for us, as well as intrinsically determinate and clear. It is this,
if [342–343] we abstract from the subordinate terms as hypothetical and as
themselves just the externalization of the inner hypotheticalness of our per-
formance (of course, we need to abstract from this hypotheticalness as well),
then we intend to turn from the form of our insight to its content, in order to
explain the form from it, as we have frequently done and as is actually a real-
istic move {Wendung}.
To repeat briefly: the content of the insight we have recently achieved
must be clearly present to you. Providing absolute self-genesis with a princi-
ple yields absolute non-genesis = being. As we just recently undertook to do,
today we have abstracted completely from the two subordinate terms. And,
looking only at our insight itself and the manner of its production, we have
justified, as was possible only in a factical manner, the absolute application of
both the maxim of self-genesis and this procedure. This is the essential con-
tent of the little, easily remembered, bit that we have achieved for our topic.
Further, at many turning points {Wendungen} we have made extremely
penetrating remarks about method, which I urge you to keep in mind, because
only with them will you find your way through the maze which we confront.C H A P T E R
23
Twenty-third Lecture
Wednesday, May 30, 1804
Honored Guests:
Providing a principle for absolute self-genesis, as we have described the light,
creates non-genesis, or being (of knowing, it goes without saying). We have
seen {eingesehen} this, have reflected about the method for producing the
insight, and have justified it factically. It is possible and actual within know-
ing, because it is possible and actual within us, and we are knowing. However,
no genetic derivation of this last point has occurred just yet.
We can justify this procedure even more deeply and from another point
of view. An objective, and absolutely compelling insight has actually arisen for
us; this process hereby has also shown itself to be coherent with the absolute,
self-producing light. Consequently, the task which we reported at the end of
the last hour as coming next (to investigate this objective insight in respect to its
true content) 1 will simultaneously be our first attempt at justifying more deeply
what has up to now been presented only factically, and perhaps even to make
the point genetic.
—And so to the content of the self-presenting objective insight! 2
1. Evidently, the absolute relation of both subordinate terms. However,
these are still hypothetical; but without them, there is no connection; it itself
is the result of hypothetical terms, and so itself is hypothetical; a fact from
which we should abstract. What is still left? Plainly nothing other than the
insight’s inner certainty; and because even the insight as such depends on the
terms, [it is] nothing more than a purely inner certainty.
The first claim on you is to grasp this certainty sharply and altogether
purely. It is not certainty about anything in particular, as the relation of the sub-
ordinate terms was in our case, because we have abstracted from that. Rather,
it is certainty pure and as such, with complete abstraction from everything.
At first, accordingly, it is immediately clear that certainty would have to
be thought completely purely. Moreover, the ground of its material what-ness
This lecture begins at GA II, 8, pp. 344–345.
166Twenty-third Lecture
167
{Washeit} lies in the “what” 3 (that it is, that it is what it is, that it is certain 4 );
but the ground for certainty cannot in any way be located there, because this
does not belong to the “what.” Therefore, certainty lies simply in itself, uncon-
ditionally and purely as such; and it is unconditionally of and through itself.
It must be thought in this way, or else certainty is not conceived as certainty.
(In passing, being is not a reality that is derived from the sum of all
possible realities [346–347] (i.e., from the possible determinations of a
“what”); rather, it is completely closed into itself, and outwardly, in its prop-
erties, is first the condition and support of every “what.” Kant was the first to
validate this proposition—misunderstood almost entirely by the old philoso-
phy. As regards the first half of the sentence: being is something living from
itself, out of itself, and through itself, absolutely self-enclosed and never com-
ing out of itself, as we grasped {eingesehen} clearly at the end of what we
called the first part. Kant merely added the second part of the sentence, “con-
dition and support of the ‘what’” 5 factically, without ever deriving it. We will
append it genetically. In brief, the task of deducing 6 appearance (as we have so
far labeled our second part) 7 is entirely the same as completely proving the
stated sentence from being {esse}. 8 Now so far we have inferred, and proved
to this extent {in tantum}, that light 9 and being are completely identical,
because we are, and are light, all of which is surrounded by facticity. Further,
in our view the light carries another qualitative character, the “from,” and
lastly absolute genesis as well. At the highest point of our speculation, we said
that the light is absolute, qualitative oneness, which cannot be penetrated
further—briefly, thus, an occult quality. Now for the first time we have
arrived at a property of light through which it shows itself immediately as
one with the being previously seen into: certainty {Gewissheit}, pure and for
itself, as such.) 10
2. How would you proceed, if I asked you to describe this certainty more
closely, to make it clear? Not otherwise, I believe, than by conceiving it as an
unshakable continuance and resting in the same 11 unchangeable oneness; in the
same, I said, and therefore in the very same “what,” or quality. Accordingly,
you could not describe pure certainty otherwise than as pure unchangeability;
and [you could not describe] unchangeability otherwise than as the persisting
oneness of the “what,” or of quality.
3. I inquire further, is the main thing for you in describing pure, bare
certainty that the “what” be something particular, or does your description not
rather explicitly contain absolute indifference toward all more exact [348–349]
determination of the “what”? Only the latter, you will say, and without doubt
you will admit that 12 the “what” remains one; but by no means what else it may
be. Therefore, in this case merely the pure form of the “what,” or quality in
general, is employed in the description, and it is the required description of
pure certainty only on condition of this formal purity.168
The Science of Knowing
4. With this, the concept of the “what,” or quality, is completely
explained and derived for the first time. This is a concept that has so far always
remained in the dark, as much in its content as in its factical genesis. Quality
is the absolute negation of changeability and multiplicity, purely as such; i.e.,
the concept is thereby closed without any possibility of further supplements.
I add here as a supplement that through this negation, the negated term
(changeability) is posited at once, purely as such, and without further deter-
mination. (Quantifiability through quality, and vice versa.) 13 “Genetically
derived,” I have said. As such, certainty cannot be described otherwise than
through absolute quality. “If it should be described, then ____ must” and so
on. Therefore, we have taken a very important step for our primordial deriva-
tion of the “what” = appearance. Now everything depends on how the descrip-
tion (i.e., reconstruction) of certainty arises.
5. We argue thus: In this way we have seen into and described certainty.
However, is our description of certainty itself then certain, true, and legitimate?
As we have always done with similar questions, let us pay attention to
our manner of proceeding. We have constructed a general “what” 14 and posited
it as unchangeable. The essence of certainty appeared therein. I ask, if we
repeat this procedure infinitely many times, as we seem able to do, could we
ever do it any differently? The “what’s” construction is entirely unchangeable,
and in all these infinitely many repetitions, it is possible in only the one way
we have described: through negating changeability. Therefore, we view our-
selves in the very same way we have described certainty, as persisting
unchangeably in the construction’s single same “what”; we are what we say, 15
and we say what we are.
6. Certainty is grounded completely and absolutely in itself. However,
according to its description, certainty is persistence in the same “what.” Thus,
in the description, the ground of the “what’s” oneness is to be posited com-
pletely inwardly within certainty itself. The oneness of the “what” lies in this,
that certainty is, and in no other external ground.
[350–351] 7. “Certainty is grounded in itself ” also means that it is
absolute, immanent, self-enclosed, and can never go outside itself; in itself it
is I; in just the very same way that the proof of being’s form was previously
conducted. Therefore, it is clear that the externalized and objectivized cer-
tainty we presented previously is not the absolute one, according to its form,
although in its content and essence it may very well be. Hence, it is clear that
in searching for the absolute we must abstract from the latter, and search sim-
ply in what manifests itself as immanence, as I (or We).
In this We, we have now surely found certainty, as the necessity of rest-
ing in the procedure’s qualitative oneness; 16 and there the matter now rests
(this enumeration demands all our attention). First and foremost, absolute cer-
tainty 17 is, in and from itself, the same as “I” or “We,” us or its own self (all ofTwenty-third Lecture
169
which mean the same thing), completely inaccessible, purely self-enclosed,
and hidden. For if it (or “We,” which means the same thing) were accessible, 18
then it would have to be outside itself, which is contradictory. As a result of
the preceding insight, we must abstract from just this: that we are actually
speaking about it now, and thus are manifesting it; and this is the appearance
(which is mere appearance, because it contradicts the truth) whose possibility
we must deduce from the system of appearance.
In a way that arises from a ground that will show itself very soon but
that is not yet clear, certainty now expresses itself immanently in itself (that is,
in us) and thus in every expression, as perception {Anschauung} of a particular,
completely unchanging process. This manifestation shows up here, but only as
an absolute fact, and so an obscurity still remains. (Process is living as living;
the process’s unchanging qualitative oneness is life’s immanence and self-
groundedness, expressed immediately in life itself.) Let us press ahead toward
clarity, to the extent that we can do so here. Immediately living and immanent
being-a-principle is light, and is intuition with an inner necessity.—I say,
“being-a-principle,” therefore just projecting and intuiting. I say the projected
term {das Projectum} is “immediately living and immanent” 19 simply in intuit-
ing, from intuiting, and out of intuiting; and the projecting is just light’s life
as “principle-providing” {Principiieren}. 20 I say “with an inner necessity,” and
thus that this necessity must completely express itself, because it is just “prin-
ciple-providing” 21 as being principle-providing. 22
It is absolute, immanent projection, and so it is a projection of nothing
else than itself, wholly and completely as it inwardly is.—Note, as it is, it is pro-
jecting, first of itself inwardly and qualitatively, but not [352–353] at all under-
stood objectively. Thus—making itself into an inward intuiting, immediately by
way of inner, living “being-a-principle,” as thinking and intuiting uncondition-
ally at a single stroke; but, in fact and truthfully, it is the latter as a result of the
former. Thus in this inner qualitative self-projecting, it necessarily projects
itself, objectively (et in virtute eius, minime per actum specialem); 23 not yet to be
sure as an objectively present I, but rather as it inwardly is: first of all as living,
one, and grounded in itself formally; but this is process in pure qualitative one-
ness. 24 This is the intuition of inner certainty and oneness of process that con-
cerned us previously. This oneness expresses itself with necessity, 25 because it is
the result of the absolute, living principle-providing. We called this process
“describing certainty” and sought a ground for it. It has been found. I ask to
wit, does such a description of certainty happen in itself and actually? But how
could it; it is nothing else than the necessary expression and result of certainty’s
life as pure “providing-itself-a-principle” {Sichprincipiierens} completely derived
and explained by us. Yet this life is unconditionally necessary in being or cer-
tainty. 26 Further—it projects itself as it inwardly is; but it is not merely living,
instead it is its own life, and, as such, it is self-projection. The life derived in this170
The Science of Knowing
way as a constructive process is therefore the construction of itself in the pro-
jection (and therefore likewise in the certainty taken objectively), which we
found as a first term at the beginning of our investigation, when we were unfa-
miliar with the higher terms. In the living description, certainty is more pri-
mordial in us than it is objectively in itself without any description. It is the lat-
ter only as a result of the construction, which is also projective.— 27
Now let us more clearly and definitely grasp the three primary modifi-
cations of the primordial light, which we have disclosed today.
Certainty, or light, is an immediately living principle, and thus the pure
absolute oneness 28 of just the light, which cannot be described further in any
way, but rather can only be carried out. If we wished to describe it, then we
would have to describe it as a qualitative oneness, which in this case does not
serve us. It is always and eternally 29 immediate I. Therefore, since we said the
foregoing, we already contradicted ourselves in what we say. 30 When we said,
“it is a living principle” 31 we began to describe it, but [to do so] primordially.
“Principle-providing” {Principiiren} 32 is already its effect, but its primordial
effect 33 in us, since we are it. It (or we, which is the same thing) describes itself
in this way. Principle-providing, if only you [354–355] think it precisely, is
projecting, 34 immanent self-projecting. To be sure, since it lies immediately and
directly in living 35 itself, it [consists in] making oneself into projection and
intuition, 36 not through a gap and objectively, but inwardly and essentially 37
through transubstantiation. Notice, because this is situated in light’s very life,
all light whatsoever is immediately 38 self-creating, and so it is like this:—thus,
it is absolutely intuiting. Even the science of knowing, in all its living activity,
cannot avoid this fate. We have not avoided it either, despite the fact that,
according to a law that we have not yet explained, it invades the principle and
becomes a self-making 39 in it, because otherwise [the principle] remains just a
being. This entire insight into the primordially real principle-providing is a
matter for the science of knowing—which is the first factor.
Living knowing intuits itself unconditionally as it is inwardly just
because it really projects itself. However, above all it is unconditionally 40 from
itself; it must therefore intuit itself as being thus, and here specifically as not
existing outside intuition. Here, therefore, the absolute gap arises, and the
projection through a gap, as a pure, rational expression of the true relationship
of things: the notion of intuition, or the concept, in its separation from
essence, not as the essence itself, but rather merely as its image, and the nega-
tion of the latter beforehand.
Hence, it is a “principle-providing,” and it must intuit {anschauen} itself
objectively and through a gap. At this point it is clear that this “principle-pro-
viding,” its process, must appear within the intrinsically immanent view as
from itself, out of itself, and through itself; as by no means grounded in it, but
rather as happening because it projects through a gap.—However, the scien-Twenty-third Lecture
171
tist of knowing, who understands {einsieht} the view itself in its arising, knows
very well all the same that this entire independence is not intrinsically true as
self-producing, but instead is only the appearance of a higher, absolutely unin-
tuitable, principle-providing. Therefore [he understands] that all the flexibil-
ity {Agilität} lying in the procedure’s appearance is not grounded in truth, so
that it can only be grasped with difficulty, just as the qualitative oneness of
intuition by no means comes to an end with it.
So then, it is an absolutely immanent 41 providing-itself-a-principle {Prin-
cipiiren seiner selbst} and indeed, as will now be explained more exactly, in
absolute fact, without any other intervening light, or seeing: as intuition. This
very intuition must itself be intuited, or projected through a gap: through
which arises the intuition of a primordially complete and persisting knowing,
what we previously called the being of knowing. The first described
[356–357] principle-providing in intuition relates itself to this “being of
knowing” as reconstruction. So, from certainty we have derived both subordi-
nate terms (which previously just stood there hypothetically) out of a deeper
insight into the essence of their relation.
We have not sought to conceal 42 where the difficulty now still remains.
That is, [we need] to ground the possibility, and justify the truth and validity
of the science of knowing’s presupposition that living certainty is genuinely
“principle-providing.”—I say carefully “principle-providing,” not “providing-
itself-a-principle.” If we prove the first, then the second follows from absolute
immanence and self-enclosure, as is completely obvious from itself. More
about this tomorrow.C H A P T E R
24
Twenty-fourth Lecture
Thursday, May 31, 1804
Honored Guests:
We have posited the primordial light as “principle-providing” living immedi-
ately in itself, and from this, we have derived three necessarily arising funda-
mental determinations of the light. At first, it was enough for us to understand
this proposition and the inferences following from it, which demanded ener-
getic thinking and inwardly living imagination above all. However, I believe
that I have succeeded in being comprehensible. Finished with this business,
we raised the question, what authorized us, as scientists of knowing, to make
this assumption? We held ourselves back from answering this, until today.
a. To begin with, consider the sense in which we then asked the question. We
know that the science of knowing is I, that light is completely I. Further,
someone could attempt to conduct a proof here in the same way we did it pre-
viously; the science of knowing as I, and therefore light, can and does; hence
the light can and does. Nevertheless, this mode of proof must fall away and
receive its higher premises. When this occurs, the place is revealed where, at
the very same time, the I that can act {könnendes Ich} and the light that can act
fall away, along with all arbitrariness, the appearance of which our presuppo-
sition certainly still carries.
So much generally. Now I request that you undertake the following
reflection with me:
1. 1 If arbitrariness is to vanish, an immediate, factical necessity of imme-
diate self-projection must show up in the science of knowing. “Immediate 2
necessity,” I say; the science of knowing must really do this, or better the
necessity must take place itself, without the science’s help. I say “necessity” and
“factical”: 3 it must be immediately intuited as necessary, but without additional
higher grounds. The remark, which we make from time to time and made yes-
terday, that we cannot escape from knowing’s projecting and objectifying, does
not suffice for our purpose. It has not been shown under what conditions and
This lecture begins at GA II, 8, pp. 358–359.
172Twenty-fourth Lecture
173
in which context we cannot escape. Thus, if we could not escape only in atten-
tive thinking, we nevertheless seem to escape in weak thinking. We should
explain just how it stands with these two contradictory appearances.
However, as I will show you briefly and without further derivation, the
following suffices and leads to our goal. I cannot predicate anything of the
light, 4 unless I just project and objectify it [i.e., the light] as the subject of the 5
predicate. “Non nisi formaliter obiecti sunt predicata.” 6 [360–361] If one merely
ponders this proposition attentively, it is immediately clear, without any fur-
ther possible {anführbaren} grounds; and if it is to belong here, then it must be
just so, according to our previous remark. 7 It is more important for us to
understand its contents properly. “I predicate of the light” (or, what is syn-
onymous, the light predicates of itself ) means that it projects itself through a
gap, by means of the enduring intuition sufficiently characterized yesterday. “I
cannot do so without projecting it” generally means: without projecting it in
the primordial, real, inner, and essential projection, which first makes it intu-
ition, as again was sufficiently described yesterday.
Let us look back at our task. We put forward the very claim about which
we admitted that it was not made in any factical knowing whatsoever (which
is always limited to already completed intuitions) but rather is made only by
the science of knowing—light’s absolutely primordial act of making itself an
intuition. This claim should be demonstrated in an immediately self-produc-
ing insight and manifestness. This is the case with the insight produced by us.
Therefore, absolute necessity, etc., is comprehended {eingesehen}.
But how is it comprehended? Not unconditionally, but under condi-
tions. Our insight brings it into connection with something else. If [some-
thing] should be predicated (that is, be intuition), then ____ must. However,
should the antecedent {erste} be? 8
In that case neither is the consequent {andere}: the first can be seen into
only under a condition—thus it is conditioned intellectually; if the consequent
happens to occur—that [would be] real conditionedness. So everything is just
about what we wished; but not the unconditionedness and oneness for which
we strove.
(a. A logical clarification: Predicate = minor premise; absolute objectifi-
cation of the logical subject = major premise. Both posit themselves uncondi-
tionally reciprocally; and so something else underlies the inference much more
deeply than the major premise, with discovery of which we began. Ignoring
this, all philosophical systems without exception fail to arrive at an absolute
major premise. Therefore, if they do not wish their thinking to stand still
somewhere arbitrarily, they must sink into a rootless skepticism.
b. With the repeated appearance of the “should,” of hypotheticalness, and
of relatedness, no one now 9 fears that we would be driven back again to the
old point that we have already abandoned. Because obviously the currently174
The Science of Knowing
related terms are higher than were the previous ones. Previously, [we under-
stood] self-construction to be the procedure that arose yesterday in the already
established {stehende} intuition, and, [we] took the being of knowing to be the
highest, the established [362–363] intuition. Now, though, the established
intuition itself is not displaced by a higher one in relation to pure and real self-
projection, rather it is simply discovered in it.)
2. Since I, a scientist of knowing, see into this connection as absolutely
necessary and unchangeable, I myself project and objectify knowing as just
this relation, as a “oneness through itself,” determined without any possible
assistance from some external factor; [a oneness] which I simultaneously per-
meate and construct in its inner essence and content. What then is the con-
tent?—Above all, something arbitrary and dependent simply on freedom and
the fact, therefore something unconditionally necessary, which, facticity, if it
happens to be called into life, grasps and determines without further ado.
Both in relation to one another, so that the one is indeed the completely
proper principle of its being, but cannot be this without in the same undivided
stroke having its principle in the other. Likewise, the other is not actually a
principle unless the one posits itself. We can best name the former “law,” i.e.,
a principle that, in order to provide a principle factically {welche zu seinem fack-
tischen Principiiren}, presupposes yet another absolutely self-producing princi-
ple. [We can name] the latter a pure, primordial “fact,” which is only possible
according to a law.
3. This is what knowing is, absolutely and unalterably, without any
exception, and it is understood {eingesehen} as such. Now, in the insight just
produced and completed, I myself, the scientist of knowing, am a knowing.
Indeed, as it seems to me, I am a free and factical knowing, since I could well
have omitted the preceding reflection, and moreover [I am a knowing] pred-
icated of knowing, describing its entire essence.—According to my own state-
ment, facticity always and everywhere takes up the law of primordial projec-
tion, therefore in actual fact it must have taken me up, albeit invisibly. Formal
projection in general rests in the law, and this is evident enough factically; I
add here only that it exists as a result of the law. The fact also rests in the law
that it [facticity] is projected just as it inwardly is, or, (as we could say more
accurately in order to be safe from all misunderstanding) nota bene, as it must
be projected 10 according to the law. However, at the heights of our speculation,
we know no other law at all except the law of lawfulness itself, that 11 it is pro-
jected according to law. This is just how it expressed itself to us earlier in fac-
tical terms; therefore, we have nothing else to do regarding this material point
than to add that this projection occurs according to the absolute law. In that
way, [364–365] by applying its own material assertion to its own form, we
have genetically derived the very insight from which we began today, and
which previously was produced factically. This is the first important result.Twenty-fourth Lecture
175
a. 12 No doubt, we will get to say more about such an application of what a
proposition asserts to the proposition itself. It is remarkable. 13 It seems to be
just the determination of the major premise by the minor premise that we
mentioned previously; and therefore truth seems to begin running back into
itself, which in all other systems would have neither beginning nor end, if they
were consistent. b. On every occasion, we have posited as the absolute of its
kind that which posits itself; and so, in the investigation just completed, we
have [posited] a law of law, or lawfulness itself. That this is the absolute law
cannot be doubted; nor that the act we carry out in accord with it is the high-
est, immediately law-conformable act of the science of knowing. (What the
stated limitations mean will become evident at the appropriate time.)
4. Now we proceed to another investigation that is of the highest impor-
tance, interesting (if I can succeed in making myself understandable to you),
and even agreeable.
Without doubt the absolute law—according to which we projected
knowing in its essence in the insight that we completed today and then ana-
lyzed—has absolute real causality on the inwardness of the act (I do not speak
of its outward form, which appears free). So that the law and the act perme-
ate one another inwardly (and indeed the act with the oneness of all its dis-
tinguishable determinations) without any gap between them. The projection
is in part formal (objectifying), and in part material 14 (expressing the essence 15
of knowing). It is by no means the latter without the former, but instead is
both at a single stroke, because it is both through an absolutely effective law.
Hence, the material expression must simultaneously express the projection’s
form. Or, to say it exactly: disregarding the prior proof of the opposite, know-
ing in projection, at least formally, cannot simply be just what it is in itself or
according to the law, without any projection. What could it be that projection
alters in it? Of course, you do not forget that we are speaking here only about
projection’s inner material form, abstracting from its outer form, on the basis
of which we merely argued and which we now drop. The inner essence of pro-
jection is living principle-providing; this [latter] must remain in [the former]
completely and absolutely as such, and can never be destroyed. What is this?
Answer: only absolute [366–367] description, as description, which stepped
between the two terms in a wonderful way. Among other things, [it is] quite
evident in the relation as such. For what is the relation except describing one
on the basis of the other? The latter must remain immanently 16 in the former
and can never be destroyed, precisely as intrinsically living “principle-provid-
ing.” Thus, it must allow itself to be perpetually renewed as such, although the
content, determined by the absolute law, remains the same. On this basis, one
can now explain the appearance of energetic reflection, and the infinite repe-
tition of the content that qualitatively remains absolutely one, content which,
to our great astonishment, has not yet wished to leave us. I said that the law176
The Science of Knowing
should determine the inner content. While we examine it, if we now think of
description as an absolutely living “principle-providing” in itself, then it is
clear at once that it must appear just through the law as the reconstruction of
an original pre-construction, and so, in a word, as an image. Or, [it must
appear] as the oft-mentioned mere statement, 17 only the enunciation and
expression, of that which to be sure should in itself be just thus. In a word,
[description is] the whole simply ideal 18 element, as which we must consider
our seeing to be, if we transfer ourselves into the standpoint of reflection, i.e.,
of “principle-providing.”
Therefore, to apply this at once, this is how it stands with knowing. The
entire form of objectivity, or the form of existence, has in itself no relation to
truth. Knowing itself, however, and everything which should arise in it, splits
itself absolutely into a duality, whose one term is to be the primordial, and
whose other term is to be the reconstruction of the primordial, completely
without any diversity of content, and so again absolutely one; differing only in
the given form, which obviously indicates a reciprocal relation to one another.
(It is really like this in every possible consciousness, if you wish to test the
proposition there. Object, representation.)
However, let us carry this reasoning further. At the beginning of our
investigation, because the looked-for absolute oneness created us, we stood (as
we later discovered) under the law without knowing either it or our 19 act as
such. Only when we reflected on this act were we able to apply the content of
our discovered insight to it as the middle term of our inference, and arrive at
an insight into the previously concealed law. Without doubt, we constructed
or described the law itself in this insight, and could catch ourselves in the act. 20
Therefore, we truly stand under the law just in case no law [368–369]
arises in knowing; and we are beyond the law, constructing it itself, if it does
arise in consciousness. Thus, our entire inference grounds itself on a mere fact,
without law, which therefore cannot be justified; and the inference itself only
speaks of a law without being or having one. Hence, this reasoning too, how-
ever much luster it shows, unravels into nothing. 21
Applying this to the preceding: the ostensible primordial construction,
which is supposed to justify the reconstruction (that admittedly presents itself
as such), is itself really just a reconstruction, but one that does not present
itself as such. However, the entire appearance disappears on closer inspection.
We must be glad that it now drops away as well, and that this stand-
point was not the highest. Because yet another disjunction lies in it, whose
genetic principle is not yet fundamentally clear, and at which we cannot
remain. Of course, it is not the outward disjunction into subject and object,
which fell away by means of the full annulment of the persistent form of pro-
jection and objectivity; rather [it is] the inner living difference between both:
two forms of life. 22Twenty-fourth Lecture
177
Just how this difficulty is resolved and where our path will go on from
there can be gathered very easily. We described and constructed the absolute
law; that is the difficulty. It must be evident that we cannot construct it; rather
it constructs itself on us and in us. In short, it is the law itself, which posits us,
and itself in us. On this subject, tomorrow.C H A P T E R
25
Twenty-fifth Lecture
Friday, June 1, 1804
Honored Guests:
We understood {sahen wir ein} that, as a condition, if knowing is to predicate
something of itself, then it must in general simply project itself. Looking sim-
ply at the form of this insight, it contains a free, arbitrary fact and an absolute
law, of which this fact, becoming actual, should immediately avail itself. This
is the first point. 1
Now in this “understanding” {Einsehen}, we—as scientists of knowing—
are also knowing, indeed a free factical knowing, and so in the same circum-
stance of which we have spoken. Therefore, we ourselves, along with this fact
of ours, come under the law to project knowing in general, and to project it as
it is inwardly, or according to law. Thus, knowing has disclosed itself in our
insight as objective and unchangeable oneness, and with the absolute mani-
festness that it behaves in just the way we have said. We add only the new 2
point that this too exists in this way because of the invisible law.
This point was argued further as follows. This projection occurs (at least
according to the material, or to the content of knowing stated in it) according
to an absolute law, which cannot not be a law and cannot not have causality.
Hence, it is an absolutely immanent projection and can never escape being so;
and, as we very illuminatingly add, the light in the projection can not be just
the same as it is inwardly or (according to the law) apart from all projection.
What does this mean? It must permanently bear the mark of the intrinsic liv-
ing principle-providing; 3 and [it must] appear in its form as the product of that
sort of “principle-providing.” Therefore, [it must appear] repeatedly to infin-
ity as a re-construction 4 related to the primordial projection through the law.
This establishes the fundamental disjunction in knowing. At this point, we
raised this objection: is not the law itself reconstructed by us? Obviously;
therefore, we never really arrive at a primordial construction and law, instead,
viewing the matter aright, [we] merely have two reconstructions, one of which
This lecture begins at GA II, 8, pp. 370–371.
178Twenty-fifth Lecture
179
presents itself as what it is, and one of which lies. However, [we] can uncover
this illusion. Thus, we still find ourselves in the realm of arbitrariness, not yet
having entered the necessary. Now to solve this.
1. In the first place, why is it, really, that we will not trust this recon-
struction of the law? Because it seems arbitrary. If it showed itself as neces-
sary, then it would show itself also as something conforming to law and as
itself the immediate inner expression and causality of the law. We received an
[372–373] immediately factical law, pervading the facts: posit the law. Now,
merely because it is projected, the projected law can always be the result of a
reconstruction; but at least the inner construction 5 of this objective law is not a
reconstruction, but rather the primordial construction itself.
2. Can this proof of a law’s reconstruction be carried out? I say: as it seems
to me, it can be done easily in the following way. The first, primordially perma-
nent projection intrinsically bears the character of an image, 6 of a reconstruction,
etc. However, the image as such refers to a content; and reconstruction 7 as such
refers to an original. Hence, the task of understanding this concept of intuition
completely and lawfully contains another [task], to posit this prior one.—Now I
ask, how then is the image an image, the reconstruction a reconstruction?
Because they presuppose a higher law, and exist because of this law, we have said;
and therefore [they] point in the image, as an image, to the law, which is already
contained there, at least virtually and in its results. 8 We, scientists of knowing,
stand presently just in the image as an image; thus, it is the law implicitly and vir-
tually present in us ourselves that constructs, or posits, itself ideally. So, what we
undertook to prove yesterday is completely proven, 9 “The law itself posits itself
in us ourselves.” Image as image is the crucial premise {nervus probandi}. 10
Notice also: 1. In order to conduct the proof we have just completed, we
first had really and intrinsically to presuppose the law as the image’s primor-
dial ground, without giving any account of how we arrive at 11 the concept or
its projection. Now we have completely explained how we arrived at it; but the
variation of this concept’s form has not yet been explained. I content myself
here with merely pointing out this as yet obscure section historically, and
adding that its clarification lies in the reply to the question of the science of
knowing’s possibility as a science of knowing, which will now be addressed
continually, but achieved only at the end. 12
2. We have now changed the following in yesterday’s sketch of know-
ing. Knowing stands neither in the reconstruction as such (representation),
nor in the original (the thing-in-itself ), but rather wholly in a standpoint
between the two. It stands in the image of the reconstruction, as an image, in
which image the positing of a law arises immediately through an inner law. 13
This permeation of the image’s essence is the primordial, absolute, unchange-
able oneness. As wholly internal, in projection, it just divides itself projectively
into a permanent objective image and a permanent objective law.180
The Science of Knowing
[374–375] I wish to be completely understood on this. The light lives
what it is in its very self, it inhabits its living. Now it is image—as image, I
have added, that is living, self-enclosed imaging. You have to attend very
closely to the latter; because otherwise, it is not yet clear how you could have
come to the former, to the image as a closed oneness that you have undoubt-
edly objectified. It is an imaging, formally immanent; 14 it images, or projects
itself as just what it inwardly is, as an image. It is intelligibly immanent and
self-enclosed. 15 Yet, the image posits a law, therefore it projects a law; and it
projects both as standing completely in the one-sided, determinate relation in
which we have thought them.
Note further here: according to yesterday’s disjunction, both primordial
image and copy should be qualitatively one, because otherwise they could not
cohere. Now, both cohere inwardly and essentially as image, positing a law,
etc., and qualitative oneness can not come into it in any way. Qualitative one-
ness is the absolute negation of variation; therefore, it can come into play only
where variability 16 can be posited. However, image, as image, is intrinsically
invariant. It is essential oneness, the law of the image is likewise oneness, and
they posit one another entirely only through their inner essence, without any
further supplement. This total removal of the material, qualitative oneness,
which so far has not left us, is a new guarantee that we have climbed higher.
Since we end the week today, I will not begin any new investigation;
instead, I will estimate what we have left to do in the coming week, during
which I intend to complete this entire presentation, should that prove possible.
Initially, it is readily clear that, if we only establish ourselves correctly in
what has just been presented, no possibility can be foreseen as to how we
should escape it and go further. Here we are immediately absolute knowing.
This is an “image-making process” {ein Bilden} positing itself as an image, and
positing a law of the image-making process as an explanation of the image.
With this everything has been unfolded, and is completely explained and
comprehensible in itself. The terms come together to form {bilden} a synthetic
cycle, into which nothing else can enter.
The particular point at which we have closed off further progress is very
clearly evident above. [We did so] through the total negation of the concept
of qualitative oneness. By negating this concept, quantifiability was estab-
lished at the same time that we opened a more comfortable way to descend
into [376–377] life and its multiplicity, as we know from the preceding. Now,
this qualitative oneness has been negated through positing a oneness of image
and its law that repels all variability, a oneness essential in itself but still merely
hypothetical. Since this oneness has become manifest absolutely, it must be
applicable here, so that we cannot possibly prevent ourselves blending this
quality in, without further ado. 17 —As to this, I wished particularly that you
had noticed for yourselves, that we had to descend again to this quality onlyTwenty-fifth Lecture
181
through a term of the necessary relation, one which is produced lawfully. (The
occult quality is entirely cut off.)
Noteworthy too is the fact that the concept of the science of knowing
as a particular knowing has now disappeared completely. What we have
derived is the one pure knowing in its absolute disjunction, which is explained
from its essence as oneness. We now are this one pure knowing; we are now
the science of knowing as well and [we] hope to become it again, so the sci-
ence of knowing is absolute knowing itself, and we are the latter too only to
the extent that the science of knowing is it.
The last consideration leads us to the path on which we can go further;
the science of knowing must emerge again as a particular knowing. Of course,
we know very well that we have not always been the one insight that we now
are and live, but that we have ascended to it by means of all our previous con-
siderations. We must retrace this, our transformation {Werden} into absolute
knowing—not empirically and artificially, as we have before; instead, we must
explain it. In short, we must once more see {einsehen} in its genesis what we
are at the conclusion of today’s lecture. This insight into the genesis—not of
absolute 18 knowing in itself, because this knows no origin; but rather of
absolute knowing’s actual existence and appearance in us 19 —this would be the
science of knowing in specie, in so far as it is a particular knowing whose
nonexistence is as possible as its existence. 20
Now, it may happen that ordinary knowing is the primordial condition
for the genetic possibility of absolute knowing’s existence, or of the science of
knowing. Hence, [it may happen] that its determinations can be explained
simply from the presupposition that the science of knowing ought to arise, 21
and the sum of our entire system resolves itself into the following rational
inference. “If absolute knowing is to appear, then ____ must, etc.” Now, know-
ing is determined like this, so consequently this should clearly happen.
I have said that all determinations must be explicable and understand-
able simply on the basis of the presupposition: “It should unconditionally
____ , etc.”; and I ask you to take this in its full strength. By way of explana-
tion, I add the following: we have seen {eingesehen} that knowing in itself is
unconditionally one, without any material [378–379] quality or quantity. How
then does this knowing descend in itself to qualitative multiplicity and differ-
ence, 22 and to the entire infinity of quantity and its forms (time, space, etc.) in
which we encounter it? We have to prove the following: simply because
absolute knowing’s being can be produced only genetically, and because it can
be so only under just the types of conditions that we find originally in living,
therefore life coheres indivisibly with the science of knowing, and with that
which it produces. Everyone must confess that, apart from the extent to which
he elevates himself to absolute knowing, his entire life would be nothing,
would lack worth and meaning, and would truly not even exist.182
The Science of Knowing
In brief: absolute affirmation 23 of the genesis of the existence of absolute
knowing (according to the preceding, no term in this description is superflu-
ous) unites both ends of knowing, the ordinary and also the absolute and tran-
scendental, and clarifies them reciprocally. We must put ourselves into this
point, as the genuine standpoint of the science of knowing in specie. With this,
I believe that I have shared with you a very important conception of our entire
procedure up to now and into the future.
Hence, the entire outcome of our doctrine is this. Simple existence,
whatever name it may have, from the very lowest up to the highest (the exis-
tence of absolute knowing), does not have its ground in itself, but instead in
an absolute purpose. And this purpose is that absolute knowing should be.
Everything is posited and determined through this purpose; and it achieves
and exhibits its true destination only in the attainment of this purpose. Value
exists only in knowing, indeed in absolute knowing; all else is without value. I
have deliberately said “in absolute knowing,” and by no means “in the science
of knowing in specie,” because the latter is only a means {Weg}, and has only
instrumental value, by no means intrinsic value. Whoever has arrived no
longer worries about the ladder.
This result, that only true knowing or wisdom has value, is very shock-
ing in our age, which counts only on external workings, and undoubtedly it
appears to us as a great innovation. It is remarkable that this doctrine, which
is an innovation to our age, is really the primordial one, as is almost always the
case with all our characteristic works and ways. I will prove this, not to sup-
port by age and authority what can prove itself, but rather to give you a par-
enthetical opportunity for comparison. In Christianity, (which may in its
essence be much older than we assume, and concerning which I have fre-
quently [380–381] said that, in its roots and especially in its charter [i.e., the
Gospels], which I hold to be its purest expression, it [Christianity] completely
agrees with realized philosophy) the final purpose, especially in the record of
it, which I hold as the purest, 24 is that people come to eternal life, 25 to having
this life and its joy and blessedness in themselves and out of themselves. 26
In what then does eternal life consist? “This is eternal life,” it says, “that
they know [recognize] 27 you . . . and [him] whom you have sent” 28 (for us, this
means the primordial law and its eternal image); merely know [recognize]. Yet,
indeed, this recognition not only leads to life, instead it is life. Thus, from that
time on, through all the centuries and in all forms of Christianity, and consis-
tently with this principle, faith in the teaching that true knowledge of the super-
sensible is the main and essential thing, has been insisted upon. Only in the last
half a century, after the almost total decline of true scholarship and deep think-
ing, have people changed Christianity into a doctrine and ethics of prudence.
This doctrine has not forgotten to teach that in genuine, truly living
knowledge right conduct 29 arises on its own, and that, even if he wishes to doTwenty-fifth Lecture
183
so, the one in whom the light has inwardly dawned cannot fail to shine out-
wardly. Our philosophy forgets it just as little. Yet, there is a great difference
in doing right from such different sources. Doing right from self-interested
cleverness, or from self-regard arising as a result of a categorical imperative,
yields cold, dead fruit, lacking blessing for both agent and the recipient. Now,
as ever, the former hates the law, and would much prefer that it not exist.
Therefore, happiness with himself and his act never arises. The latter cannot
animate and bring to life that which fundamentally has no life. Only when
right action arises from clear insight does it occur with love and pleasure. Only
then does the act, self-sufficient and requiring nothing else, reward itself.
Discussion. 30C H A P T E R
26
Twenty-sixth Lecture
Monday, June 4, 1804
Honored Guests:
Absolute knowing has been presented objectively and in its content, and if
that were the only issue, our work would be completed. However, the question
still remains how we have become this knowing, since we have become it; and,
if this has further conditions, what are they? The first [is] the transcendental,
in all its determinations, the second [is] ordinary. Our proper standpoint [is]
the reciprocal determination of both. 1
Based on our previous experience, I assume that we are sufficiently con-
strained by the shortness of time from easily comprehending the difficulties.
Therefore, today we will pass immediately into the indicated central point. In
the subsequent sessions, this will make it easier to fill out the remaining gaps
between the center and the designated end-points. I tell you in advance that
this investigation is important and raises entirely new topics. 2
1. How, I say, have we arrived here? It is quite clear that we do not want
to see the steps that we have taken historically retold, and that this entire
“how” contains the de jure question rather than any historical one. Without
further ado, it is clear that, in the knowing just presented, we have, and are the
same as, absolute knowing only on the condition that we are certain of this, or
that we are formally certain in general and that we express this formal cer-
tainty within ourselves in fulfilling this knowing. Otherwise, it could always 3
be certain for someone else without being so for us, and we would in no way
have elevated ourselves to the absolute, although we are able to repeat 4 the
word quite properly. This is the first point. 5
The content presented should itself be the expression of absolute know-
ing, or certainty. Therefore, thought in relation to what has gone before, the
following two modes of inference are possible. We are certain, and so what we
say in this situation of certainty is certain, i.e., it is an expression of inward cer-
tainty; or, this is certain, we see into it; therefore we are certain, or expressing
This lecture begins at GA II, 8, pp. 382–383.
184Twenty-sixth Lecture
185
certainty. Whichever of the premises one chooses, one presupposes that the
nature of certainty is known. Without doubt, we will carry out one of the two
inferences. Hence, we must presuppose that the essence of knowing is known.
2. If it is known. Previously we designated it as resting in itself, thus a
oneness of qualities. We cannot assume it here as this, since this [384–385]
oneness has been completely lost to us in the foregoing. Indeed, it has been
lost in relation to an inner essential quality, as in an image, or in the law of an
imaging, which excluded, in the second power, all variability and hence all
qualitative similarity. Here we must characterize certainty with this higher
concept. Thus, it is essential, immanent self-enclosure (as absolute being was
previously seen to be).
3. Now, for the judgment that we wish to pass, that “we are certain,” we
presuppose a primordial concept, or description, of this immanent self-enclo-
sure. How is this concept or description primordially possible? That means,
the law for such a description should be presented, i.e., be understood {einge-
sehen} by us simply as a primordial description.— 6
Since in fact I rise here to a higher term than all those presented so far,
I cannot connect what is to be presented to these, nor can I explain it from
them. Rather, I can rely only on your creative intuition, which I merely hope
to guide. Everything will become completely clear in context, if only one can
make a beginning toward clarity based on the parts. 7
I say that, description, as such, is inward, immanent projection 8 of the
described. Above all, this in no way means an objective projection through a
gap. Instead, [it means] a projection that recognizes itself as a 9 projection (i.e.,
as something superficial) immediately in itself and negates itself as such; and
so posits something described only through this self-negation, an inwardness.
In a word, it is just pure ideal seeing, or intuition, 10 permeating itself simply as
such. I do not say that this permeation exists in itself. Instead, I commend it
to you, as the “We” of the science of knowing, as that through which alone you
become this “We.” You permeate it, if you grasp it as self-negating, yet posit-
ing in this negation something described, positing itself and us as the exter-
nal, and so negated. “Seeing X” means not regarding the seeing as X; thus
negating it. 11 In this negation, seeing becomes a seeing, 12 and something seen
arises simply if it abstracts from X. This, which must now be grasped in intu-
ition much more deeply than I can express it in words, is what I want to point
out. It is the inner essence of pure seeing, as such, which we have actually real-
ized here in its essence, as surely as you have followed my intuition. This is the
first point. 13
Here now should be a description of certainty, as absolute enclosure
within itself. If we just stay with the elicited intuition of pure seeing, or the
purely formal description, [386–387] then we can derive from it alone this
content: that it is a description of an enclosure within itself. Being is what it186
The Science of Knowing
intuitively projects in this negation of itself in which it still exists {ist}. Because
it projects this within, and by means of, its own ineradicable essence, this
essence is an expression, an intrinsically self-expressive being, i.e., an intrinsi-
cally living and powerful being.
(Because I think that I can cast the highest clarity over the whole, I
make the following remarks:
1. Seeing (which reveals itself intrinsically as material and qualitative external-
ity and emanence, as anyone can discover for himself, if he notices what see-
ing means), thereby annuls itself in the face of absolute immanent being. I
do not say, and have not proven, that it does this on any particular grounds,
but, according to the presupposition, it has done so in our own person.
2. Yet in this self-disclosure in essence and self-negation, it is, because we
continue to know; and it exists with its unchangeable basic property, as
expression. Hence, the being in front of which it negates itself is none other
than its own higher being, in the presence of which the lower, which is to
be objectified as seeing, vanishes. This, its being, therefore carries its pri-
mordial feature, expression, which, since it has now become absolute,
expresses itself.
3. Therefore, when it negates itself as seeing, 14 seeing becomes actual seeing,
inwardly and truly effective, or, as is better here, pure light. Thus, pure light
does not come to be as absolutely inward self-expression, power, and life,
instead it just is. It emerges only in insight; and in that case it comes to be
through seeing’s absolute self-negation in the presence of being.
4. To what extent then, and on what grounds, is seeing negated? Answer:
because it is the expression 15 of another and is set over against another as
what it sees (which other, as a result of its own self-negation, lies within see-
ing itself ), therefore, what is negated is simply absolute intuition as such.
5. This intuition is likewise present in our recently generated insight into
absolute light as a living self-expression, and it is to be negated in relation
to inner essence and truth, despite the fact that it may always factically per-
sist. We could point out its origin, and we have done so in the third remark.
It lies in this insight’s genesis from the negation of another.
6. Now, to make a review of the foregoing. The insight that was so significant
at the end of our first part—that being was an absolutely self-enclosed, liv-
ing, and powerful esse—was nothing other than merely the factically fulfilled
absolute insight whose genesis we have realized here. [388–389] Indeed, it
had come about factically in just the way we have understood {eingesehen}
here that it had to come about, by abstracting from everything, or factically
negating all intuition. Thus, our requirement at that time to abstract, as
regards validity, from the objectivity that did not factically weaken as a result
of intuition, was completely correct and has now been proved.)Twenty-sixth Lecture
187
We continue, and on good grounds we stay with the standpoint origi-
nally held: what seeing intuitively projects in this self-negation, and has been
seen into by us as intuitively projecting, is being; and a powerful being at that.
However, being as such, or objectivity, is self-enclosure. Moreover, a living
self-enclosure {in sich Geschlossenheit} (= living self-enclosing {in sich
Schliessen}) posits a principle of coming out of oneself, which is annulled by
the former. Therefore, this principle is better named a “drive.”
And so we would have realized what we sought, a description (i.e., live-
liness and construction of certainty) as a self-enclosure based simply on the
closer and deeper description of seeing, of ideality, and of intuition as such. 16
In this description, four interrelated terms {eine Vierfachheit der Glieder}
disclose themselves, and this enumeration can clarify for you what has been
presented. First, [there is] seeing in its essence, as absolute externality in rela-
tion 17 to another, thereby as self-negating, and, in this negation, regarding this
other as inwardly 18 self-enclosed being. This yields two terms. 19 Things would
have rested with this pair of terms, if we had not added that, despite this par-
tially valid self-negation, seeing continues in the insight into being’s necessary
arising, and remains what it ineradicably is, an inward externality. 20 Therefore,
being must intrinsically carry its characteristics of life and externality. Pursue
this a bit deeper. Strictly speaking, here there is a two-fold view of seeing,
from which there follows a two-fold view of being; or perhaps the other way
around. Thus, [there is] seeing, as intuition, which therefore intuits itself, from
which mere dead being follows; and [there is] seeing in its inward essence as
absolute externality, whence follow the potency and life of being. 21 It is clear
that, in being, both should exist together completely in one, a living self-
enclosure, or an intrinsic self-occlusion, and therefore being must likewise
merge into knowing, from which it has been deduced. A self-negation result-
ing from the construction of its essence as ____ , and a projection [390–391]
of this essence into being as a result of its negation of itself. Finally, the fourth
term, the principle of coming out of oneself through the self-enclosure’s live-
liness and energy, which is simultaneously posited in its being and negated in
its effect—this principle is immediately clear and needs no further proof.
It is not difficult to add the fifth term, which should belong here
according to historically well-known rules so that a complete synthesis can
arise. That is to say, the arising and persistence of these four terms in our
insight depend on the fact that one understands seeing in its essence as exter-
nality, as we have understood it, and that one negates it as absolute imma-
nence (that is, as intrinsically valid) in the face of the recognized being. Still,
[one must] observe it as itself factically enduring above and within being. This
seemed to us a possible insight, and, if it is brought about, an absolutely evi-
dent one, which in general we freely produced and to which we were sum-
moned. Further, the continuance of these four interconnected terms as a188
The Science of Knowing
whole (I say “as a whole” because if one is posited, the others follow as well)
depends on this, that within the insight that we have described (and which is
properly a description of certainty) one preserves what appears there too as
dependent on our freedom.
Enough about this. I commend it to your memory for future use.
1. Prevented by time from introducing new syntheses, I will merely pre-
pare you for them. If being’s living self-enclosing {in sich Schliessen}, which we
mentioned, really is literally a living enclosing, as we express the matter, then it
posits an act, and indeed an absolute act qua act, an effect of the drive to come
out of itself. Now it should annul this latter without resistance. Both can very
well be united if we consider that this entire “livingness” is only a result of see-
ing’s inner expression and projection, as such and as intuition. This latter has
been overthrown as valid in itself, although factically it cannot be eradicated.
Therefore, the former can always exist and remain in factical appearance,
while in truth the second does not occur in any way. In this way, the unifica-
tion of truth in itself and its appearance first provides again a true ground for
the oneness of quality. Previously, we assumed this only factically, then lost it
in the face of a higher truth, and now have been working to reinstate it. The
drive for coming out of oneself [392–393] always emerges in appearance and
is struck down; undoubtedly, it can form {bilden} the appearance of freedom
and genesis themselves, but it does not enter into truth. Just so this drive pro-
vides the true real ground for quantities. 22 Through the unification of these
quantities with the principle of quality, we hope to arrive at a deduction of the
phenomenon of ordinary knowing’s primordial form.
2. This being, then, resides simply in absolute seeing itself and it is to be
sought only there. As much as in appearance, it does so in the truth, which can
of course remain inconceivable and indescribable in its oneness. [It remains
so] not simply because it lies outside of knowing, which is the old fundamen-
tal error; [it does so] rather because absolute knowing itself is inconceivable,
and absolute conceiving is by no means the same as absolute knowing.
According to the standpoint we have adopted, then, where should we look for
it? This clearly confirmed itself in the insight we completed previously.—For
no reason, we saw very directly on the basis of the matter itself that seeing is
an externality. We perceived equally immediately that this was clear to us, or
that we ourselves saw in the very seeing of seeing, and in the insight into the
latter; it was certain for us. Yet, certainty remains to be described, and it is
described as an inner self-enclosing {Sich-schliessen} against the principle of
coming out of oneself. It has come out of there, seeing sees itself, coming out
as coming out. It sees itself just as externality; it was enclosure {Geschlossen-
heit}. Everything was one in the same oneness; simply seeing into self, which
cannot be described further, but can only be lived. Hence, it was actually not
just a re-description, like the first one presented, rather it is the primordialTwenty-sixth Lecture
189
description, i.e., authentic completion of the one certainty, or of knowing. For
the reasons given, we could later also add to it being arising from negation,
which of course lies within knowing itself as a part of its original description,
along with that which flows further from that in the preceding synthesis. But
the last thing which flows from it was the necessity that knowing hold and
sustain itself within its one unchangeable standpoint, which now for the first
time fully points out knowing’s independence, which does not surprise us. The
one knowing would be completely enclosed by means of this complete syn-
thesis in its formal oneness, and by means of the multiplicity which proceeds
from its essence, if this entire reflection did not appear to be freely created and
(if we wish to pay attention to ourselves) appear as objectified by us our-
selves—as was just mentioned and as it consistently seems to us. Objectivity
and genesis 23 are entirely one, as we have seen previously. However, genesis is
this inner externalization and principle of this oneness. Hence, if our appear-
ance is correct (and it must [394–395] of course hold good and be explained
as an appearance), there must be yet another specific manifesting principle, to
which everything just recapitulated is related as to a oneness that contains it
all. It can only be derived from this as the enclosing oneness, and this, as
enclosing oneness, can only be derived from it.C H A P T E R
27
Twenty-seventh Lecture
Wednesday, June 6, 1804
Honored Guests:
Our continuing task is to show that, if transcendental knowing (= the exis-
tence of absolute knowing) is to arise, then another, ordinary knowing must
be presupposed.
The permanent result of the last hour is this: seeing permeating itself as
seeing 1 necessarily surrenders itself as something independent, 2 and posits an
absolute being. The latter enters into a further synthesis, which yields a
description of certainty as an enclosure into itself.
I presented this as the center, and did not fail to recall that continuity of
connection had not been achieved from there out to the endpoints, but rather
that gaps remained. Today’s project is to fill these gaps from the midpoint out-
ward, according to absolute knowing in its oneness. Clearly, this requires new
terms that are not to be developed from the foregoing. Meanwhile, this new
investigation can thus be tied to the previous one. Undoubtedly, it was we, sci-
entists of knowing, who created this insight into the essence of seeing, and
therefore the other as well, that seeing is necessarily a form of intuition, etc. I
also add that another entirely different insight, which we have overlooked, is
contained in this mastering of seeing’s essence, and we will present it now.
1. Namely, I say: if seeing is posited as seeing, then it follows that seeing
actually takes place; or, seeing necessarily sees. So that you do not overlook this
sentence’s importance because of its simplicity, I can justifiably make it more
difficult, by therefore adding the following preliminary remarks. Evidently,
this sentence is the completion of what is required, but never achieved, in the
scholastic proof of God’s existence as the most real being—a proof with which
you are all familiar—to infer the existence of a thing from the bare thought of
it. Seeing posited as seeing means that it is thought, i.e., formed and con-
structed in its inner essence, just as we have done in the last hour; and this is
done hypothetically, as follows immediately from this construction, without
This lecture begins at GA II, 8, pp. 396–397.
190Twenty-seventh Lecture
191
further ado. Because, from positing what something is, nothing at all is estab-
lished regarding its existence or non-existence, and instead, it remains hypo-
thetical. As surely as this latter alone is posited, and if “seeing-into” and cer-
tainty are enclosed in that, nothing can be said about it definitively. 3 However,
we reasoned as follows: “Seeing 4 necessarily occurs, and so its existence is posi-
tively established and expressed.” This is the first point. 5
[398–399] Thus, we genetically derive being there {Dasein}, the true
inner essence of existence {Existenz}. Providing this derivation has been our
task, and indeed, the sole immediately derivable existence is that of seeing.
This is the second point. 6
2. For either the proof in its full significance or the insight intended here
to arise requires strenuous attention. For instance, one could well say life lives,
which is of course correct; but whether life is situated in seeing, completely as
such, as we have most intimately intuited it in ourselves, that is the ques-
tion.—I handle it like this: seeing’s self-permeation is an absolute negation of
itself as independent, and a [matter of ] relating itself to something external.
It exists only in this self-negation and relating, and otherwise not. Yet this
negating and relating is an act that exists {ist} just in itself and in its immedi-
ate completion. Hence, it is necessary, immediate, and actual; and if the whole
is to be, then it must be and exist. Seeing cannot be posited, except as imme-
diately living, powerful, and active.
3. The insight that we have just completed is now the absolute insight
of reason, that is, absolute reason itself. In this insight, we have immediately
become absolute reason, and have dissolved into it. Let us analyze this insight
more closely, and with it reason as well. First, reflect purely on the last thing,
the insight itself, abstracting from the way we came to it: it itself is a seeing,
insight, thus something external. It is a seeing of seeing, 7 thus seeing itself as
existent, absolute and substantial, a seeing grounded in itself and as such self-
expressive, [a matter of ] certainty and self-enclosure by virtue of inner neces-
sity. Thus, as a result of its absolute inner essence, it is a light that determines
and declares itself, that simply by itself absolutely cannot not be, and
absolutely cannot not be what it is, just self-positing. Self-positing, I say, but
not as itself. We have added the last in a way that must now be explained. Rea-
son’s absolute insight sees, just the seeing under discussion, but not immedi-
ately again its first-named seeing, because this is an absolute insight. Proposi-
tion: the absolute insight of reason brings absolute existence {Dasein} (of
seeing, in fact) with itself. It does so immediately [in the course of ] perform-
ing this action, and as the expression of doing so. Said differently, absolute
reason permeates itself as reason, in exactly the way we have indicated, as
absolute effect. The proposition is important: I think [it] is immediately clear,
and if for this one or that one it is not, I can help in the following way. In the
previous hour, formal seeing permeated itself as something [400–401]192
The Science of Knowing
absolutely outward, negated itself therein, and thereby posited being; not
immediately, but rather by means of ____ , etc. It is otherwise here in absolute
reason; this happens without any mediation. Thereby, reason shows itself in
us, as reason’s reason, thus as absolute reason in this self-permeation and being
permeated by itself.—A doctrine of reason {Vernunftlehre}, through itself, from
itself, and in itself, as this duality {Duplicität} can be expressed. A doctrine of
reason, as the first and highest part of the science of knowing, which does not
become, but rather is unconditionally in itself, and is that which it is.
Let’s go to the other side of this insight, that is to the genetic aspect. On
the condition that seeing has been permeated in its inner essence, reason
posits absolute existence, and permeates itself as positing. I ask: were we
intended to assume a seeing and a permeated-ness of this seeing, apart from
what we have here in reason itself and on its basis? I hope not, because then
we would have two absolutes, which certainly could not be brought together.
Thus, this itself is the expression of reason, grounded in reason itself, occur-
ring in order to arrive at its result of positing being and permeating itself in
this positing. Therefore, reason is intrinsically genetic; and it is intrinsically
completely consolidated, necessary, and lawful in unchanging oneness. How-
ever, it is genetic insofar as it is actually and truly living, and expresses itself
actively. Now, however, it exists necessarily and cannot not be, because it is
absolutely self-directing existence and authentic existence. Therefore, it is
inwardly and secretly what it asserts outwardly about external seeing, and its
inwardly grounded existence and life consists in just this [act of ] speaking
about seeing. But further:
4. [To say] that reason itself is the ground of its own proper existence,
here inward, living, and active, would mean that it posits its life and existence
unconditionally and inseparably from itself. No living and existence is possi-
ble apart from this one, and it is impossible to escape it. Now, however,
although we say that it cannot be escaped, we ourselves have obviously done
so. Consequently, reason’s genetic life, which we have described in positing
seeing’s absolute existence in case this latter is merely thought, is still not rea-
son’s primordial and absolute existence. That it is not this, also follows for
another reason. We appear to ourselves as having completed the very insight
just completed, with freedom. We indeed are reason, because reason is simply
the I, and cannot be anything else than I, according to the proof conducted
long since. Thus “we seem,” or “reason seems,” are completely synonymous.
Now we appeared here [402–403] actually as the ground of the insight’s exis-
tence, but not as absolute ground of its actuality, but rather as the ground of
its possibility, or its free ground. Thus, we are only a mediate, factical appear-
ance and insight, and hence not absolute reason.—This mere possibility is
also evident in another way. Evidently, this entire insight appears as the recon-
struction of an original construction in reason, and we objectify this primor-Twenty-seventh Lecture
193
dial reason, along with its construction. We appear to ourselves as yielding
ourselves to the original law of reason by a free act, and now, gripped by it,
as made into manifestness = certainty. This surrender and this manifestness
appear to us as able to be repeated indefinitely. Further, the following in par-
ticular should be noted and made more precise. The premise for the
absolutely completed, rational insight into seeing’s inner essence as some-
thing absolutely external (despite the fact that we do not let it serve as an
absolute premise, but have instead derived it from pure reason itself ), is still
an intrinsic genetic power for the rational insight into existence, which has
been brought about. It conditions the latter just as the latter has conditioned
it. Hence once again the circle between reciprocal conditionings remains, and
thus the old hypothetical character. The absolute condition has by no means
been disclosed. That hypotheticalness remains can also be shown in another
way, and it is interesting and instructive to give this demonstration.—“If see-
ing is, then thus and so are as well, and from this follows ____ , etc.” Precisely,
therefore, what lies in the conclusion is presupposed hypothetically in the
premise, and so it holds only on condition that the first is given. Now, to be
sure, we say that absolute reason 8 posits it; but it is only we who say it, that is,
arbitrariness and freedom. Of course, reason speaks in the connection; but we
have provided it with speech in the first place, and should not trust it. Rea-
son itself must start talking directly.
5. In what we have achieved so far, we have at least learned where this
might be discovered. “Reason itself is immediately and unconditionally the
ground of an existence, indeed of its own existence, since it cannot be of any
other.” This means: this existence {Dasein} has no further ground. It is not
possible, as before, to provide a genetic premise, on the basis of which it can
be explained further, since otherwise it would not be grounded in absolute rea-
son, but in a reason that would have to be interrogated in its turn. Instead, we
must simply say, it is grounded in reason. [It is] a pure absolute fact.
Now just what is this fact? It has accompanied us all along, and thus in
this last investigation as well. It is expressed in the formula: “Reason is the
absolute ground of its own existence.” If anyone does [404–405] not yet see it,
this is because it is too close to him. We ourselves have continually objectified
reason, and therefore posited its existence, as existence, in the “form of outer
existence.” Now we are reason. Therefore, in us ourselves, that is in itself,
absolute reason is the absolute ground of its existence {Daseins}, that is, of its
existence {Existenz} as such. This is a fact from which we cannot ever escape,
and that cannot be explained or understood from any further genetic premise,
as must be the case if it 9 is to be an expression of absolute reason. (Someone
has encountered the enduring objectification of the absolute, and has opined
that therefore there is no absolute. How could that be, if, as is the case here,
true absoluteness is not found in what is objectified, nor in the objectifying194
The Science of Knowing
[agent], but rather in the immediate event of objectifying itself? From time to
time, one sufficiently remembers and makes clear the fact that the absolute
must not be sought outside the absolute, and especially that we would cer-
tainly never grasp the absolute, if we did not even live and conduct it.)
6. Now in this case we are not just the bare fact, but instead we are
simultaneously the insight that this fact is reason’s pure original expression
and life; that the fact is origin and the origin, fact. The characteristic mark of
the science of knowing consists purely in this synthesis. The bare fact, as fact,
yields ordinary knowing. So initially, this self-presentation of reason itself
occurs through reason, 10 we say: for us, then, how is the latter reason? I think,
if we view it properly, [it is] objectified in the same way as the former, and
undoubtedly the same one reason as the former. Thus—our life, or the life of
reason, does not escape from its own self-objectification. 11 And that which is
one implicitly (and probably for ordinary knowing as well) divides itself
within the science of knowing’s insight into a twofold appearance. Indeed, this
division arises only from the insight, 12 or the genesis of that which in itself is
an absolute fact. This is the first point. 13 Further, and most significantly—we
have not just found this insight (that reason alone is the absolute ground of its
existence) immediately in our proposition. We could not have done so. This
insight alone is what has turned us into scientists of knowing. If we had, the
highest fact would have required another genetic premise above itself; and this
in turn another, and so on to infinity. Then our system would have met the
fate of other systems and found no beginning. Instead of that, we have found
it in the more profound insight carried out at the beginning of today’s lecture.
Now it is merely applied as a universal and absolutely valid proposition.
[406–407] Hence, the absolute insight, on which everything depends here,
our science as well, is possible only through a deeper one, which likewise is
possible only through another, etc. In this way, the genuine ground for the
connection between knowing’s various basic determinations is opened up and
made accessible. The deduction can start from the principle: if the insight
obtained in this way—i.e., the science of knowing—is possible in its princi-
ples, then must ____ , etc. This was the second point. 14 Finally, reason is the
ground of its own existence as 15 reason. For, note well, this and nothing else
was the absolute fact with which we are concerned. However, it is so only in
the insight we have created, because it is duplicated (i.e., is reason as reason)
only within this insight. Therefore, this very insight, or the science of know-
ing, is reason’s immediate expression and life: the single life of reason unfolded
immediately in itself and permeated by itself. Precisely because, just as one
lives in it there is seeing, or it appears, thus reason itself lives and appears. In
its existence, the insight appears to be possible only by means of freedom; it is
also thus actually and in fact. That is, in this way reason shows 16 itself as freely
self-expressing. That freedom should appear is simply its law and inwardTwenty-seventh Lecture
195
essence. If it appears as conditioned, then it actually is so; that is, it must
appear thus, as it is conditioned, etc. This is all unconditionally true and cer-
tain, but [it is] true and certain to the exact extent that it lies within reason as
its necessary appearance and expression. By no means is it as certain as reason
in itself, to which it in no way comes, except in its expression. The task on
which the truthfulness and certainty of our insight depends is simply this: to
see everything in its context, and to articulate it in this context. Appearing 17
truly exists when it is conceived as the absolute appearing and self-expression
of absolute reason, and without the latter qualification, it is not true. It is true
that reason appears 18 thus and so, e.g., as inwardly free, only insofar as it also
appears as inwardly necessary and actually existing. Without this qualifica-
tion, it is not true, etc.
Absolute appearance, or genesis, has been presented. The law for deriv-
ing it, and for deducing inferences from it, has also been presented, and the
deduction can proceed. With this, I stop at the border of a philosophia prima
(I would have wished to regard these lectures as being this), [408–409] pre-
senting only appearance’s first basic distinction, which in its oneness consti-
tutes the concept of pure appearing as such, without any further determina-
tion of the latter. In this work, I can either go to work on the details (in which
case I will not end this week), or I can present the main point briefly and
forcefully in a single lecture. To do the latter, I will need more time for the
preparation of this lecture. I prefer the latter, and ask you to permit me to
defer tomorrow’s lecture, so that I can end on Friday.C H A P T E R
28
Twenty-eighth Lecture
Friday, June 8, 1804
Honored Guests:
Here is where we stood: reason is the unconditional ground of its own exis-
tence and its own objectivity for itself, and its primordial life consists just in
this. I assume that this intrinsically dull sentence has become familiar to you
in all its importance in the preceding lectures, especially the last.
Without doubt, we see {sehen ein} this, ourselves intuiting and objecti-
fying 1 reason—as the logical subject of the sentence, which according to us
should now further objectify itself—as its predicate. Numerically reason
occurs here twice, once in us and once outside of us. I ask, which one is the
sole absolute? That is, does our projection perchance result from the primor-
dial projection of the external reason, about which we have been speaking, so
that we ourselves said what it is that we were? Or is the reason projected out-
wardly the result of its own immediate self-projection in us? In a word, rea-
son exists in duality, as subject and object, both as absolute. This ambiguity
must be removed.
1. We can ground 2 this entire existence most effectively with the formula
already previously used and proven: reason makes itself 3 unconditionally intu-
iting. I say “it makes itself,” 4 and not “it is intuiting.” 5 It positively does not arise
with intuition’s being, and to the extent that it appears to arise there, we must
abstract from it. I say, “it makes itself unconditionally”; 6 i.e., as we express our-
selves meanwhile: this making is in no way accidental for it. Instead, it is
entirely and completely necessary. If its being is posited, the latter [intuiting]
is posited as well, and being unfolds in it. This is its proper, immediate, and
inseparable effect. (All these propositions are easy, if one just takes them up in
the strength of energetic thinking. Taken weakly and inconsistently, they are
confusing and lead to nothing.)
NB: pure, intrinsically clear and transparent intellectual activity 7 consists
in this “making itself unconditionally intuitive” 8 in genuine vivacity and exis-
This lecture begins at GA II, 8, pp. 410–411.
196Twenty-eighth Lecture
197
tence. It is raised above all objectifying intuition as itself the latter’s ground, and
it completely fills the gap between subject and object, thereby negating both. 9
Further: This self-making 10 is just effectiveness {Effectivität}, inner life and
activity; indeed, the activity of self-making, 11 [is] thus a making oneself into
activity. 12 Here arise at one and the same time an absolute primordial activity
and movement, and also a making (or copying) of this primordial activity as
its image. (The [412–413] former fundamentally explains the self-making and
presenting manifestness that has grasped us in every investigation; the latter
explains our reconstruction of this [manifestness] as it has also consistently
appeared to us.) Now, though, we have to stand neither in one nor the other,
but rather in the midpoint between both, that is in the absolutely inwardly
effective self-making, real through itself, without any other making or intu-
ition. We must abstract from the fact that we go on to objectify this midpoint
as well, and should in no way admit the validity of that intuition. Because oth-
erwise we have in no way explained the subjective and the objective, but only
added it factically. This was the first point. 13
Reason is this, as certainly as it is: but it is unconditionally. Now, reason
is an absolute, immediate, self-making; thus, it sets itself down as existing,
objectively, and as making itself, objectively. The permanent object and the per-
manent subject, which initially come into question; neither one [exists] by way
of the other, as we initially thought, but instead both [exist] by means of the
same original essential effect in the midpoint. This was the second point. 14
Now, the effect through which it [i.e., reason] throws down {hinwirft} a
permanent object is the same one through which it threw down {hinwarf }
objective living, and therefore the primordial construction, objective reason,
devolves onto this objective life. In addition, the effect through which it threw
down the persisting subject is the same one through which it threw down
imaging as imaging, and so this imaging devolves onto the subject. This was
the third point. 15
Result: reason, as an immediate, internal, self-intuiting making, and to
that extent an absolute oneness of its effects—breaks down within the living
of this making into being and making; into the making of being as made and
not made; and into the [making of ] making as likewise primordial, existing,
and not primordial, i.e., copied. This disjunction, expressed just as we have
expressed it, is what is absolutely original.
2. Now, in the business we have just concluded, we indeed have objecti-
fied and intuited the one reason as the inwardly self-intuiting making; but we
realize that we must abstract from this, if we want to recognize reason as the one.
We are also aware that we can indeed abstract from it. That is, although we are
not able to exempt ourselves from it factically, we can indeed think of it as not
valid in itself. In this way, reason is actually seen into. That is, as the primordial
self-making, it fully merges into our imitative image {nachmachenden Bilde}. 16198
The Science of Knowing
Therefore, it is the very same relation 17 immediately in us that we have just pre-
sented objectively. We, or reason, no longer stand either in that objective reason
nor in subjective reason; since [414–415] we must just abstract from both—
instead [we, or reason, stand] purely in the midpoint of the absolutely effective
self-making. 18 In the science of knowing, reason is completely, immediately alive,
opened in itself, and has become an inward “I,” [both] periphery and center,—
since this is the site of the science of knowing—as it had to come; and all this
has happened completely through abstraction. Absolute reason is absolute 19
(accomplished) thinking {Intelligieren} of oneself. 20 Thinking oneself {Selbstin-
telligieren} as such, is reason.
I can then objectify this very absoluteness 21 of reason (or the science of
knowing), produced in this way, or I can not objectify it, because I abstract from
the objectivity, which factically ceases without my help. If I do the latter, then
everything ends here, and reason is enclosed in itself. If I abandon myself to the
former, then I give myself up to a mere fact that is completely cut off and with-
out any principle. 22 It does not arise out of reason, since only those things
derived here arise from it; just as little does it arise from something else, since
it is simply absolutely factically given. 23 Therefore, it is completely inconceiv-
able, i.e., without a principle; and, once surrendered to it, nothing remains for
me except just to abandon myself to it. It is pure, simple appearance.
3. I want to surrender to it. Now if I objectify this absoluteness, then it
appears to me as an objective 24 situation, to which I raise myself 25 through free
abstraction from my initially objectifying reason. “I want to surrender,” 26 I say; 27
hence, within appearance I find an objectively complete “I” as the original
ground and original condition for this situation. (Consciousness, and self-con-
sciousness, as the original fact underlying all other facticity. The thing whose
validity has so far been completely denied [is] once more derived and justi-
fied.) From the preceding knowledge of reason, I know very well what the “I”
is in itself. I need not to learn it from appearance, which would not give me
any information about its essence. It is the result 28 of reason’s self-making.
Consequently—a very important conclusion, and the only possible one if
appearance itself is to be traced back to reason—appearance itself is inacces-
sible only to me, to the “I” projected absolutely through appearance; it itself is
a self-making of reason, of reason’s primordial effect, and indeed as an I. (The
imprint of reason’s effect rests solely in the I, as such). 29
However, a. that I must simply abstract, if this consciousness is to arise,
and b. that I could either do this or not, are found purely in appearance, that
is, in reason’s effect, which is inaccessible to me only in its principle. Thus, that
I am free.
4. What arises through abstraction for me, as a result of the assertion of
appearance? Reason as absolute 30 oneness: this arises, and appears as arising.
However, [416–417] all arising appears as such only with its opposite. TheTwenty-eighth Lecture
199
opposite of absolute oneness, which in this oppositional relationship in turn
becomes qualitative oneness, is absolute multiplicity and variability. Therefore,
if this oneness is to appear genetically, then consciousness, from which one
must abstract as the beginning point, must itself appear as absolutely change-
able and multiform. This is the first fundamental principle. It is also the first
application of our basic principle: if the science of knowing is to arise (i.e., if
the science of knowing is to be precisely an arising, 31 an origin) then such and
such a consciousness must be posited.
5. The I of consciousness in appearance is an inconceivable effect of rea-
son, above all materially. This inconceivability as such enters immediately into
the primordial consciousness, which genesis presupposes, which is uncondi-
tionally changeable, and which exhausts itself in infinite multiplicity. It does
so explicitly as inconceivable, i.e., as real. Reality in appearance is the same as
the primordial effect of reason—the one eternally selfsame. 32
6. The I of consciousness is reason’s effect in the conceptual form of this
effect, which we conceived earlier in the four presented terms. So far, we have
presented these four terms collectively only insofar as we have penetrated rea-
son as an inner oneness. As a genetic abstraction, [this oneness] presupposes
external oneness just as before, and as a primordial effect of reason within
experience, [it] presupposes inner multiplicity. 33 —However, the inner multi-
plicity and separateness must necessarily consist in the lack of correlation in
realizing these four terms, thus in their apprehension as separate principles.
Then, in an absolutely necessary differentiation of principles, we would have
four fundamental principles:
1. In the enduring 34 object, and indeed in what is absolutely transient: the
principle of sensibility, 35 belief in nature, 36 materialism. 2. In the enduring sub-
ject: belief in personality, and, given the variety of personalities, the identity
and equality of personality, the principle of legality. 37 3. Holding to the
absolutely real forming 38 of the subject, which, since this forming is connected
to the enduring subject, conceptually makes the latter into a oneness, leaving
multiplicity just to the former. 39 [This is] the standpoint of morality, as an
activity that proceeds purely from the enduring I of consciousness, continuing
through infinite time. 4. Holding to the absolute imaging and living of the
absolute object, 40 which [418–419] becomes a oneness for the same reason
introduced in 3. above. The standpoint of religion, as belief in a God, who
alone is true for all lifetimes, and alone is inwardly living. 41
Now all of these standpoints are effects of reason, although not realized
as such in their principle, and reason is completely as it is, wherever it is.
Therefore, it is evident that, simply as effects of reason, the three other stand-
points insert themselves immediately into each of the four, without any appar-
ent free act of abstraction. Yet they do so colored by, and in the spirit of, the
dominant principle. Thus, among religious people morality exists as well, but200
The Science of Knowing
not, as with those having the latter as their principle, as their proper work.
Instead, it is a divine work in them, which brings about in them will and ful-
fillment, as well as the resulting desire and pleasure. Other people also exist
besides themselves, and a sensible world. Yet, they do so always only as an
emanation of one divine life. Likewise, too, God exists in morality as a prin-
ciple, but not for his own sake; instead, so that he maintains the moral law. If
they had no moral law, they would not need a God. Further, other people exist
apart from themselves, but only just so that they can be, or become, moral, and
they have a sensible world simply as a sphere for moral action. Likewise, God
is present for the legal standpoint, but he exists only in order to manage the
higher police work that lies beyond the reach of human policing. Legality also
has a morality; however, it coincides with outward integrity in relation to oth-
ers, and ends there. Finally [there is] a sensible world, for the purpose of civic
industry. Fourth, God surely has a place within the principle of sensibility, if
it is left to itself and does not become unnaturally depraved through perverse
speculation. That is, God exists so that he can give us food in due season. 42
With this comes a certain morality, namely to spread one’s pleasures wisely, so
that one always has something more to enjoy, especially not to ruin things
with this God who provides. Finally, there is also something analogous to rea-
son and spirituality and that is also to enjoy these pleasures in the right order
and with prudence. Hence, there are four basic factors in each standpoint (five
if you also add the unifying principle). Together these yield twenty main fac-
tors and primordial fundamental determinations of knowing; and, if you add
the science of knowing’s indicated fivefold synthesis, which we have just com-
pleted, this becomes twenty-five.
It has already been proved that this division into twenty-five forms
coincides with the absolute breakdown of the real, [420–421] or [with the
breakdown] into absolute multiplicity of that effect of reason, which is imme-
diately inaccessible in its oneness. This follows because multiplicity in general
arises out of the genetic nature of reflection on oneness. However, this reflec-
tion on oneness immediately breaks down into five-foldness. Therefore, the
manifold from which it is necessary to abstract breaks down in the form of
fivefoldness, by the same rule of reason.
Above, we have made it immediately clear factically that this entire con-
sciousness is posited and determined solely by the genesis of the science of
knowing. If we abstracted from our rational insight, as a condition objectify-
ing itself immediately in us, then things would rest there, and we would not
attain anything further. The insight into the law of all consciousness arose for
us only insofar as we reflected on it, and thus posited the science of knowing
as genesis, as something that ought to arise, because of our absolute decree. 43
The task we had assumed is therefore completely finished, and our sci-
ence has closed. The principles have been presented with the greatest possibleTwenty-eighth Lecture
201
clarity and determinateness. Anyone who has truly understood and penetrated
the principles can carry out the schematism on his own. Making many words
does not contribute to clarity; it can even obscure things for the best minds.
Perhaps there will be time and opportunity this coming winter for applying
these principles to particular standpoints, for example to religion, which
always should remain the highest, not only in the partiality and sensible form
in which it was grasped previously, but in our science’s inherent spirit, and
from there to the doctrine of virtue, and of rights. Until then, I ask for your
benevolent remembrance for myself and for science, and I give you my thanks
for the new courage and the new happy prospects for science, which you have
given me so amply in the course of these lectures
